% =====================================================================
%  Harald Høffding, Shorter Works. Second Series (1905)
%  Translated from the Danish: Mindre Arbejder. Anden Række. København: Gyldendalske Boghandel (Nordisk Forlag), 1905.
%
%  TRANSLATION -- GENERATED; DO NOT EDIT BY HAND.  Built by ebook-method/mindre-arbejder-2/tr_parts.py build
%  mindre-arbejder-2 --write from the parts in ebook-method/mindre-arbejder-2/tr/ (GLOSSARY.md, TRANSLATOR-BRIEF.md).
%  Made from the collated transcription texts/hoeffding/mindre-arbejder-2/transcription.tex (2026-10-16); [N] in the
%  margin = the page of the original.  Terms: GLOSSARY.md (choices marked ◆ there for Hans to review); for two pieces
%  Høffding's own English (St Louis 1904; Mind 1905) was used as the terminology anchor, the sentences being translated
%  from the Danish.
% =====================================================================
\documentclass[12pt,a4paper,openany]{book}
\usepackage[utf8]{inputenc}
\usepackage[T1]{fontenc}
\usepackage{libertinus}
\usepackage{libertinust1math}
\usepackage[german,french,english]{babel}
\usepackage{textalpha}
\usepackage[protrusion=true,expansion=true]{microtype}
\usepackage{setspace}
\usepackage[margin=1.2in]{geometry}
\usepackage{fancyhdr}
\usepackage{marginnote}
\usepackage{graphicx}
\usepackage{hyperref}
\hypersetup{pdftitle={Shorter Works. Second Series}, pdfauthor={Harald Høffding}, hidelinks}
\pagestyle{fancy}
\fancyhf{}
\fancyhead[LE,RO]{\thepage}
\fancyhead[C]{\textit{H. Høffding: Shorter Works. Second Series (1905)}}
\fancypagestyle{plain}{\fancyhf{}\fancyhead[LE,RO]{\thepage}}
\setstretch{1.3}
\DeclareUnicodeCharacter{0254}{\reflectbox{c}}   % ɔ (id est)
\renewcommand{\thefootnote}{\arabic{footnote})}
\newcommand{\opage}[1]{\marginnote{\footnotesize\textit{[#1]}}}
\newcommand{\essay}[2][]{\chapter*{#2}\addcontentsline{toc}{chapter}{#1}}
\newcommand{\streg}{\begin{center}\rule{3em}{0.4pt}\end{center}}
\newcommand{\stjerner}{\begin{center}*\hspace{0.4em}*\\[-0.9ex]*\end{center}}
\pdfstringdefDisableCommands{\renewcommand{\footnote}[2][]{}\renewcommand{\opage}[1]{}}

\begin{document}

\frontmatter
\begin{titlepage}\centering\vspace*{2cm}
{\large\textsc{Harald Høffding}}\\[1ex]\rule{3em}{0.4pt}\\[3ex]{\Huge SHORTER WORKS}\\[2ex]{\large SECOND SERIES}\\[1ex]
{\small [Mindre Arbejder. Anden Række]}\par\vfill
{\small COPENHAGEN}\\[0.5ex]{\small GYLDENDALSKE BOGHANDEL}\\[0.5ex]{\small --- NORDISK FORLAG ---}\\[1ex]1905\\[3em]
{\small Translated from the Danish}
\end{titlepage}
\chapter*{CONTENTS.}
\streg
\begin{flushright}\small Page\end{flushright}
\noindent\begin{tabular}{@{}p{0.86\textwidth}@{\hspace{0.5em}}r@{}}
Philosophy and Life. (1901) & 1\\
The Present State of Psychology and Its Relations to the Neighboring Sciences. (1904) & 21\\
The Concept of Analogy in Its Philosophical Significance. (1904) & 33\\
Preface and Postscript to My Philosophy of Religion. (1905) & 47\\
On Some Recent Works in the Philosophy of Religion. (1903) & 58\\
Parmenides of Elea. (1901) & 141\\
Holberg as a Philosopher. (1903) & 150\\
Henrik Steffens (1902) & 162\\
Henrik Steffens in Danish Intellectual Life. (1902) & 172\\
Søren Kierkegaard. A Characterization. (1903) & 181\\
Fichte's ``Addresses to the German Nation'' and the Danish North Schleswigers (1899) & 192\\
Harnack on the Essence of Christianity. (1901) & 203\\
Speech for Georg Brandes. (1891) & 209\\
Speech for Henrik Ibsen. (1898) & 213\\
Review of Bjørnson's ``Poul Lange og Tora Parsberg''. (1898) & 217\\
Julius Lange's Letters. (1902) & 224\\
Goethe. (1899) & 226\\
\end{tabular}
\par\medskip\noindent{\small The page numbers are those of the original, given in the margin as [N].  The contents and
the heads of the pieces differ in places in the original; both are translated as printed.}
\mainmatter
% --- p. 1 ---
\setcounter{footnote}{0}%
\opage{1}\essay[Philosophy and Life (1901)]{PHILOSOPHY AND LIFE.}
\streg

Once in his young days Søren Kierkegaard met F. C.
Sibbern on Gammeltorv and asked him how philosophy
is related to life in reality. Sibbern was, as he
himself has told, surprised at the question, because all his philosophy
aimed at studying life and reality; but he understood
that one who was strongly influenced by Hegel might ask
in this way. Yet the question does not fall away because life and reality
become objects for philosophy. On the contrary, it is
thereby posed all the more sharply. For life is, after all, first and foremost
to be lived, not studied, and reality sets us first and
foremost tasks for our willing and acting. In order to study and
think one must place oneself outside life, look at it as at
an object. When Wieland dissuaded the young Arthur Schopenhauer
from studying philosophy, the latter replied: ``Life is a precarious
thing; I have decided to spend mine thinking about it.''
Schopenhauer thus agreed with Sibbern that the philosopher
is to study life; but he is glad that this
study can lay claim to the whole of life, so that one escapes
the vexations of ``real life''. Schopenhauer could have used
as the motto of his philosophy Goethe's words: ``Was im Leben
uns verdriesst, man im Bilde gern geniesst''. It depends, then,
on the relation between the image of life and life itself. The disagreement
between Sibbern and Schopenhauer consists in this, that Sibbern
% --- p. 2 ---
\setcounter{footnote}{0}%
\opage{2}is enthusiastic that through philosophy he stands in a relation to
life, whereas Schopenhauer wants to use the occupation with the image
of life to escape from life itself.

Perhaps the whole question can be illuminated most clearly if
the weight is laid on this, that philosophizing is itself a form of
spiritual life. The conflict that seems to exist between science
and life arises for the most part from the fact that those
who most often have the word ``life'' in their mouths overlook that thinking
and research too are a form of spiritual life, while the thinkers
and researchers are inclined to place the life of thought above
or at any rate outside the other forms of spiritual life. In
the question of the relation of philosophy, or of science generally,
to life, one is thus really asking about the relation between
two forms of spiritual life.

\section*{I.}

The reason why one, in the case of philosophy rather than of other
sciences, comes to ask about its relation to life
lies in this, that philosophy occupies itself in particular with questions
that concern life in its whole and lie on the border
between the life of thought and the other forms of life. Philosophical problems
concern the nature, conditions and limits of knowledge, the
essence and worth of existence, and the basis for the valuation
of human actions and social conditions. It is the task of philosophy
to gather and use everything that can serve
to illuminate these questions, and thereby to orient us toward
the ultimate riddles of life and existence. In its attempts at this
philosophy appears as an independent form of intellectual life,
which reacts determinatively upon the other forms of intellectual life.
Thus the Vedanta philosophy has been of great importance for the
Indian intellectual development; Platonism has powerfully influenced the
development of the Christian church in antiquity and the Middle Ages, and in the century
just past the ideas of Spinoza and Kant have influenced, both directly
and indirectly, the mental development, even among those
who thought they had nothing at all to do with philosophy. There
is here constantly a work of thought to be done, without which intellectual life
will lose clarity and depth. And this work can only be done
% --- p. 3 ---
\setcounter{footnote}{0}%
\opage{3}properly when it is done with full freedom, --- when the work of thought
stands as an independent form of life, not merely as a means,
but as an end, that is, as that which has immediate value.
Philosophy is to seek to settle the account with regard to man's
position in existence as knowing, feeling and willing,
and such a reckoning can only have significance when
it is made impartially, ruthlessly, with full freedom, --- whatever result
it may lead to. Thought is the part of life that
is to cast light on the other parts of life, and it cannot then do
to have these other parts of life call themselves ``Life''
and want to lord it over thought. The experiences thought needs
in order to be able to treat its tasks it must itself seek and itself
acknowledge: only by the aid of thought can one distinguish between
real experience and mere ideas or illusions. In its
attempt to arrive at clarity, that is to say at full consciousness
of man's position in existence, philosophy works
hand in hand with all other scientific research and has, in
its most significant periods, not separated its cause from the latter's.

But how difficult it is to arrive at full consciousness!
Behind our clearest consciousness there always lie unconscious
tendencies and dispositions, which can be brought forth only later.
Our consciousness always works its way forth out of a dim chaos,
whose scattered elements are shaped by an involuntary comprehending.
Even our clearest thinking has therefore its Achilles' heel, namely the
point where the unconscious forces are at work. This Achilles' heel can
later be discovered, and thinking may perhaps avoid being wounded
more often in this place. When Thetis dipped her son in fire in order
to make him invulnerable, she held him by the heel; she was not so
wise as to dip him twice and the second time hold him by the
other heel. This wisdom philosophy has won on its way,
the more it has freed itself from dogmatism. Every philosophical
attempt rests upon certain presuppositions, of which the philosopher
cannot be fully conscious, and which express the dependence
in which thought, despite all its energy, always stands to other sides
of life. But such presuppositions later
thinking can bring forth, and on this depends --- apart from the new
material that continued outer and inner experience can bring  --- the
% --- p. 4 ---
\setcounter{footnote}{0}%
\opage{4}progress of philosophical research. Achilles must ever anew be made
invulnerable, since new battles always lie ahead, and therefore there
will always be an Achilles' heel, as long as life and development last.
The philosopher cannot avoid being a child of his time and his people.
In his most active work of thought he is often led by an involuntary
need to find expression for peculiarities and tendencies
to which he pays no heed --- precisely because they lie so close to him. Seen from
this side, philosophy stands in a passive, receptive relation
to the other sides of intellectual life. The philosophical thoughts here become
symptoms or symbols, however much labor their formation
has cost. Plato's theory of ideas stands over against the obscurity
and unrest of the world of experience, as Olympus in its eternal
splendor stands above the struggles and vicissitudes of human life, and in
Kant's ``Kritik der reinen Vernunft'' the eternal striving stands
as the redeeming word, no less than in Goethe's ``Faust''.
What myth and poetry express in the form of the image, the
philosopher expresses in his ideas. And where it is essentially the
dim need of life that drives thought forward, while it believes
it follows its own laws, --- there it can become difficult to decide
where thinking leaves off and where poetry begins.
Perhaps the last riddles can be answered only in the symbolic
form of poetry. It then becomes the business of criticism to take care that the border
between thinking and poetry is observed.

It is thus a two-sided relation in which philosophy stands to
life, --- an active one in the conscious work of thought in the domain of the border problems,
and a passive one in the involuntary symbolizing.
Both sides of the relation have their great importance, and
both will always be demonstrable where the work of philosophy is practiced
on a larger scale. According to the peculiarity of the various philosophical tendencies
the weight will come to lie now more on the
one, now more on the other of these two sides. In the philosophy
of Romanticism --- despite the considerable work of thought that was carried out,
and which nowadays one does not always sufficiently acknowledge --- symbolization
gained the preponderance over actual thinking, whereas
Positivism strove to keep as close as possible to the definite,
given experiences. Between these two tendencies the struggle
has stood in the philosophy of the nineteenth century. We can now at
% --- p. 5 ---
\setcounter{footnote}{0}%
\opage{5}the close of the century look back upon this struggle without partiality
and understand it from the historical conditions under which
the philosophical work had to be done. What quite particularly
characterizes the century now past, in comparison with
its predecessor, was indeed precisely its eye for this
historical conditionedness of all intellectual life, thus also of the life of thought.
It has impressed upon us the relation of dependence in which the single
individual always stands to the whole race. Romanticism and Positivism
each in its own way brought forth this dependence. Johann
Gottlieb Fichte, who can be regarded as the founder of
the philosophy of Romanticism, declared that in the choice of the principles that
ultimately determine our world-view it is the thinker's personality
that decides. A philosophical system was for him
not a dead piece of furniture that one could in an external way acquire
once and for all, but a world of thought that was born
in the innermost soul of man. What philosophy one has depends
on what man one is. Hegel carried this line of thought
further in his brilliant but obscure work ``Phänomenologie des
Geistes'', which was to show how the kind, degree and
direction of thinking are determined by the stage of development to which the intellectual life
had attained. In Comte's doctrine of the three stages it is impressed
that the method and results of knowledge become different according to
man's stage of development: in the ``theological'' stage other modes of explanation
are employed than in the ``metaphysical'', and the ``positive''
stage, the standpoint of the newer empirical science, poses
its questions and demands its answers according to other presuppositions
than the two earlier stages. What in the one stage
is a good explanation, is not so in the other. The procedure of knowledge, its tasks
and results thus always depend on
the general development of intellectual life and of culture. The most comprehensive
point of view on this is given in the philosophy of evolution, which
regards all knowledge in the light of the struggle for existence and
of the necessity of adaptation. The conditions of life and of the struggle of life
determine, through many intermediate links, the forms in which the problems are
posed, and the solutions that one thinks to find.

That philosophy has become so clearly aware of this marks
a great progress in the understanding of the relation between
% --- p. 6 ---
\setcounter{footnote}{0}%
\opage{6}the life of thought and the other forms of life. --- I now go on to
examine this relation more in detail.

\section*{II}

The philosopher is first and foremost a logician. One has even
attempted the definition of philosophy that it is the science of
human knowledge. This definition is not exhaustive,
but it brings out a main point. An investigation
of the nature and origin of our concepts, judgments and inferences
and of the whole way in which they are connected with the
nature and mode of working of our consciousness will always be of decisive
importance in philosophy. It is ultimately the logician's task to
say what it is to \textit{understand}. What is required for the full understanding
of a matter, and how far is full understanding possible? It
is the endeavor to understand that leads to the development of concepts,
judgments and inferences, and when we seek to understand something,
we transpose it into the forms of our concepts, judgments and inferences,
in that we form as coherent series of thoughts as we
are able, on the basis of the content given in experience. But
with what right do we thus transpose the content of existence,
which we have after all not ourselves produced, into the forms of our thinking?
If existence were to be fully intelligible (rational), all
its content would have to be capable of being presented as a great system of concepts,
judgments and inferences, and each single event in the world
would have to be capable of being placed in the same relation to other events
as the conclusion in an inference stands in to the premises.
In all domains of science thought works tirelessly
through observation and analysis to reduce the content of existence
to such simple points of view that a purely logical formulation
becomes possible. The problem of knowledge is posed when the
question is raised whether the content of existence can be transposed in
an exhaustive way into the forms of thought.

For logic only that has validity which can be ordered
without contradiction as a member of a series of thoughts. But even if
such a purely logical formulation were possible of everything in existence,
the content and character of existence would not thereby
be exhaustively characterized. We could at most
% --- p. 7 ---
\setcounter{footnote}{0}%
\opage{7}say that existence can be conceived or understood \textit{as if} its
states and events were related to one another, as
the members of a series of inferences are related to one another. But
the states and events always contain Other and More
than the purely logical relations according to which they can be ordered in our
understanding. The special properties (qualities) with which they
appear are not explained because we can demonstrate with the strictest
accuracy their relation to other properties.
The special color qualities, e.g., are not explained in their peculiarity
because the lawlike relation between the series of colors
and the degrees of refraction of light has been demonstrated. And there remains in any
case the great difference between a series of thoughts and
a series of events or states, that in the series of thoughts
one member is related to the other merely as ground to consequent,
but in the series of events one member is related
also to the other as cause to effect. The event we call
the cause precedes in time the event we call the effect.
The temporal relation has here decisive significance, whereas
in the series of thoughts the conclusion is really given as soon as we have
the premises, however long a time we may need to see it. In
pure logic the temporal relation has no significance. But in existence
something \textit{happens}; the difference between past and future
has here real significance. There is here a ``not yet'' and a
``no more''; there arises thereby an opposition between memory
and expectation, and the question becomes whether we are with
right grounding our expectations on our memories, and whether
the expectations we thus ground are confirmed by
the following events. As soon as philosophy --- with Hume
and Kant --- became clear about the essential difference between the relation
between ground and consequent on the one hand and the relation
between cause and effect on the other hand, the problem of knowledge was
posed in earnest, and it became a great question
whether existence could be assumed to be intelligible, rational
in itself.

As logician the philosopher takes account of only a single spiritual
want, the need for understanding. It is his task
% --- p. 8 ---
\setcounter{footnote}{0}%
\opage{8}to draw the boundary between \textit{the} intelligible and the unintelligible.
But there are other spiritual wants than the need
to understand, and it belongs to the tasks of philosophy also to
investigate these other wants and their relation to that
purely intellectual need. The philosopher is to be not merely a logician,
but also a psychologist. The very concept of understanding is a psychological
concept, since all understanding comes about through a process of consciousness.
All understanding is an utterance of \textit{the personal} life, whose
manifold forms and stages the psychologist is to follow step by
step from its most elementary utterances to its highest. Here
the philosopher appears as an observer and is to practice an art
of description, which may at one time demand abilities of a kind similar to the poet's,
at another leads to the application of the experimental method.
The often delicate and half-hidden nuances in which mental phenomena
appear must be sought out by widely different ways. And
when observation and description succeed, there still remains the explanation,
the demonstration of the causal relations in the mental
domain. Here the philosopher has the nearest-lying
example of a content of existence which cannot be exhaustively
presented in purely logical form. And yet he works constantly toward
such a presentation. Even for the highest mental states he seeks
to find laws and conditions; even their hidden and quiet
growth he seeks to trace. He seeks a continuity in the development of
\textit{the personal life}, which can make it possible to conceive
the preceding states \textit{as if} they were related to the
following as premises to conclusion. But his own direct
experience shows him clearly here that the real states are not
exhaustively characterized because their mutual relation
can be presented in this way. --- ---

Both as psychologist and as logician the philosopher comes into
contact with other sciences, in that he makes use of natural science
and the historical sciences to illuminate the development of the mental life,
and in that he finds in the special scientific methods
applications of the general logical principles. And in
both capacities he will also in the end feel himself placed face to face with
\textit{the metaphysical} problem, the problem of the innermost
nature and essence of existence. The way in which this problem is
% --- p. 9 ---
\setcounter{footnote}{0}%
\opage{9}treated will depend essentially on the result one has come
to in the treatment of the problem of knowledge, and on the conception
one has formed of the nature of mental life and of its relation
to other sides of existence. And if one wholly rejects
this problem, then that too will essentially be for reasons
drawn from logic or psychology. In this connection it will
be of special interest to see what significance an attempt
to treat the problem of existence, that is, to form a
metaphysical hypothesis, can have.

Such a hypothesis can serve as an idea or a principle,
a guiding thought, which poses to us further-reaching questions
and problems where an anxious restriction to the single
experiences at hand would stop research. It may give us a larger horizon
for our work in the special domains, --- a warning against
stopping too early. Thus the assumption that existence, from both its material and
its mental side, forms a great continuity, a great unbroken
system of causal series, will naturally encourage us to go on
seeking until one has in the particular found the definite intermediate links
and thereby filled the gaps that seemed to be in the series of outer
or inner events. On the other hand spiritualistic
and indeterministic assumptions will contain a temptation to
stop research too early and to rest content with the interference of the ``soul''
or of ``the free will'' in the course of events. This
is the one significance a metaphysical hypothesis can have.

Another significance it has by being an attempt to
settle our intellectual account. We draw in the metaphysical
hypothesis the ultimate consequences of our knowledge and
see whither they point. Such a reckoning will be of great
importance for the whole of intellectual life. It is a fundamental orientation
that is undertaken. The history of thought shows us that there
are here certain main systems between which the struggle stands. Their
number and their special distinguishing marks are not the same at
every time. Some metaphysical possibilities may drop away as
invalid or unnecessary, and others come in their place, or
the older hypotheses must be reworked in order to harmonize with
the demands that new experiences make. Perhaps it is an eternal strife we
% --- p. 10 ---
\setcounter{footnote}{0}%
\opage{10}here face, but then it is nevertheless of importance that we become
clear about what the strife concerns, why it is necessary, and
how it must be conducted.

In the border region in which the metaphysical hypotheses move,
and where thinking and poetry can so imperceptibly
pass into one another, the philosopher's individuality, his ``personal
equation'', will inevitably come to play a part. In face of
the human tendency to dogmatism, --- that is to say,
to make one's presuppositions the only right or only
possible ones, --- it is important to be attentive to this.
But it does not follow that every attempt to build up a
world of thought on the border of our certain knowledge is unjustified
and useless, when the attempt is made both with critical
sense and with deep feeling of the sting of the problems. A narrow-minded
skepticism can be just as dangerous for the clarity of intellectual life
and its health as a narrow-minded dogmatism. Neither new doubts
nor old dogmas ought to hinder us from looking freely around us, even
when we have come out into the border regions. Only the inescapable
demand holds, that a harmony must be aimed at between
the scientific understanding of existence and the world of thought
the single individual builds on the borders of science. In the long run
a conflict between the demands of science and of the personal life
cannot persist. --- There is one thing one as a rule overlooks when
one speaks of the opposition between ``faith and knowledge''. Scientific
research is not a purely impersonal concern. It
has never been so for the great researchers to whom the fundamental
thoughts of the sciences are due. In the endeavor to find
coherence in the course of nature and in the life of the mind, and between
both, a personal need expresses itself, a spiritual want of life.
Our soul's involuntary striving to preserve unity
with itself through all changes that are lived through and experienced,
and to comprehend all the various elements that stir
in it, is the ultimate explanation that can be given of the tirelessness
with which work is done to find the causes of
everything that happens. The more the new can be conceived as continuation
or transposition of what was known before, the better
is that need satisfied. The wish to find natural, that is to
% --- p. 11 ---
\setcounter{footnote}{0}%
\opage{11}say: in experience known and demonstrable, causes of events,
is therefore not merely forced upon knowledge from outside, but
expresses an essential side of the nature of our intellectual life. Scientific research is a work of personality, just as art
is, and just as practical development of character
is. There is an analogy, and more than an analogy, between thinking
and character: both demand coherence and consistency
and combat the sporadic, the isolated and the fragmentary.
Only by this view is understood the deep joy that arises
when the work of thought advances and when understanding is won.

\section*{III}

The philosopher is not merely logician and psychologist; he is also
ethicist. It is not merely his task to construct series of thoughts
and find causal series, but also to test values, to
carry through a coherent valuation of human
traits, actions and forms of society. All valuation is determined
ultimately by a fundamental value, an immediate value, in
relation to which all other value holds. What for man
has immediate value, that is to say: carries with it immediate
satisfaction, he seeks to hold fast, and if it cannot
immediately be held fast or attained, he seeks for means
to realize it and fasten it. The ends man
sets himself are determined by what has immediate value for
him, and by the ends he sets himself are in turn determined the norms,
the standard, which he applies in the valuation of traits,
actions and conditions of life. The concepts: value, end
and norm lie thus in the same series. From a certain fundamental value
thinking can advance with logical necessity,
when the requisite experience is at one's disposal to establish the
way in which a trait, an action or a social order
intervenes in human life. The question
becomes whether the trait, the action or the social order
lies in a series that in its continuation leads to the realization
of the fundamental value.

Ethical thinking is thus only a special application
of logical and psychological thinking. The relation between end
% --- p. 12 ---
\setcounter{footnote}{0}%
\opage{12}and means is really not different from the relation between ground
and consequent and between cause and effect. The value of the end
contains the ground of the value of the means, and the means is cause of
the realization of the end. He who has set himself an end,
because something stands before him as valuable, cannot without coming
into contradiction with himself omit to choose the means that show
themselves necessary for the attainment of the end; in any case
he must attribute to these means a derived value. All that has
value for us we seek involuntarily, yet not always with consciousness,
to assert, and this asserting expresses itself first and foremost in
an acknowledgment of the value of the traits, actions and
social orders that can be means and ways to the persistence of
what is valuable. Most perfect will this asserting be
when the very endeavors to find such means and ways,
nay, when these very means and ways can acquire immediate value
for us. The distance between value and norm and between end
and means becomes then as small as possible.

It is of importance to bring out this inner kinship
between ethics, logic and psychology. The prejudice has not yet
vanished that on the transition to ethics one must put off
one's otherwise usual way of thinking and inquiring. Ethical
thinking introduces only the special presupposition of a fundamental value,
but needs otherwise no peculiar laws or principles.
The difficulties for ethics rest partly on this, that the
extent and certainty of our experience do not always permit us to draw quite exact
consequences, partly on this, that every ethical line of thought only
holds on the presupposition of the validity of definite fundamental values.
The ethical investigation stops when the question of
these fundamental values themselves is raised. How we must philosophically
place ourselves when the discussion has reached this point, we shall
show in what follows.

\section*{IV}

Both as logician, as psychologist and as ethicist the philosopher
meets with difficulties that are connected with the limitation of our knowledge.
It will be of interest to see how
he must place himself toward these difficulties.

% --- p. 13 ---
\setcounter{footnote}{0}%
\opage{13}It has already been remarked above that existence cannot be
transposed into a logical system, the most perfect form of understanding
we possess. Exact natural science approaches
this intellectual ideal in the domains where precise measurement
and calculation can replace mere observation and description,
and where it is therefore possible to trace the laws for the course of events
back to a series of equations. But even if
such a reduction of the qualitative to the quantitative could be
carried through in all domains, the difference between the
logical form and the real content would nevertheless always remain. The
real, qualitative content must always be presupposed and taken as
given, but cannot be derived from the logical form itself: it is
thus richer than this form, however significant this form may
be. There exists ultimately an irrational (in the mathematical
sense irrational) relation between the form of our knowledge and
its content. Thought is itself a part of existence, and we use
this part as a form in which the whole of existence is to be
expressed. But this task cannot be exhaustively solved, since
thought always seeks and finds new content, content it must
take as it comes, and cannot conjure forth out of its pure
forms. Even though the content it has so far had to work upon
has fairly willingly let itself be expressed in the forms of thought,
there lies in that no guarantee that it will go as easily for
the content given in the future. The straight lines by which thought
describes the movement of existence might perhaps be vanishingly
small parts of an infinite curve, whose formula it
is not possible for it to find.

It was this irrational relation between the form of knowledge
and its content that Kant in a somewhat unfortunate way expressed by
setting up the much-disputed ``Ding an sich''. In the definite
shape in which he set up this concept, it cannot
be maintained. But as an expression of the irrational relation
between the form of knowledge and its content it cannot vanish
from philosophy, but will, however eagerly one may seek to cover
over it, stick its head up at some point or other in every
philosophical conception. It expresses the fact that there
can be no proof adduced for the rationality of existence. It might
% --- p. 14 ---
\setcounter{footnote}{0}%
\opage{14}indeed be that our whole proud science became possible only because
existence for long times had, so to speak, turned only its
rational side toward us, but that we, if it once turned
one of its other sides toward us, would face a complete
impossibility of setting up coherent series of thoughts
that could be applied to it. The logical ice age of existence would
then begin.

Also by another way the irrational possibilities which
a thorough discussion of the problem of knowledge discovers can be illuminated.
All science proceeds from life, from man's need and will
to orient himself in existence. Truth is in this respect
a product of will. But is truth necessary? Is it the
highest? Can our will not free itself from the rules and norms
it has itself produced? --- The justification of such
questioning (which is not so new as one has thought in recent
times) rests on this, that what for us is truth always corresponds
to our definite stage of life or development, and that there is needed
a new truth when a new stage of life has been reached. Even the most
exact science we now possess corresponds to a definite stage of life,
whose time will one day be past, so that it must be superseded by a
new knowledge. This follows directly from the fact that all knowledge is grounded
on experience. All our science arises through the interpretation
of the experiences that have hitherto been made. From this it does not,
however, follow that one with due good or evil will could without further ado
break out and place oneself outside the truth that has been attained. He
who has really reached a new stage of life will have enough to do with
finding the new truth that corresponds to this stage. Let him
then present it, in order that the others, who have not reached the ``higher''
stage, may get an inkling of the events that lie ahead in the world of truth! If he cannot do that, it is only impotence when he insists
upon his will, and it is a question whether it is really a
will that stirs in him. And those who still stand on the
``lower'' stage will in any case try whether it might not after all
be the old logic and the old truths that
he is using in his protests against ``truth'', and perhaps
even uses in an ``all-too-human'' way.

Where a new stage of life is really reached, knowledge will
% --- p. 15 ---
\setcounter{footnote}{0}%
\opage{15}immediately begin to work in order to find a continuous connection
between the new stage and the earlier ones. Every new religion seeks
thus to regard the earlier religions as preparatory
stages, as ways that had to be walked before the fullness of time could
come. And scientific knowledge seeks by a biological,
psychological and historical (sociological) way to find the laws
and the ways of the transition from one form of life to another. Thus we proceed
when we compare man's knowledge
with the animal's, Greek science with the notions of the savage peoples,
modern science with the Greek. It is
this method that Hegel, Comte and Spencer each in his own way
have applied in order to show the dependence of knowledge on the stages of life.
In such a way knowledge, when it is no longer
possible to build a coherent series of thoughts, will always
go further back and investigate whether there does not at the very first
source of the series of thoughts show itself a connection that can both
explain how the series of thoughts became possible, and why it
could not be continued further. ---

What thus holds for our knowledge in general, holds
in particular for ethical knowledge. Ethics operates, as we
have seen, with the concepts value and end as fundamental.
But here arises then the question: \textit{which} immediate values,
which highest ends are to be taken as the basis of valuation
and determine its norm and standard? There is
possibility for as many ethical systems as there can be
fundamental values or fundamental ends, from which a consistently
carried-through valuation of human traits, actions
and social conditions is possible. But from the one fundamental end
or the one fundamental value to the other no purely intellectual
way leads; the transition takes place through a psychological and historical
process, which may be long and involved and take place through
many displacements and transpositions. Here ethics, as before
the theory of knowledge, goes back to biology, psychology and sociology.
In vain will one seek to evade this difficulty
by defining ethics as the doctrine of society's necessary
demands upon the single individual. If the single individual has power, cunning and
consistency, he will not let himself be impressed by this. Even if one
% --- p. 16 ---
\setcounter{footnote}{0}%
\opage{16}can perhaps force him to admit that he does best to
conform to the demands society makes upon him, he will
not thereby let himself be forced to acknowledge that they are good, justified
and just. This acknowledgment arises only in
him when the persistence and development of society acquires immediate
value for him. He must discover the value of social life. And
this discovery does not occur by mere reasoning. Again here
we come back to the different stages of life: what is ethically
consistent at the one stage need not be so at the
other. The significance of ethical thinking consists in that it draws
all consequences of the fundamental values that make themselves known
at a given stage of life; but it is not sufficient to lead directly
and definitively from one stage of life to another.

But the connection which ethics cannot procure for us, because
the fundamental values, and with them the starting points of the ethical
series of thoughts, can stand in conflict with one another, can be found by
going back to biology, psychology and sociology. Here we can
study the origin, development and disappearance of values. The history of values
can have continuity, even though the ethical systems,
which are projections from the single stages of life, in their principles
stand sharply against one another. ---

But --- the goblin moves along. For biology, psychology and sociology
are after all themselves kinds of knowledge and build upon the
ultimate principles of all knowledge. Also the way in which a biological,
psychological and sociological knowledge at a given time develops
will depend on the stage of life that has been reached. We seem then to move
in a circle: the very science of the development of the stages of life
depends on the character of the stage of life to which the development
has led! This circle nevertheless marks no error in our line of thought.
It marks only that we here stand at a relation of reciprocal action
that makes itself continuously felt: knowledge depends
on life, and the very conception of life one forms depends
on the way in which knowledge is determined in
its dependence on life. This circular movement does not exclude an
advance, --- provided only that the circle we here constantly anew must
describe, again (as epicycle) moves along a larger
circle, and this perhaps again along a still larger, and so
% --- p. 17 ---
\setcounter{footnote}{0}%
\opage{17}on. We stand here at a process that can be continued to
infinity, but whose law we can discover already within our
experience.

Even if we disregard this main problem, there can be
difficulties enough in understanding the development from one stage of life
to another. Continuity cannot everywhere be demonstrated in the particular.
There are points where research again and again stops.
There is, however far the motions of the globes and of the atoms
are traced back, always a certain direction of the motion and
a certain distribution of the acting forces, which could not be derived from
the general laws of motion. In the composition of substances
new substances arise with new properties, which do not
find their explanation in the general law of the conservation of matter
through all changes; on the contrary, precisely
this law presupposes that changes take place, that is to say, that there
arise new properties. The origin of life and of consciousness stands
as an unsolved riddle, and similar riddles we face when the
appearance of the great historical personalities leads human life
into new paths. Here are knots that at least provisionally cannot
be untied. Does it mean that continuity is really broken
at these points, or does it mean only that the connection is
so intricate that we cannot unravel it, or that in the knots
more threads are interwoven than those we know from our experience?
--- However much weight one may wish to lay on such
unsolved or insoluble questions, research nevertheless lets itself
not be stopped by them. Where it cannot explain, it observes and
describes. And it rejoices in the fullness of existence, which
promises us that our intellectual work shall not lack material
and tasks, either in the mental or in the material
domain. Our sense may be dulled, and our thought grow tired, but
existence is ever young, and even when it becomes winter for us,
a new spring can begin for it.

\section*{V}

There awaits the philosopher still a great task. As logician
he operates with the concepts ground and consequent, as psychologist
with the concepts cause and effect, and as ethicist with the
% --- p. 18 ---
\setcounter{footnote}{0}%
\opage{18}concepts end and means. But in the end these different
points of view cannot fall outside one another. They must all
come into their own in our concluding conception of existence.
In the metaphysical systems one of these points of view has as a rule
had the preponderance. Existence has been conceived as a great
logic, or as a great mechanism, or as a realm of values. Can
all these points of view now be united in one idea?

About this is waged the great, perhaps eternal strife between
the world-views. Only then has this strife philosophical interest
when it is waged with clear thoughts and in such a way
that work is constantly done on a definite basis of experience and
scientific understanding. It is here that science, philosophy and
religion collide. It is the task of the philosophy of religion to
investigate the conditions and significance of this border dispute.

An attentive glance at the development of the problems through
the ages leads, in my view, to the result that both religious and antireligious
dogmatism are drawing toward their fall.
No dogma, that is to say: no idea that directly
or indirectly has sprung from the religious feeling,
explains existence for us, either in the whole or in the
particular. There are gaps in our series of thoughts and in our causal series;
but they cannot be filled by appealing to the
interference of gods or souls. Such a gap is an unsolved problem,
neither more nor less. But even if the gaps are filled,
then, as we have seen above, the content of existence
cannot be exhausted in any series of thoughts or causal series. It
is certainly an important contribution to the characterization of existence that
its states and its events can to a large extent be calculated
and predicted through the laws our knowledge can find.
This moment the monistic systems emphasize quite particularly:
since all understanding rests on there being unity and coherence
in everything that happens, an identity behind all the differences, one
concludes that since existence is intelligible to as high a degree
as it is, there must also be unity, coherence,
identity in its essence. But since the possibility is not excluded
that existence stands in an irrational relation to our
% --- p. 19 ---
\setcounter{footnote}{0}%
\opage{19}knowledge, no monism can become more than a speculative
hypothesis.

The unity that logical thinking and causal explanation
trace in existence will be the same, however it may fare
with what we call values. The logical consistency and
the causal connection are the same, whether the starting and
end points of the series of thoughts and causal series we can
form are valuable or not. Understanding is in this respect independent
of valuation.

But thereby the fact does not lapse that human beings
in their personal and practical life work with concepts such as
end and value. The fact that such concepts
are formed and must be formed must be taken into consideration in our
concluding consideration of existence. It is here, as with
knowledge, a part of existence that enters into a relation to
the rest of existence. In the religious feeling a judgment is pronounced
upon existence regarded as a whole. The values man
knows have their origin in real conditions. Their fate
is interwoven in the great connection of existence. The religious
feeling is determined by the relation man experiences between
the values and reality, and religious faith is a firm
assurance that what truly has value will also persist,
whatever surge it must pass through, and whatever changes
and displacements it must undergo. The difference
between religious standpoints will by the nature of the case partly be able to
rest on what it is one attributes value to, --- partly on how far
one knows reality, --- partly on what relation
one has experienced between the valuable one has found and
the reality one knows. Here are causes enough for strife.
And since both the values, reality and the relation between them
are in constant development, the strife will again and again be able to blaze
up in new forms.

The task of the philosophy of religion is to find out what the strife
really concerns, and to investigate the psychological possibility of
uniting faith in the persistence of value with the scientific
understanding of existence and with the ethical work on the discovery
and production of values. Philosophy cannot produce
% --- p. 20 ---
\setcounter{footnote}{0}%
\opage{20}faith. Faith is not thought out and is not constructed, but grows forth
out of life. Faith is an object for philosophy, not a fruit of philosophy.
Philosophy investigates its psychological possibility and its
ethical significance, as well as its epistemological justification.
In this last respect it becomes especially the task to
determine the border between thinking and poetry.

Perhaps the highest is accessible to us only in the form of poetry,
and a poetry of life, in which the personal life gave expression
for its deepest experiences, would more and more show itself to satisfy
the need which formerly could be satisfied only by myths
and dogmas. Then Goethe will be proved right:

\begin{verse}
Und deines Geistes höchster Feuerflug \\
Hat schon am Gleichniss, hat am Bild genug!
\end{verse}

If life is to preserve its health and energy under the advancing
division of labor, then the divided forces must be able to be
united in cooperation. New concentrations must take place, if
the great values of the past are to be preserved in new forms and under
the altered conditions. For this are required great, newly forming personalities,
in whom the poet, the prophet and the thinker work united
on the same work. In relation to such a work it is
only a one-sided labor to philosophize; but it is a labor
that shall and must be done, and which indirectly and by many detours
can contribute to the great process of concentration, out of which
a new stage of life emerges.

% --- p. 21 ---
\setcounter{footnote}{0}%
\opage{21}\essay[The Present State of Psychology and Its Relations to the Neighboring Sciences (1904)]{THE PRESENT STATE OF PSYCHOLOGY\\ AND ITS RELATIONS TO THE NEIGHBORING\\ SCIENCES\footnote[1]{A lecture at the international scientific Congress in St. Louis, 23 September 1904.}.}
\streg

\section*{I}

In the conception of the nature and significance of science that prevails in our day, one conviction asserts itself in
particular: that the great and complex connection of reality can indeed be understood by us only when we isolate, in a
more or less artificial way, individual elements and investigate their mutual relations, but that the elements which our
science thus distinguishes may not without further ado be regarded as reality's own constituent elements.
In other words: the conditions of understanding and the conditions of reality do not coincide.
Our way to understanding is not necessarily nature's way to production. That is the old fundamental thought of critical
philosophy. It has pressed through slowly, especially in recent times, not only among philosophers, but also among
natural scientists who have discussed the first principles of their science.
There then appears the possibility of an irrational relation between thought and reality, --- the possibility that the
analysis of thought cannot let the great synthesis of reality come into its full right. Yet the validity of science is
not thereby shaken, for the analyses and distinctions that we undertake in order to attain
% --- p. 22 ---
\setcounter{footnote}{0}%
\opage{22}scientific understanding must be grounded, point by point, in observation of living and concrete reality.

In no field of experience does this relation assert itself so strikingly as in psychology.  The elements that in the
material field are separated out of the whole do after all lie in space outside one another, however intimate the
interaction between them is assumed to be. But such an external relation of separation cannot be assumed in the field of
mental life. Here the individual element is so interwoven in the whole that it receives its whole determinate character
thereby, and the whole cannot here be regarded as a mere product of the elements. Analysis and distinction therefore
acquire in psychology a more artificial character than in natural science. We dare attribute absolute reality to the
psychical elements even less than to the material atoms. And yet there is, after all, no other way to scientific
understanding of mental life than analysis on the basis of observation and experiment.
Moreover, observation shows us that mental life at its very beginning and at its limit more and more assumes a sporadic
character, appears in flashes and as apparently isolated sensations. Thereby it stands in sharp opposition to the manner
in which it expresses itself where it has reached its full development and maturity; here it decidedly bears the stamp of
the whole, of synthesis, and exhibits an intimate connection with all the individual elements. There arises hereby an
antinomy of significance for the whole character of psychology. For we are unable to conceive mental life as having
arisen through mere combination of the elements between which analysis distinguishes, nor yet through that of the
sporadic flashes to which it seems to be reduced at its limits; and yet there must be an inner connection between
concrete consciousness with its synthetic stamp of wholeness and the elements to which analysis and observation of
limits lead back.

This antinomy has intervened deeply in the course of development of recent psychology. It appears in particular in the
conflict between the two great tendencies, the one founded by Hartley and Hume and continued through the psychology of
association and the Herbartian school, the other founded by Leibniz and Kant and continued through the German Romantic
school and its
% --- p. 23 ---
\setcounter{footnote}{0}%
\opage{23}after-effects. The former tendency leads to a psychical atomism, whereas the latter maintains the synthetic character
of mental life. I shall not here pursue the history of this conflict in detail. There is no psychologist whose general
standpoint and special views are not determined by the position he takes on the relation between atomism and concrete
reality in the field of mental life. On both sides there is a temptation to dogmatism, in that one either one-sidedly
fixes the single point of view as exhaustive, or posits an irreconcilable opposition between the points of view.

There is a psychological atomism that regards the elements of psychological analysis as absolute constituents, whose
mechanical composition produces mental life as its product. It overlooks that the whole psychological problem comes into
force again in each of the psychical atoms. For just as it has shown itself in natural science that the atoms at which
one believed one could stop were whole worlds, in whose interior movements and currents took place that made new
investigations necessary, so closer consideration of the conditions under which our simplest sense sensations occur shows
that they themselves must be conceived as synthetic phenomena, as concrete wholes, in that they correspond to a more or
less complex physiological process.

There is, on the other hand, a dogmatism for which mental life constitutes an indivisible unity, perhaps even a
substantial unity, that is to defy all analysis. It overlooks that mental life, like all life, persists under a constant
struggle with oppositions, inner and outer, and that already for practical introspection decisive differences present
themselves, e.g. at every choice between possibilities. The various struggling dispositions or tendencies can make
themselves known in consciousness in an emphatic way. Mental life persists under a constant striving to harmonize the
conflicting elements, and the energy of synthesis appears in many different degrees of intensity. It is thus not
scientific observation and analysis first that produces the differences, but these belong to life itself. As scientific
analysis again and again by its one-sidedness compels us to go back to the great synthesis of life,
% --- p. 24 ---
\setcounter{footnote}{0}%
\opage{24}so in turn the inner relations of tension of life compel us to acknowledge the reality of differences and
oppositions, and to that extent endorse the significance of scientific analysis. We can only analyze what life has
previously connected; but this connection does not exclude relations of difference between the elements that are
connected.

In American literature the relation between analysis and synthesis in psychology has been energetically discussed by two
such eminent investigators as William James and Hugo Münsterberg.
Perhaps, however, the latter in particular has been led to maintain a sharper opposition between psychology and life than
it is consistently possible to hold fast, since after all it is always life that is to yield psychology its material, and
since it is established that introspection does not begin first with scientific psychology but is a practical want that
presupposes the possibility that real differences exist within the whole of life. ---

The opposition between analysis and synthesis in the psychological field stands in a certain connection with the
opposition between intellectualism and voluntarism. For it is within the field of sensations and ideas that one can most
easily distinguish different elements. It is the most articulated side of mental life, and it is also from this side that
mental life lies most open to observation and experiment. In opposition to it, feeling and volitional life present a
stronger concentration; the synthetic character emerges here in a more pronounced way. Overvaluation of the results of
the analytic method will easily lead to an undervaluation of the significance of feeling and volitional life in reality,
perhaps even to an attempt to regard feeling and volitional life as a mere result of the interplay of sensation and idea.
Against this it militates that not only is it impossible to derive the concentration from the interaction of absolute
elements, but that on closer investigation the whole direction of the development of sensations and ideas is determined
by the interests, values and ends to which feeling and volitional life bear witness. The voluntarism deriving from Fichte
and Schopenhauer, which has found an important support in the biological theory of the
% --- p. 25 ---
\setcounter{footnote}{0}%
\opage{25}struggle for life, also wins more and more adherence as the most essential psychological point of view.

\section*{II}

Both by emphasizing the ever incommensurable relation between analysis and synthesis and by stressing voluntarism in
opposition to intellectualism, we lessen the inclination to assume that the concept of personality could once and for all
be closed, as speculative philosophy and theology have often believed. Positivism and empirical science are often blamed
for denying the concept of personality, and in recent times one often sets the historical conception and the doctrine of
freedom in opposition to empirical psychology. I believe that precisely the empirical, experimental and analytic
tendency in psychology rests on a serious and sincere acknowledgment of the reality of personal life. It rests on the
conviction that the value of the life of the mind is not diminished by being tied to determinate conditions and subjected
to determinate laws. Therefore it investigates the life of the mind boldly along all the paths that stand open to science.

The difference between the various psychological schools really depends on where one finds the problem, and how the
burden of proof is distributed. Is the psychological riddle how there can be unity and connection in consciousness, or
does the enigmatic arise there where the life of consciousness assumes a sporadic character, appears in isolated flashes?
This is the main question. It divides naturally into many particular questions. Because I am convinced that the nature of
consciousness expresses itself in an inner connection and gathering together, tasks enough nevertheless arise: namely, to
find the special forms of gathering together, which can by no means be derived a priori from the general nature of
consciousness, --- and moreover to explain how the life of consciousness under certain conditions can assume a sporadic
character. On the other hand there are also many kinds and degrees of isolation, --- and moreover it is to be explained
how unity and connection can appear in consciousness. And every psychologist must concede that as long as the life of
consciousness persists, a constant struggle takes place between the synthetic
% --- p. 26 ---
\setcounter{footnote}{0}%
\opage{26}and the sporadic tendency. Where the stamp of the whole predominates, it becomes a question whether the whole is a
purely mechanical result or aggregate, or whether it does not rest on more deep-lying conditions.

Pathological psychology seems decidedly to speak for the synthetic conception. A constant work is required, without which
the ``psychological tension'' (tension psychologique, Pierre Janet) cannot be sustained. Mental illness shows itself in
this: that this tension, without which the unity of consciousness cannot subsist together with a rich content of
qualitatively different elements, can be sustained only with great and painful effort, and that it demands especially
strong influences and occasions if disruption or slackness are not to set in. The isolated and the sporadic in mental life
always sets a task that is to be solved. It appears in the incipient forms and in the forms of dissolution, where the task
often cannot be solved.

The concept of psychical energy we may determine as follows. There is so much the more psychical work to be done, the more
elements or tendencies there are that are to be connected in the same state, --- the more diverse they are, --- the
stronger each of them is by itself, --- the more intimately they are to be connected, --- and the more remote in time they
lie from one another and from the present moment. In the individual case it will not be easy to apply the concept of
psychical energy that here shows itself possible. Only by careful observation and all-sided knowledge of the special
determinate and individual conditions can one give a determination of these five factors. The number of the elements or
tendencies, the degree of their difference, the strength of each of them, the degree of connection and the magnitude of
the temporal distance --- all this is difficult to ascertain, and it varies from case to case. And to this is added the
velocity with which the connection of elements or tendencies takes place. Those five factors correspond to mass in the
determination of physical energy.

There shows itself here an inexhaustible wealth of possibilities owing to the highly different relations that can obtain
between these five circumstances. There is here a great field for observation, experiment and comparison. A
% --- p. 27 ---
\setcounter{footnote}{0}%
\opage{27}comparative psychology of individuality is still in its first beginning. Only in the field of the psychology of
religion, especially here in America, have beginnings been made in this direction. In the end no theory will be able to
give an exhaustive presentation of how the various elements or tendencies work together in a single state or in a single
personality. As everywhere, so here too the perfected individualization is not a matter for science, but for art. Only
art can in the end give a synthesis that tolerably lets the great and deep synthesis of life come into its own.

\section*{III}

New problems meet us when we try to characterize the relation of psychology to its neighboring sciences. Psychology
stands in a peculiar relation to the natural sciences on the one side, to the human sciences on the other. And the
relation can be briefly stated thus: that in comparison with the natural sciences it has a decidedly synthetic character,
whereas in comparison with the human sciences it has a decidedly analytic character. Thereby are determined the problems
that appear everywhere where different sciences meet.

As already mentioned, the simplest psychical elements we can distinguish correspond to complicated physiological
processes. What stands as simple psychologically is physiologically a multiplicity. There must thus in the single
psychical element be gathered together what physiologically fills several moments and spreads over a whole region. To
this is added another circumstance. The psychical elements are qualitatively different from one another, whereas there is
reason to believe that the corresponding brain processes present only differences in degree, direction and combination.
What psychologically appears as qualitative differences must therefore from the standpoint of natural science be conceived
as quantitatively different. From the natural-scientific standpoint it becomes easier to demonstrate continuity than from
the psychological standpoint. The old proposition, that nature makes no leaps, psychology is not able to carry through so
completely as natural science is able to.

From these circumstances some investigators have thought
% --- p. 28 ---
\setcounter{footnote}{0}%
\opage{28}that they must conclude that a science of the psychical was possible only if, in place of the relation between
psychical states, one could put a relation between the corresponding brain states. Psychology would then, in order to
become a science, have to be transformed into physiology. Otherwise it would --- one holds --- be impossible to approach
the ideal of scientific understanding, namely the demonstration of continuity and of relations of equivalence between the
phenomena. In face of this view I must first remark that in all fields we begin by demonstrating causal relations between
qualitatively different phenomena; it then becomes the task of later investigations to let this elementary causal
relation be replaced by a more perfect causal relation whose terms stand in relations of continuity and equivalence. Now
it is indeed so that in the psychological field we can scarcely get beyond the elementary causal relation, because we
cannot find psychical units and so cannot carry through any quantitative method. But that does not exclude the
justification of assuming at all a causal relation between psychical states. And this becomes moreover necessary: for if
there are causal relations between the corresponding brain processes, there must also, at least mediately, be such a
relation between the psychical states. Second, it is for the time being only purely schematic constructions that one is
able to give of the brain processes corresponding to psychical states, constructions in which the analogy with the
immediately observed and analyzed psychical states is guiding. We infer from the differences of the psychical states to
differences between the corresponding physiological processes. And this inference is not always easy to draw; moreover it
necessarily presupposes that psychological observation and analysis is perfect. The independence of psychology is thus
presupposed.

Possibly both the simplicity and the qualitative character of the psychical elements are connected with the fact that
they are themselves fruits of a hidden synthesis, so that if we could go farther down into the field of psychical
differentials, go down to differentials of the second or third degree, the whole problem would stand clearer before us.
But here we finally strike the great problem of the relation between the psychical and the physical.
% --- p. 29 ---
\setcounter{footnote}{0}%
\opage{29}Here too the relation between analysis and synthesis becomes of decisive importance.
For the difference that is made between the psychical and the physical is due precisely to an analysis, a distinguishing
between elements which reality always presents in connection. We perform here an abstraction; we break up the whole of
reality and then wonder afterwards how difficult it is to unite again the pieces into which we have divided it. This
point has been illuminated with great force especially by Wilhelm Wundt and Roberto Ardigo. Reality always stands as the
great fundamental synthesis within which we move with all our analyses and abstractions, --- as a full-toned sound of the
whole in contrast to our laborious spelling-out. But there is, as said, no other way to knowledge than the one that begins
with analysis. Our attention goes from point to point and seeks only afterwards to work together its individual results.
And just as little as we dare ascribe absolute validity to our distinctions, just as little dare we regard it as an
accident that our striving for knowledge compels us precisely to make these determinate distinctions. It is a
contribution to the characterization of reality that it can be understood only when careful analyses of its content are
undertaken.

I shall not here enter upon a detailed discussion of the hypotheses concerning the relation between the psychical and the
physical, but only remark that since brain processes cannot be subject to other scientific points of view and methods than
other material phenomena, the burden of proof rests on him who would maintain an intervention of the soul in the course of
the physiological processes. Such an assumption would abolish the independence of physiology, just as the denial of a
causal relation between states of soul would abolish the independence of psychology. There is no reason to depart from
the principle that natural science has followed for centuries, and to which it owes its triumphs: that material phenomena
are to be explained by material causes. Spinoza's words: ``To say that the soul sets the body in motion is at bottom to
say that one does not know how this motion comes about'' --- stand in full force to this day. The only working
hypothesis that makes a cooperation between physiology
% --- p. 30 ---
\setcounter{footnote}{0}%
\opage{30}and psychology possible without encroachment from either side, is that which both within the psychical and within
the physical holds fast to as much continuity as possible, and which conceives the very relation between the two series
of phenomena as a functional relation in the mathematical sense of the word. A concluding metaphysical interpretation then
remains always open, but does not concern psychology as such. The criticism to which parallelism, or as I prefer to call
it, the hypothesis of identity, has been subjected has essentially attacked it as a metaphysical hypothesis. But it is
first and foremost a working hypothesis, and the only one that is actually followed and can be followed at the present
stage of science. It certainly holds of our working hypotheses what I remarked above of our analyses and distinctions: we
have no right to assume it to be a pure accident that our understanding of existence makes precisely these determinate
working hypotheses necessary. A metaphysics that would construct a world-view without any regard whatever to the working
hypotheses that have shown themselves necessary would have no philosophical value.

\section*{IV}

Whereas psychology in opposition to natural science stands with a decidedly synthetic character, it stands in opposition
to the human sciences with a decidedly analytic character. The human sciences deal with human works or human ideals, but
psychology deals with the elements and general laws of human mental life. If we call the doctrine of human works cultural
science, the doctrine of human ideals ethics, then psychology's relation to cultural science and ethics is determined by
the relation between elements, works and ideals. All three investigations advance side by side, and they draw from the
same deep source: the immediate and spontaneous life of the mind, the real, concrete life that is lived, and which no
analysis will be able to exhaust, as little as it will ever have expressed itself completely in any work or in any ideal.
All research here stands over against the infinitely concrete and tries from different sides to characterize its essence,
to
% --- p. 31 ---
\setcounter{footnote}{0}%
\opage{31}express its fullness in determinate forms. But the tones of life are so manifold and sound so intimately together
that no science's system of notation can fully express them. This holds of the cultural historian and the ethicist as
well as of the psychologist. But within this kindred position there is again an interaction between psychology and the
human sciences. In order to find the elements and laws of mental life it is not enough to study the single individual in
its individual states. This requires also the study of human works and ideals, in which the essence of mental life reveals
itself through the ages. There exists no mental life in general; it appears in different forms at different times and in
different places, and in each single one of these historical forms it expresses itself as far as possible according to its
full extent, although the interplay and its elements acquire a different timbre. Here psychology has a significant group
of facts to analyze. The sociological method here steps up beside the introspective, the experimental and the
physiological. Its final end has a more abstract character than that of cultural science and ethics, in that the element
is more abstract than the work or the ideal. Here there is need for the inverse deductive method used by Comte and
described by Mill, in that first a course of development is established that has led to the arising of certain works or
ideals, and thereafter a derivation and explanation of this course of development from general psychological elements and
laws is attempted. By pure deduction one will never be able to derive the special human sciences from general psychology;
a reduction, not a deduction, is the highest we can expect to attain. We cannot e.g. deduce pedagogics, esthetics and
ethics from general psychology; but by observing how the spontaneous art of life unfolds itself in education, esthetic
production and ethical action we can find governing points of view and ideas, which can in turn be understood according
to general psychological laws. Such, too, is psychology's relation to the theory of knowledge and to the science of
religion. Concerning this last relation I have in my Philosophy of Religion expressed myself thus: ``In the religious
field men have made some of their deepest, most intensive
% --- p. 32 ---
\setcounter{footnote}{0}%
\opage{32}mental experiences. In religion (if it is genuine and original) all the elements of mental life work with an energy
and an interplay such as hardly in any other field. Therefore the study of religious mental life is of such great
importance for psychology as a whole''. A reciprocal relation will here certainly establish itself more and more, so that
insight into the nature of the psychical elements and into the laws of their interaction will have a guiding and
clarifying effect upon the work of the special human sciences. A closer consideration of the history of the cultural
sciences and of ethics will also show that directly or indirectly such an influence is constantly at work. If psychology
has a future, this influence will be able to become still more significant and fruitful than it has hitherto been.
Psychology stands in a great debt to its neighboring sciences, to the human sciences as well as to natural science. Let
us hope that it will be able to pay on this debt, even though this --- if the other sciences continue to advance --- will,
one hopes, ever be incurred anew!
% --- p. 33 ---
\setcounter{footnote}{0}%
\opage{33}\essay[The Concept of Analogy in Its Philosophical Significance (1904)]{THE CONCEPT OF ANALOGY AND ITS PHILOSOPHICAL\\ SIGNIFICANCE\footnote[1]{A lecture to the Jowett Society in Oxford, 26 November 1904.}.}
\streg

\section*{I}

When one follows the direction that the discussion of the first presuppositions of science has taken in recent years,
one cannot avoid noticing that the ultimate criterion of their validity is more and more found in their applicability.
The justification for setting up a principle is based on its making it possible to gain clarity and find connection in
observations that would otherwise stand obscure and scattered before us. The truth of principles therefore does not
consist in their agreement with an absolute order of things persisting in and for itself. An order of things first comes
into being for us when, by the help of principles, we bring connection among the observations. We therefore produce
truth ourselves when we are able to find the principles that to the greatest extent and in the most intimate degree can
connect our observations. A critical or dynamical concept of truth is thus in the process of working its way forward, in
contrast to the dogmatic concept of truth, which may be called static, since it presupposes a resting, persisting order
of things that thought had merely to reproduce. This concept of truth was already demanded by critical philosophy, when
it taught that objective validity consists in the full, lawful connection among the observations. But it
% --- p. 34 ---
\setcounter{footnote}{0}%
\opage{34}is of no small significance that not only philosophers but also natural scientists, such as Clerk Maxwell, Ernst
Mach and Heinrich Hertz, are more and more going in this direction. It is not only Dogmatism but also Skepticism that is
thereby displaced. For Skepticism arises through one's presupposing the dogmatic or static concept of truth and then
showing that it cannot be satisfied.

What matters, then, is to set up such principles as make it possible to reach the goal one has set oneself, namely to find
the greatest possible connection within the greatest possible number of observations. The principles need not have other
properties than just those that are necessary in this respect. And in particular they need not, as regards the ideas
contained in them, express or reproduce what in itself or ultimately might be contained in the given observations. Their
task is to supply points of view for the thorough ordering of these, and this task they can solve even if they make use of
figurative expressions, provided only that the inner connection among the images that are employed is such that the
connection of the observations can be expressed and predicted by means of it. It is therefore not a relation of identity
but a relation of analogy between principles (axioms, ideas) and observations (experiences, phenomena) that is
required in order for truth to be gained. Analogy is likeness in the relations within different objects, not likeness in
the single properties. The critical or dynamical concept of truth thus becomes at the same time symbolic. It requires a
system of images or signs by whose help we can read off the relation among the observations and be led to new
observations. We regard the events that experience presents to us \textit{as if} they were merely alterations in our
figurative ideas. It will then be natural to choose the simplest and clearest circle of images that can fulfill the task.
Thereby one arrives at what the thinkers of the Renaissance called the principle of simplicity, and what in our own time
has been called the thinking of the world according to the principle of the least measure of force.

That we have no right to ascribe to our knowledge more
% --- p. 35 ---
\setcounter{footnote}{0}%
\opage{35}than symbolic validity follows first and foremost from the fact that we have no right to ascribe to it properties
that are not necessary for it to be able to reach the goal it has set itself. When analogy is sufficient for this, one is
not permitted to presuppose an identity, which moreover cannot be proved. For the possibility is never excluded that other
principles (axioms, ideas) than those we employ could solve the same task. The transition from ground to consequence
is always more certain than that from consequence to ground, because different presuppositions can lead to the same
conclusion. One and the same goal can be reached by different paths, and that our science has found the absolutely only
path to order and connection in the world of experience will never, ever be capable of proof.

But also for other reasons our knowledge must become symbolic, when it operates essentially with analogies.

It lies in the nature of our mind that we cannot think without images. What we in a special and strict sense call thinking
is always the activity by which the relation among different images becomes as clear and unequivocal as possible. Every
thought, however abstract and sublime it may be, we spontaneously strive to illustrate by bringing forth images, in whose
mutual relation it comes to light. The image may be now full and concrete, now a mere scheme or a sign. It is not only in
quiet thinking for ourselves that we proceed in this way, but also, and quite especially, where it is a matter of
presenting what is thought to others. We think in examples, even if we reduce these examples to quite simple schemes or
signs, and still more clearly is it the case that we must use examples and signs when what is thought is to be
communicated to others.

Our thinking is always an interpretation of something given. It is an activity exercised on the basis of receiving. The
opposition between acting and receiving, between activity and passivity, is admittedly relative, depends on the point of
view and on the stage of development. What from one point of view is passive may from another point of view and in
relation to a lower stage of development be regarded as active. But in each single case, in each single act, thinking
carries out a conversion of
% --- p. 36 ---
\setcounter{footnote}{0}%
\opage{36}something received into its own form, and this work of conversion cannot be proved to be of purely formal
significance, such that there should be absolute identity between the given and the very elements of thought. When
thought, e.g., analyzes a material or a mental phenomenon, no proof can be given that the elements it distinguishes in
order to clarify for itself the given content also exist as independent parts in the given. It is sufficient for the
validity of thinking that the relation among the elements of thought \textit{agrees with} the relation among the elements
of experience.

But thinking itself does not consist merely in a conversion from passive to active form. In each single case it will
moreover prove that thought works by holding different domains of experience together and letting the one illuminate the
other. All our knowledge, both the spontaneously developed and the scientific, is therefore permeated by analogies. This
too is connected \textit{with} the fact that the work of thought is the solution of a task, an attempt to realize an end.
When thinking proceeds to new tasks, it does not choose wholly new means and paths, but seeks as far as possible to help
itself with those used earlier, especially when they have the advantage of being clear and distinct. Thus language has not
formed wholly new expressions for mental phenomena, but has helped itself by transferring expressions that were
originally formed for the material phenomena that had first attracted attention. And since on the whole there appear
decisive differences among the domains of experience,
--- since the facts are not uniform, but form different groups, each with its own peculiarities, --- thinking too is
compelled to proceed by using one domain, one group, to illuminate another, and then particularly in such a way that the
domain that orders itself for thought in the simplest and at the same time most fruitful manner is made the basis in the
understanding of the other. If existence did not present qualitative differences, there would be no necessity for this.
But since the different parts of existence, as we know them through experience, are not uniform, analogy becomes a
necessary path to understanding.

% --- p. 37 ---
\setcounter{footnote}{0}%
\opage{37}I shall now try to illuminate the character and significance of analogy in the different domains of thought where
it is applied.

\section*{II}

Every formation of a general concept takes place by means of analogy. The marks that the general concept contains vary in
the different phenomena for which \textit{it} is to hold, and the general concept therefore becomes possible only through
the relations of the properties agreeing, even if they are each different. The general concept ``horse'' becomes, e.g.,
possible only because the brown color stands in the \textit{same} relation to one horse as the gray color to another
horse. Human properties vary from human being to human being, but what is constant is a certain relation among these
properties, and thereby the general concept ``human being'' becomes possible. We cannot form any concept whose marks occur
in a thoroughly identical manner in the single phenomena. The concept is formed through the possibility of constant
transition from example to example by means of analogy.

This has a peculiar interest, since from the time of Aristotle analogy has been set in opposition to concept. For
Aristotle analogy was the bond between different domains that could still hold after the most comprehensive concepts had
been formed. Kinship by analogy was thus not a conceptual kinship. Strictly speaking, this opposition does not hold when
the very formation of concepts rests on analogy. In place of this opposition there must be set an opposition between
different degrees and kinds of analogy.

In the formation of concepts different examples are held together, and a transference takes place from the one to the
other, such that the simplest and clearest example is made the basis and illuminates the others. When it is a matter of
finding propositions that express the relations in experience and can lead to finding them, the formation of analogy
emerges still more clearly.

In its strictest form analogy is quantitative and is then called by its Latin name: proportion. The relation between two
numbers can be \textit{the same} as the relation between two other numbers, although all four numbers are mutually
different. The analogy is here
% --- p. 38 ---
\setcounter{footnote}{0}%
\opage{38}identity of relation. In such cases one can always find one of the terms of the relation by the help of the
others. Quantitative analogy is constructive. It brings the unknown forth by the help of the known relation to the known,
without needing to consult experience. Where numerical relations can be demonstrated, one is able by their help to read
off the course of events in the world of experience. But this applicability itself rests on an analogy --- an analogy
between the places of numbers in the number series and the places of events in the temporal series. The number series is
the simplest and most exact ordering we can form, and it is therefore natural that one strives to bring all other
orderings and series into as close a connection with it as possible. This analogy has been brought forward by Maxwell and
Cournot.

Since now the changes that can be determined most completely by means of numbers are the movements in space, it is
understood why modern natural science since the days of the Renaissance has aimed at explaining all changes in nature as
movements. It pushes the qualitative differences and changes aside or charges them to the account of the subject; at times
it expresses itself \textit{as if} they were mere semblance. It is nevertheless sufficient to presuppose an analogy between
quantitative and qualitative relations, such that the qualitative relations can be determined and predicted by the help of
the quantitative ones. Science then conceives the qualitative changes as if they were merely quantitative changes, and one
thinks of the changes as taking place in something that presents only quantitative differences. Thus arise concepts such as
matter, atom, ether. They are images, or rather schemes, that can be determined throughout by numbers, --- that make a
consistent application possible, --- and that can lead to reliable conclusions with regard to new observations. --- This
analogy has proved so fruitful and so significant that one has been inclined to confuse it with identity, or at any rate to
regard it as the only possible analogy. But for this there is no justification.

Of still more fundamental significance than the numerical relation is a relation that makes itself felt in all activity of
thought, whether this concerns the quantitative or the qualitative, namely
% --- p. 39 ---
\setcounter{footnote}{0}%
\opage{39}the logical relation between ground and consequence. It is the most intimate connection that can obtain among our
thoughts, apart from the pure relation of identity or coincidence. It is therefore natural that we seek as far as possible
to bring all our observations, all relations among events as experience shows them to us, closer to this relation among
our thoughts. We seek to carry through such a conception of the relation among the events in existence that the one comes
to be related to the other as ground to consequence. As far as possible we order the events in such series that the
preceding terms can lead to an understanding of the later, perhaps even to determining these in advance. But since the
relation between ground and consequence is a purely logical relation, whereas the relation among the events in existence
is a temporal relation, it is by means of an analogy that we apply the law of grounding to the events of experience. This
is the fundamental thought in Kant's well-known doctrine of the analogies of experience, whereby the decisive significance
of the logical principles for real research was maintained, while at the same time the dogmatic identifying of ground and
cause (ratio and causa) by older systems was annulled. Hegel later made a romantic attempt to replace analogy once more by
identity; but even by those who in our days ascribe the greatest significance to Hegel's thinking, this attempt is hardly
regarded as successful. It is one thing that the laws of thought can lead us on our way through the world of experience,
another that existence should be exhausted and fully expressed in the laws of thought.

The number series, movement and the relation of grounding stand as significant examples of the scientific significance of
analogy. Special examples could be taken from the individual sciences. Thus the idea of a wave motion has been transferred
from one physical domain to the other by virtue of the likeness of relations. At times one does not even know what it is
that undulates, only that a wave motion takes place, as with the so-called ``ether''. The significance that analogy has
within biology, in that different organs in different living beings can come to be of the same significance for life,
Aristotle has already
% --- p. 40 ---
\setcounter{footnote}{0}%
\opage{40}brought out. But to illuminate the significance of analogy further I shall refer to the endeavor that has made itself
felt in sociology to obtain, through the analogy between society and organism, a guiding thought, a basic scheme for this
science. A series of eminent authors have sought to lay this analogy as the foundation, but they have striven in vain to
make the image a consistent and fruitful scheme for sociological research. In the first place, society is a connection of
organisms, and the relations within one and the same organism cannot be transferred to a group of organic beings. Next,
there is within the organism a significant difference between elements to whose function consciousness is attached and such
in which this is not the case; in society, on the other hand, all beings are conscious, and they therefore stand as ends in
a quite different sense from that in which the individual parts of the organism can be said to do so. But the zeal with
which one has sought within sociology for an analogy that could be laid as the foundation testifies to the significance
that analogy has for the development of science. This example also shows what demands are made of the analogies that
science can use. They must fit not merely at single points but at all points, so that one can draw the full consequence
from them, without the applicability to the domain that is to be illuminated ceasing. The symbol or image must be able to
become a scheme that one can confidently follow in the particular, in that it constantly leads to new, advancing
investigations. It is the good fortune of natural science that it has such a scheme in mechanical movement. For three
centuries that scheme has led natural science forward to new tasks and new discoveries. When speculative philosophy and
religious thinking have not led to such results, the reason must be sought in their not being able to remodel their symbols
into working schemes\footnote[1]{In a wider sense the schemes, even the numerical signs and the numerical ordering, are
also symbols. In this sense Leibniz used the word symbol when he set symbolic knowledge in opposition to intuitive
knowledge. (\textit{Meditationes de cognitione.} 1684). In the distinction between scheme and symbol I follow Kant
(\textit{Kritik der Urtheilskraft.} § 59).}.
And the reason for this must in turn be sought in their not,
% --- p. 41 ---
\setcounter{footnote}{0}%
\opage{41}like the knowledge hitherto discussed, seeking to explain one part of existence by means of another part, but to
cast light upon the whole of existence by means of what can be known about a single part.

\section*{III}

If speculative or metaphysical knowledge is to have any scientific significance, its task must be to work out a world-view,
a view of existence or of its innermost ground, such that it constantly builds on what empirical science teaches. Now
existence as a whole is never given, and cannot be thought of as given to any knowledge. The investigation that can be in
question here can then only aim at whether any of our concepts won through experience might be suited to be laid as the
foundation of a conception of existence as a whole. Our concepts are now all formed on the basis of determinate
experiences, and thus each for itself expresses a part or a side of existence. The question thus becomes whether existence
as a whole can be understood and expressed by means of a single part or a single side. It is not the parts that are here to
illuminate one another, but the single part that is to illuminate the whole. The problem of the possibility of knowledge
becomes here considerably sharper.

Where we pass from one domain of experience to another, it will be possible to specify the changes that the concepts must
undergo when they are transferred from the one domain to the other. But the change that a concept must undergo when it is
to be transferred from part to whole must be quite otherwise far-reaching than when the transference takes place from part
to part. The question is whether a transference from part to whole is possible at all, since the whole may have other laws
than those that hold for any single part considered by itself alone. The whole may moreover have sides or parts that do not
come forward in our experience at all. Perhaps we know only some of the elements of which existence as a whole is
composed. In any case it is here impossible to verify the use of analogy, since the whole is never given. Therefore it is
no wonder that thinking here comes to no conclusion. Neither
% --- p. 42 ---
\setcounter{footnote}{0}%
\opage{42}matter nor mind, neither multiplicity nor unity, neither being nor becoming are suited to be final expressions of
existence; it constantly bursts anew, through the richness of experiences, the final form in which one wants to gather their
content. And yet there is a constant need in the human mind to go to the limit, --- to apply its chief thoughts as
expressions for existence considered as a whole. This need will hardly be able to be eradicated; but it must be joined with
the critical consciousness that we here constantly apply analogies that cannot be carried through, symbols that cannot
become scientific schemes.

Among the attempts in thought that can be made here, philosophical idealism has a peculiar significance. It appears under a
twofold form, which it is of interest to examine, since the one form (which I shall call speculative idealism) will not
concede that it builds on analogy. This form of idealism holds it a necessity that whatever exists, whatever has reality,
must be at once unity and multiplicity; the greater the fullness and the more intimate the unity present in a given case,
the higher a degree of reality do we stand before. In this standard for true reality the very essence of the Absolute
proclaims itself. In a significant and profound manner this idea has in most recent times been developed by Francis
Herbert Bradley. With regard to this attempt in thought I have two remarks to make. First, it is not constructive or
properly speculative idealism that is presented here. It is not a system but a standard, an ideal, that is developed. We are
reminded more of Kant and Maimon than of Hegel. Second, this very standard, this ideal, is taken from the human mind. In
our mind we constantly find a given content whose multiplicity is to a certain degree worked together into unity, and the
greater the fullness is without harm to unity and concentration, the higher a degree of development do we ascribe to the
life of the mind. It is this standard which that idealism extends from the life of the mind so as to hold for all reality. It is
in reality an analogy that lies at the basis of this extension: the assumption is that what holds for the part of existence
immediately accessible to us must hold for the whole of
% --- p. 43 ---
\setcounter{footnote}{0}%
\opage{43}existence. We are thereby led over to the other form of idealism, which from the very first acknowledges that it
builds on analogy, and which I shall call metaphysical idealism. The most important representatives of this standpoint are
Leibniz and Lotze. Leibniz is the one who first with clarity and energy declared that every idealistic world-view is based
on an inference from analogy. From the standpoint of critical philosophy the question now becomes whether it is compellingly
necessary to make this inference from analogy. It might seem as though one were bound to choose between this possibility and
materialism, between conceiving the material in analogy with the mental or the mental in analogy with the material. In a
certain sense this is also so. But this boundness perhaps hangs together with our knowing by experience only two forms of
existence, the mental and the material. There obtains here, however, nothing contradictory but only a contrary relation of
opposition, and the possibility is not excluded that existence (as Spinoza conceived it) has far more basic forms
(attributes) than the two our experience shows us. On this ground alone the question must constantly remain open, quite
apart from the difficulties that in any case lie in the transference of a concept from part to whole. Even if existence as
a whole could be expressed by a single part or a single side, it would not follow that it was one of the two we know that
was suited to this.

Whereas the second kind of idealism has an advantage in the clearly conscious way in which it refers to analogy as its
foundation, the first kind distinguishes itself by specifying a standard of reality and thereby also a method for research.
The scientific significance of idealism consists in its, by finding the standard of reality in the connection of
multiplicity and unity, setting a work in motion for the extension of experience and for the working together of what is
found, uninterrupted analysis and synthesis in intimate interaction. Its practical or ethical significance agrees with this;
it demands both the extension and the elevation of life, proclaims a Plus ultra! and an Excelsior! at once.

% --- p. 44 ---
\setcounter{footnote}{0}%
\opage{44}Metaphysics is perhaps more an art than a science. It is then the task of philosophy, by setting up a standard such
as the one mentioned, to exercise criticism of the spontaneously developed world-views and thereby bring them into
connection with the properly scientific work. In any case there shows itself here a likeness between the highest points of
view of philosophy and the principles of exact science: in both places the value and justification consist in this, that
the principles set up are able to initiate a fruitful work of thought, lead to new tasks and new solutions.

\section*{IV}

Religion too works in its sphere of ideas with analogies. It maintains in its own way, like metaphysics, a conception of
existence as a whole. But whereas the metaphysician is guided by the intellectual interest of winning a final thought and
leaves it to the consistency of thinking and experience what the content of this thought shall be, the essential interest
in religion is what fate the valuable (wherever one then finds it) meets in the great struggle of existence. Religion is
conditioned by the experience that values arise and values perish, and the religious problem consists in the question
whether it is possible to hold fast a continuity under this constant change and beyond the limitation to which our
experience is subject. Where the religious consciousness is original and living, it is uninterruptedly occupied with
gathering its experiences and finding a satisfactory expression for them. It maintains that in the values that are
experienced, in the qualities of feeling, just as well as in the qualities of sense, a great order of things reveals
itself, and it seeks to form this conviction in definite thoughts. For the religious consciousness values cannot stand as
sporadic flashes, as by-products in the great process of existence. But in the course of its development the religious
consciousness becomes clear that what persists cannot be the single, empirical values; their time may be past, and yet ---
this is the religious faith --- the realm of values persists. It costs a hard struggle before this knowledge breaks through.
People let go only reluctantly of
% --- p. 45 ---
\setcounter{footnote}{0}%
\opage{45}the experiences lying nearest. If symbols are to be used, those used by preference are the most immediate and
nearby facts. Above all the religious consciousness would like to be spared the concession that all its ideas are symbolic.
Dogmatism has here its proper home; but its features are entirely similar to those it shows where it expresses itself in
the scientific domain.

Religion, like metaphysics, is compelled to express what is to be an absolute whole by ideas that are taken from a single
part or side. Its difficulties are in this respect of a similar kind to those of metaphysics. But the struggle between
thought and image becomes more far-reaching in religion than in metaphysics. The metaphysician struggles to make the image
or symbol into a scientifically usable scheme. This interest does not stand in the first line in religion. What matters are
great, gripping images that can express the mood during the struggle of life and keep up courage and hope. The possibility
of carrying the symbol through consistently recedes, and it is as a rule a more or less hostile criticism that examines it.
When the image merely expresses an essential side of the experience of life, it is held fast, and the inconsistency is
kept away by excluding comparison as long as possible.

The future of the religious problem depends on whether it is possible to hold fast the conviction of the continuity of
values under advancing experience and under clear awareness that every formulation becomes symbolic in a pronounced sense.
However it may go with this, the religious problem presents an analogy with the other chief problems of the human mind.
And if one is right that the properly religious axiom is the proposition of the persistence of value, then it also shows
itself here that the significance of principles depends on the mental work they call forth. That axiom can namely be
maintained only through a constant striving to discover and produce values in the world of reality. The religious question
therefore becomes in the end an ethical question. It is
% --- p. 46 ---
\setcounter{footnote}{0}%
\opage{46}in religion no more than in science the single images or symbols, but the use that is made of them, that is the
essential. There is a significant analogy between religion and science, where they become clear about their essence,
namely in the use that each for itself makes of the analogies with which they are now once compelled to work.

% --- p. 47 ---
\setcounter{footnote}{0}%
\opage{47}\essay[Preface and Postscript to My Philosophy of Religion (1905)]{PREFACE AND POSTSCRIPT TO MY\\ PHILOSOPHY OF RELIGION}
\begin{center}(1905).\end{center}
\streg

A discussion of the religious question, not merely in its relation to science but also in its relation to personal life, cannot avoid being colored by the personality that conducts the discussion. In my presentation of the philosophy of religion I make no secret of this either. The personal element, which necessarily asserts itself, may work both to the benefit of the clarity and thoroughness of the treatment of the problem and to its detriment. Which of the two holds with regard to my attempt, I naturally cannot decide myself. But I may perhaps give a small contribution to the decision by looking back on my personal and scientific course of development.


I have sometimes seen the remark that my present attitude to the religious question finds its explanation in the fact that in my young years I studied theology, and that something from that time still sits in my blood, which makes me stand somewhat differently on this question from many other freethinkers. For myself it stands thus: however sharply I feel the opposition between my present views and those on which I found sustenance in my youth, the continuity in my development nevertheless strikes me most. I have the
% --- p. 48 ---
\setcounter{footnote}{0}%
\opage{48}conviction that in the course of the years I have won fuller and more inward satisfaction of a personal need that in my young years I sought to have fulfilled by entirely different ways than now. What I believe I have preserved from those days is the sense of the significance of the inner life and the need to win expression for the longing and mood of the mind in the course of life --- to get the value of life and the fate of life expressed in great and gripping thoughts or images. It was this need that in its time led me to hold on as long as possible to the presuppositions of positive religion and to let them fall only step by step. I saw at that time in Christianity, as the Church presented and represented it, the only form in which such a need could be fulfilled.

It was the reading of Søren Kierkegaard's writings that first shook this assumption. It became clear to me that it was a great self-deception when the Church and ``modern Christianity'' thought they had reached a harmony between the life of personality demanded in original Christianity (``the Christianity of the New Testament'') and life as it is led in family and state, in art and in science. Theology had sought to throw a bridge across here, but this attempt rested on a great misunderstanding.

Heavy and dark were the years that now followed. I do not look back on my youth as a happy time. With the doubt as to where I should turn with my need to understand myself there joined itself the doubt where I should find a field of work. My work could not procure me satisfaction if it could not become of essential significance for my personality. But the break with theology led me into a duality, a dualism between the inner and the outer life, which in the long run was unendurable. It was life itself that led me beyond this duality by teaching me from my own experience the value of purely human relationships. As I had, under the influence of Kierkegaard's problems, flung my theological books to the floor in impatience at their unclear blending, so it was now Kierkegaard's books that I cast from me, because I could not reconcile their content with what life now taught me. I came into a living relation to life
% --- p. 49 ---
\setcounter{footnote}{0}%
\opage{49}outside the inner chamber. And once the way was open, one step brought the next with it. Soon science too became for me a reality that demanded its independent place in every serious settling of life's account. The study of Greek thinking in those years was of great importance for me in this; it taught me to see the possibility and value of a free and human understanding of life and conduct of life. I now had to find, in life and in actual experience, a possibility for the satisfaction of the need that had led me first to theology, then to the Kierkegaardian dualism. I saw that the way must be long, and from the beginning I have not had great faith that I would get very far along it; but I held fast to the assurance that here was a problem that could not be pushed aside.

In 1870 I defended my dissertation for the philosophical doctorate (with a treatise on ``Den antike Opfattelse af Menneskets Villie''), and in the short autobiography that was printed on that occasion in a university program I expressed myself in the following way about my course of development and my endeavor:

``As my future study I had in my school days for some time thought of philology, on account of my interest in the ancient languages and history; yet the need to develop and clarify the simultaneously awakening religious sense became the strongest, so that with desire and zeal and with the end of one day entering the service of the Church I took up the theological study. But there soon arose a doubt in me, which hindered the satisfaction I had expected to find in this study, both in the religious and in the scientific respect, so that it had a depressing and inhibiting effect instead of a clarifying and liberating one. The study of S. Kierkegaard's writings and of philosophy, so far as I could then pursue it with the time and ability at my disposal, gradually opened my eyes to what was justified in the criticism that was directed with such great emphasis against theology from the philosophical chairs of the University\footnote[1]{Rasmus Nielsen's and Hans Brøchner's.}. I concluded my theological course in the summer of 1865 with the conviction that I must seek by other ways
% --- p. 50 ---
\setcounter{footnote}{0}%
\opage{50}the truth I had not found there. I had given up the thought of entering clerical activity; I now occupied myself especially with history, literature and the ancient languages, in which subjects I also worked as a teacher. Classical antiquity came forward for me in a fresher and clearer light, precisely through the opposition to the scholastic period that I seemed to myself to have lived through; I was truly glad to read its historians, its poets and soon especially its philosophers. The deep conviction I thereby won of the independent significance of science and of the humane life, together with the fact that I still saw the true expression of my religious faith in the ecclesiastical dogmas, led me to believe that I had found a solution to my doubts in a view which just then came forward with greater energy before the public, but had met me from my first student years, and which sharply distinguished between religion and science as two independent realms\footnote[1]{Rasmus Nielsen's doctrine of faith and knowledge.}. I have always conceived this view more from the skeptical, negative side, more as a means to a personal solution of the difficulties, than as a positive doctrine of its own. A little work that I published at that time, ``Filosofi og Teologi'' (1866), sought from this standpoint to determine the relation of that view to those that had earlier appeared in our literature. What has later led me away from it is partly the contributions, some of them very thorough, that appeared in the further discussion of the matter, partly that the new view in its further development seemed ever more to turn from a skeptical or critical one into a positive, dogmatic one. To this came now also my own, gradually more independent study of philosophy. The need for unity and harmony in view of life and world-view grew together with confidence in the power of the human mind to attain such. Hereby there is indeed as yet only an ideal set, but an ideal that precisely must spur to ever renewed research.''

\stjerner

Of essential importance for me was later the study of English philosophy. The poles between which my development has
% --- p. 51 ---
\setcounter{footnote}{0}%
\opage{51}swung are Kierkegaard's doctrine of personality and French-English positivism. Brøchner, with whom I stood personally close in his last years of life, had moreover drawn my attention to Spinoza, and I remember how the reading of Spinoza's letters on a walking tour in Jutland in the summer of 1872 gripped me strongly and gave me a central standpoint that I have since held fast.

That I have succeeded, after my years of apprenticeship, in finding a new foothold I dare not say. But I have found a foothold where I myself can stand, and from which I can perhaps help others to find where they should stand. I have won the conviction that even the highest and noblest in the world of the mind is subject to the laws of a great natural order, which it is the task of science to find out, and that what is truly valuable is not violated but gains by having its firm roots in the connection of reality. Even the most beautiful tones are not too good to be expressed in notes. But my score is not finished, and presumably never will be. I can hear more and more beautiful tones than I can express by my notation, and I find in experience many signs that I cannot decipher, but that possibly are a notation for tones I cannot hear. That does not shake my conviction, and above all it does not shake my method. The problems have become greater and sharper for me in my wanderings in the world of thought. Some call that my skepticism or my negativity. I myself find it natural that the sky gets higher the higher one climbs on one's mountain ascent.

\stjerner

For my conception of the religious problem my attempt to present psychology and ethics on the ground of experience, setting aside speculative and religious presuppositions, became of essential importance. Thereby my knowledge was sharpened of the necessity of psychical equivalents, i.e. of new spiritual forms of life that could take the place of those that vanish, and bring compensation for them, so that life in its development does not lose its value (perhaps also its energy). I found now further --- in the studies that prepared the publication of
% --- p. 52 ---
\setcounter{footnote}{0}%
\opage{52}my work on the history of modern philosophy --- that nearly all eminent philosophers since the great critical and revolutionary breakthrough at the end of the 18th century had been attentive to this point. In a lecture that I gave in the spring of 1895 in the Students' Association of Kristiania, and which is printed in the first part of my ``Mindre Arbejder'' under the title: ``Kampen mellem gamle og nye Tanker. Et Tilbageblik'', I sought to show how the problem of psychical equivalents had constantly stirred in the thinking of the last century. In a university address (1897) on the religious problem (which is likewise to be found in the first part of ``Mindre Arbejder'') I stressed that the decisive question for the future of the religious problem is whether the personal view of life shall and can be grounded solely on what can be derived from the scientific results, or whether elements and tendencies must not always assert themselves in the final settling of accounts about the significance of life that cannot be justified by scientific means, even though they do not conflict with the results of science.

These studies and reflections led me to the conception that I have presented in my Religionsfilosofi, which appeared in the year 1901.


Philosophy of religion has no other methods at its disposal than the general philosophical ones. Its method is a combination of epistemological, psychological and ethical investigation. --- For the epistemological consideration it appears more and more that religion does not arise from purely intellectual motives, and that religious ideas cannot be maintained by purely intellectual means. If religious ideas are still to have value, it must be as figurative expressions of an experience of life and a need of life of a more special and personal kind than the experience on which science builds. --- In the psychological section it is shown how the religious feeling gets its natural place in the life of consciousness, when it is conceived as conditioned by the experiences that are made or are believed to have been made concerning the relation between life value and reality. It is the relation of tension between value and reality that again and again poses the problem and will continue to pose it, even if the time of the positive religions should
% --- p. 53 ---
\setcounter{footnote}{0}%
\opage{53}be past, indeed even if one no longer wants to use the word religion. A historical investigation of the deepest motives to which the religious consciousness appeals proves the proposition that the religious need is a need to maintain the persistence of value. --- In an ethical section the value of religion itself is finally discussed, and it is shown that the religious need has the significance of leading to the discovery of new values and to the maintenance of those already found. Although religion is not a necessary presupposition for ethical conduct of life, it can nevertheless supplement and continue this. --- The result is that religion is a form of spiritual life that --- if the power of life is not to be weakened --- cannot disappear without a new and equivalent form of life being developed. The great question is whether this is possible. My consideration turns partly against secure dogmatists, partly against self-satisfied freethinkers. ``Prophete rechts --- Prophete links!'' The problem lies midway between them, and they both shun it, each in his own way.

\stjerner

In thinkers with a theological past this past may work later on in different ways. --- There may be a mere echo, wistful but weak, and ever weaker, an echo that can hinder one from drawing consequences that would otherwise stand as unavoidable. --- A contrast effect may assert itself, which manifests itself in one's later being unable to find words strong enough to express the opposition, the antipathy, perhaps the repugnance. This type sometimes reminds one of schoolboys who enjoy their freedom by jeering at the school they --- not always equally thoroughly --- have gone through. It is an entanglement, just as much as the echo. --- A third type may be designated as that of critical continuity. It finds a more deeply lying connection behind the strong opposition by going back to that which in different forms must form and always has formed the foundation of every serious understanding of life and its fates. Even if the ideas in which men in earlier times expressed this understanding have for us now at most poetic significance (that is not even always the case), there may nevertheless be a kinship present, and
% --- p. 54 ---
\setcounter{footnote}{0}%
\opage{54}it is in any case the task of the psychology of religion to investigate how matters stand in this respect.

It is naturally difficult for the individual to decide to which of these types he himself belongs. But I have no doubt which of them is the most fruitful and the most justified. It must be the one that, without being ensnared by echo or contrast effect, uses the experience that has been won by living in an opposite camp to illuminate the problem as it presents itself to us now and for the time being must present itself. --- It is of no small interest that investigations in this direction are under way throughout the civilized world. (Cf. below the treatise ``On Some Recent Works in the Philosophy of Religion''). It is likewise not without interest to see how the religious problem, when it is conceived in the manner here indicated, proves to be analogous to the other great problems that occupy human thinking. It is indeed probable that what becomes a problem for human thought in the different domains must present analogous features (See on this my University program of Nov. 1902: Filosofiske Problemer). ---

Whether my work can be claimed by one side or the other is rather a matter of indifference to me. It testifies, however, either to great undemandingness or to naive calculation when even from the Catholic side my book has been presented as a favorable sign of the times. I myself find the bridge ever more decisively broken between me and every kind of dogmatic religion. But perhaps one is from the theological side so accustomed to find religious questions treated with scorn and mockery or with schoolboy arguments that the mere fact that the problem is discussed with a certain seriousness may give the impression that a rapprochement has taken place. Yet not all theologians have conceived my presentation thus. In ``Kristeligt Dagblad'' there appeared some years ago a sharp but chivalrous critique of my views. The critic ended by testifying to his opponent his personal esteem. The then chairman of ``Inner Mission'' reproached him for this, declaring that I was a man for whom no Christian could have personal esteem, --- and the then
% --- p. 55 ---
\setcounter{footnote}{0}%
\opage{55}editorial board of ``Kristeligt Dagblad'' associated itself in a note with this declaration. Such things reassure, when one sometimes feels uneasy at being claimed by a certain side.

\stjerner

I have tried to summarize the essence and significance of religion in a single proposition, which I have sought to ground psychologically and historically: the proposition of the persistence of value. From various quarters I have met with objections to this hypothesis. (I have not published it as more than a hypothesis).

It has been said that such a general proposition cannot express the essence of religion, since religion is something living and concrete. This objection has always seemed to me somewhat childish. Every general proposition can only be a concise expression of what occurs in manifold different forms in reality. None of these manifold forms can be immediately expressed in the general proposition, but its significance must consist in providing a clue for finding that in which the individual phenomena, despite their difference, resemble one another. The general proposition in which we define a vertebrate cannot contain anything about the lion's roar or about man's speech, but it can nonetheless very well hold both for the lion and for man. The meaning of that proposition in the philosophy of religion is naturally not that the religions themselves speak directly of values and of their persistence, but that in all religion there manifests itself a need that in our psychological theory can be designated as a need to assume a persistence of the valuable that man at the given stage of development knows. I have not wanted to reduce religion to abstract metaphysics, but I have spoken as a philosopher about something that is other and more than philosophy. One could with the same right reproach the physiologist that he puts an abstract theory in the place of real breathing.


Next, the essential in religion has been found at other points than that on which I lay the weight.

From a theological or at least theistic standpoint it has been said
% --- p. 56 ---
\setcounter{footnote}{0}%
\opage{56}that the essential in all religion is that man feels himself in relation to one or more higher personal beings: religion consists in a relation between an I and a Thou! --- Against this assertion I ask (and have already in my book asked): What is it that makes the gods divine, and what need is it that leads to faith in them? If the gods were not assumed to maintain what men recognize as good and great, would men then worship them? Even where the gods are worshiped merely as the mightiest beings that are known, there lies a need at the base to hold fast to a continued persistence of that which man places highest (whatever this may be). When religion has hitherto predominantly been faith in personal beings, this is due for a great part to the influence of animism, i.e. the tendency to conceive the events of the world as effects of the intervention of personal beings, a tendency that manifests itself very early. Religious ideas are very often transformations of ideas that have developed from other than religious motives. In the faith in the transmigration of souls, which is peculiar to many religions, there is an example of the same. Another example is offered by the manner in which Greek philosophy was used by Christian theology. The question is whether in such ideas we have the essential in religion, or whether the essential must not be the need that is felt to take up and shape them, a need that under other conditions could lead to erecting entirely different spiritual edifices.

From the standpoint of the history of religion the essence of religion has been found in its being a kind of magic, a mediate magic, which arises when man gives up the hope of immediately changing the course of nature according to his wishes. By prayers and sacrifices he now believes he can induce mightier beings (whether these are conceived as personal or not) to intervene and bend the course of nature, so that he can attain his wish. Against this conception I ask in a similar way as before: \textit{why} does man seek to win a power over nature that he cannot win by his own labor? Is it not because he experiences the limit of his own power and yet feels a need to win security and safety for that which lies nearest his
% --- p. 57 ---
\setcounter{footnote}{0}%
\opage{57}heart (whatever this may be)? Here again this need shows itself as the decisive thing, and to it one must therefore go back when one wants to trace the innermost forces in the religious life. Only thereby will one be able to get a basis for the discussion of the religious problem in its whole extent. The need can persist although the forms change, and neither magic nor animism need to all eternity be bound to the essence of religion. Should the time of religion be past, because it is declared that love is better than sacrifices, or because no image any longer stands as wholly satisfying? --- The opposition between old and new forms therefore does not become meaningless. At a given time it can call forth struggles of life and death in the outer world and the greatest upheavals in the inner world of the mind. The philosophical investigation can concern only the \textit{possibility} of new spiritual forms of life; the actual production of them is a matter for life.

% --- p. 58 ---
\setcounter{footnote}{0}%
\opage{58}\essay[On Some Recent Works in the Philosophy of Religion (1903)]{ON SOME RECENT WORKS\\ IN THE PHILOSOPHY OF RELIGION.}
\begin{center}(1903).\end{center}
\streg

\section*{INTRODUCTION.}

By philosophy of religion I understand not an attempt to put speculation in the place of religion, still less an attempt to ground religious assumptions by the road of speculation, but an attempt to throw light on the position of religion in the intellectual life of the age.
Religion is here the object of investigation as an actually given form of spiritual life, whose nature, significance and justification it is our task to determine. If one now asks what points of view philosophy commands for use in such an investigation, there can surely be talk of only three points of view, --- but these are really also all the points of view that philosophy commands at all: the logical (or epistemological), the psychological and the ethical. Since every religion --- whether one takes the word in a wider or a narrower sense --- finds expression in certain ideas, the question naturally arises of the validity of these ideas, their value as knowledge. Do they contain a consistent and valid knowledge of existence, and can they be brought into agreement with the ideas that science by its own roads has arrived at holding fast as valid? Even if such an \textit{epistemological} investigation should lead to the result that the
% --- p. 59 ---
\setcounter{footnote}{0}%
\opage{59}significance of religious ideas cannot consist in procuring for us a scientific understanding of existence, the question of their value and justification is not thereby settled. It may rest on their relation to other sides of mental life than the purely intellectual. And direct observation also shows, indeed, that feeling, need and will play quite as great a role in religion as knowledge. There is therefore every occasion to undertake a \textit{psychological} investigation of the religious life, as observation teaches us to know it in the present and in historical tradition in the past. This investigation aims, first, at finding the constantly operative motives in religion and separating them from the changing forms and connections in which they appear under different conditions. Only in these motives can the essential in religion be sought, --- that which, so long as it possesses vitality, will always be able to express itself in new shapes. There exists no religion in general; but experience shows us certain fundamental motives in a wealth of different connections or variations. Therefore it is, second, the task of the psychology of religion to seek out and describe the most important types of religion and religious life. If, then, motives and types have thus been demonstrated, it is the business of the \textit{ethical} investigation to decide what value religion, as a peculiar form of human spiritual life, has for human life in its whole, great connection. The standard here is the same that we also apply otherwise when we discuss the value of a human quality or action. We examine their relation to certain fundamental values, in the last instance to a single fundamental value which for us is the decisive one. What this fundamental value is, it is the business of ethics to find out. If there is disagreement about the fundamental value that determines what we in the individual cases call valuable, there will also be disagreement about the valuation of religion. The question would of course be decided in advance if one assumed that religion (i.e. the religion one has or believes oneself to have) is itself the fundamental value by which everything is to be measured. But then there would likewise exist no religious problem, nor would there be any possibility of or need for a philosophy of religion, just as little as there would be any possibility of or need for a theory of knowledge, if
% --- p. 60 ---
\setcounter{footnote}{0}%
\opage{60}one assumed that some religion or other contained the absolute truth, by which everything else one would call true was to be measured. A problem arises here, on the other hand, as soon as one becomes aware that in acknowledging religion one presupposes a certain standard as valid, a standard that is not drawn from religion alone, and especially when one discovers that in the religious domain one applies a standard that differs from that which one applies in other domains.

Of the three points of view from which the religious problem can and must be treated, it is the psychological that in the most recent times especially attracts attention. It is of course nothing new that religion is conceived as a psychological phenomenon. Every more profoundly reaching religious movement had to lead to it, especially by making it clear that what matters is not merely to stand in a relation of ideas to the Divine, but also to be moved by it in inward feeling. More and more it is pointed out from the religious side that the ultimate foundation of religion is human experience of life. The psychology of religion fully acknowledges this, and it then sets itself the task of describing this experience.
It proceeds here from the view that, if the appeal to experience is to have any meaning, one must understand by experience immediate experience, and it then matters to determine what and how much can be immediately experienced. It is necessary to distinguish between what is really immediately given and what is added as explanation. The explanation is drawn from a circle of ideas that lies close at hand to man. And from the explanation the mode of expression must again be distinguished. But these three things, experience, mode of expression and explanation, it is extraordinarily difficult to keep apart from one another in the case of religious experience, which is essentially an experience of feeling. For the psychology of religion it therefore matters to seek out simple examples where the difference is clear, and with their help to explain to oneself the more complicated cases. But so much will be clear at once, that there is a long way from an immediate religious experience to the often very subtle dogmas or to the poetically shaped myths and legends that the historically given religions retain in their
% --- p. 61 ---
\setcounter{footnote}{0}%
\opage{61}traditions. And a quite special distinction must be drawn between experiences that are evoked by such traditional ideas in the attempt to convert them into self-won experience, and experiences that at first hand lead to the formation of certain determinate ideas. Only this latter kind is original religious experience. And even with it one must still ask whether the one and the same experience could not be expressed and explained by other ideas and in other forms than is actually the case. Only if such a counter-test could be made, and if it then led to a negative result, would the proof from experience have been furnished.

I believe that this psychological method will acquire very great importance in the continuing religious discussion. It is with full right that appeal is made from the religious side to personal experience of life as the last criterion against the often very superficial appeal to the supposed results of reason or science. But if one does not treat the very concept of experience with strict criticism, distinguishing sharply between the given and its interpretation, one introduces a dogmatism that is far worse than the old rationalistic dogmatism on which one loftily looks down, as surely as it is worse to deceive oneself and others with regard to the most inward relation of one's life than to commit logical errors in an intricate piece of reasoning. In any case psychology cannot allow itself to be prevented from attempting to analyze the phenomena, immediate and uncompounded at first glance, that are appealed to. There lies here a great material ready to hand in the utterances of the religious consciousness at different times and under highly different cultural conditions, attractive and refreshing especially for those who have had enough of the hitherto current way of treating the religious questions, both among the assailants and among the defenders.

It is this that I have sought to the best of my ability to impress and carry through in my Philosophy of Religion, whose most extensive section is precisely psychological. I have also had the satisfaction that a theological critic (Pastor L. P. Larsen in Studenterhjemmets Julebog 1901) has declared that my method is the only right one. But my astonishment must then be all the greater when I see that this same critic disputes my results --- by means of
% --- p. 62 ---
\setcounter{footnote}{0}%
\opage{62}a quite different method than that which he has thus acknowledged! For he uses the concept of experience precisely in the vague and wide sense in which it is often used in everyday speech. He puts the question throughout thus: do certain traditional dogmas have value for life? But a consistently carried-out investigation in the psychology of religion goes also the opposite way and asks: do life values really disappear because one does not accept certain dogmas, or because one puts poetic symbols in the place of the dogmas? I myself dare (although I do not hide my own personal standpoint in this matter) not give a definite answer to this question; I have expressly emphasized the possibility that there might here be a loss of value. This possibility no science can refute, and therefore the assumption of the persistence of value rests on a faith, even if one does not extend this assumption beyond the limits of the human world. ---

In a domain where different tendencies still clash --- and perhaps for all eternity --- it will be of great importance to investigate whether the procedure employed is the only one that can come into question. Therefore I have wished to bring forward a series of attempts in the philosophy of religion, most of which first appeared after the completion of my book, in order to see whether there are other points of view to be adopted than the three mentioned (the logical, the psychological and the ethical), and specifically to see whether the attempted determination of religious experience and of the fundamental idea of religion should need supplementing or correcting.

Not only as regards the method but also as regards the results, such a comparison with the contemporary works of other investigators may be of importance. I therefore permit myself to recall that the results at which I believe I have arrived are mainly the following three. --- The religious problem falls under the general problem of valuation, and the method for its treatment will therefore more and more become ethical valuation alongside psychological description. Neither of these two procedures can be dispensed with, since one can of course only value accurately what is accurately described. --- The religious need may be characterized as a need to assert the persistence of value, and religious experience consists of what man
% --- p. 63 ---
\setcounter{footnote}{0}%
\opage{63}has experienced concerning the relation between the values he knows and the reality he knows. The special religious assumptions indicate --- whether they appear in the form of myth, dogma, legend or symbol --- the ways in which men have sought partly to express, partly to explain what they have thus experienced. --- In the conditions of intellectual life in modern times there is nothing to prevent such a need from continuing to express itself and to find satisfaction, even if it must more and more happen in other forms than in earlier times, --- and even if it is not necessary that it express itself in everyone. It is a hypothesis of the history of religion that the proposition of the persistence of value is the fundamental religious idea, i.e. that it is this which expresses the pervading tendency in the historically given forms of religion. But the proposition of the persistence of value can also be set up as a philosophical hypothesis. It denotes a view of life that may claim to be \textit{compatible with} a strictly carried-through scientific world-view, even if it cannot be \textit{derived} from it. The view of life stands in everyone with a certain independence in relation to the world-view. There takes place a constant influence of the latter on the former; but never is the view of life a mere product of the scientifically grounded world-view. For the view of life contains a plus which derives from the personality's valuation of the existence it seeks to understand with the help of science. Now there may indeed be found within valuation greater or lesser consistency, and the points of view that are laid at the basis of valuation may be more or less natural and essential. But in the domain of valuation we are to a far higher degree compelled to be empiricists than where understanding and explanation are concerned. The irrational (in the mathematical sense), which our thinking must always leave standing over, has a greater extent in the domain of valuation than in that of pure theory. Beyond the \textit{possibility} of a view of life that builds on the fundamental idea of the persistence of value, philosophy cannot lead us. The fundamental motive of every view of life springs from practical experience of the fate of values in the world of reality; and if this experience leads to faith in the persistence of value, then
% --- p. 64 ---
\setcounter{footnote}{0}%
\opage{64}this faith will not need to stand as a mystical hypothesis, a mere object of passive contemplation, but will naturally become a working hypothesis, in that it spurs one to find new values when the time of the old is past, or perhaps even oneself to produce such. It then gives courage to work both in the domain of thought and in that of reality in order to win progressive confirmation. Perhaps not all individuals will be led by their experiences to harbor such a faith; it may suffice if it maintains itself in the race. For the life of the race, in order to be rich and vigorous, needs many different individualities, perhaps even a constant struggle between spiritual tendencies. Many individuals may be so absorbed by the special work of thought and of reality that they do not have time, strength and interest left over to form any view of life expressed in determinate thoughts. Here, as in other respects, a division of labor may take place. Why should the division of labor be restricted to the purely external conditions of life? Is there not reason to assume that it intervenes in the most inward concerns of men? Life will become richer thereby, just as the world of art becomes richer the more the individual artist follows his individual need and capacity. The very proposition of the persistence of value perhaps even requires that not all believe in it or think of it.

These last remarks I lay particular weight on in face of an objection that has been directed against my book: that it presents a purely subjective conception. I have not hidden my own personal position toward the questions I discuss; but the main thing for me has been to investigate the possibilities that present themselves under the cultural conditions of the present in a religious respect, and here it was not enough to confine oneself to a single individual. It is a comparative investigation that I have wished to make. I have wished to undertake a comparison between the position of the religious in the past and in the present, and between its significance for different personalities and standpoints in the present. What matters in a philosophy of religion is not the philosopher's own faith, but the possibilities of faith that he finds in the intellectual life of the race, as it can now be carried on
% --- p. 65 ---
\setcounter{footnote}{0}%
\opage{65}in clarity and in depth. Less and less has a philosopher the right to assume that the psychology of the race is exhausted with his own individual psychology.

A related objection has been raised against my presentation by an energetic German scholar, whom I have the honor to be able to call my personal friend. In a detailed review of my book in the journal ``Das freie Wort'' (1902) FERDINAND TÖNNIES has criticized my conception of the religious problem in a manner that seems to me capable of illuminating the relation between the history of religion and the philosophy of religion.

``For Guyau,'' he says, ``and I confess that in this I side with him, religion is in the first instance an object of social psychology or sociology. For Høffding it belongs almost entirely to individual psychology, which indeed coincides with what one is accustomed to understand by psychology at all. He calls Buddhism and Christianity `the higher popular religions': in them the inner psychological drama is to appear as the real world drama\dots{} Now it is indeed unmistakable that in both these highly important religious systems the individual need and the individual thinking have acquired an overwhelming significance. But whence does this struggle of the soul derive? Is it an original and essential constituent of religion as popular belief and popular custom, or at any rate of the religious feeling that `bears' belief and custom? \textit{Or} does it, precisely as regards its purer and higher content, issue from another source, namely from a thinking, a philosophizing, that by its nature is alien to popular belief and popular custom,\dots{} although it has certainly attained in those two popular religions an adjustment with popular and wild-growing (naturwüchsige) ideas and feelings, indeed in Christianity has even concluded an eternal league with them?''

The consequence of this line of thought seems bound to be that Buddhism and Christianity are not really religions. Indeed, one would really have to go further. Is there any positive religion whatever in which elements of thinking cannot be demonstrated? The most elementary, not to say the most primitive religions build on animism, man's most elementary philosophy, and thus at all stages of
% --- p. 66 ---
\setcounter{footnote}{0}%
\opage{66}development a people's religion will bear the stamp of its antecedent ``philosophy''. Historians sometimes have rather romantic ideas of the immediacy and ``Naturwüchsigkeit'' with which men enter into their cult, myths and dogmas. There take place, as the history of religion itself shows, constant combinations, supplementations and metamorphoses, in which a more or less profound work of reflection co-operates. But what for me is the main point: it is misleading here to speak merely of thinking and philosophy. It is first and foremost personal experience of life that exerts its influence in these changes, which are often new formations. The personalities who have been decisive for the origin and development of the higher popular religions have first and foremost laid their own experiences of the valuable and its fate in the world at the basis of the changes they were led to make in what was handed down. No religion makes itself. Religion arises and exists as religion only in the minds and thought of men. It must have been certain determinate human beings who shaped cult, myth and dogma, just as it must have been certain determinate human beings who composed the ``popular poetry'', and just as ``society'' at all times consists of individual human beings. To that extent all psychology must be individualistic, even if it zealously applies the sociological method (alongside other psychological methods), which there is quite special occasion to do in the psychology of religion. --- At the end of his review Tönnies, moreover, after having repeated that in religion he sees especially an object for social psychology, says: ``Ich bekenne meinerseits, dass mich Høffding gelehrt hat, die Tragweite der individuellen Entscheidung besser als bisher zu würdigen''. ---

Only when one lays weight on the subjective element in religion does it become possible to find an inner continuity between the changing religions and views of life. And on this element rests also the significance for the continued life of the race that a religion can have, even though its time is once past. It is the great human treasures that the religious spirit has been able to raise in the inner life of men that constitute the most precious inheritance that the religions of the past leave behind. A poem like BUNYAN'S
% --- p. 67 ---
\setcounter{footnote}{0}%
\opage{67}``Pilgrims Progress'', which has been called the divina comedia of Protestantism, has thus acquired a lasting significance (at least for the English race) beyond the limits of the definite confession within which it came into being. And the reason is that it expresses what moved in the soul of a human being deeply gripped by one of life's oppositions. An excellent historian of literature has rightly said of it: ``All that is best and most characteristic in Bunyan's writings proceeds from the inner drama in which the actors were God, Satan and a lonely human soul. Institutions, churches, ordinances, rites and ceremonies help him little or not at all\dots{}. Let everyone faithfully hold fast to his most inward experiences in the great concerns of life, and his words will awaken manifold echoes in other souls; \textit{the deepest fellowship is found not in institutions or} \textit{corporations or} churches, \textit{but in the} lonely \textit{heart's} \textit{secret}.'' (DOWDEN: \textit{Puritan and Anglican}. London 1900.
p. 234). In my Philosophy of Religion I have also pointed out how the sense of the significance of the purely subjective elements in the religious domain can pave the way for a deeper feeling of community than that which is conditioned by union around common dogmas. To one and the same dogma or one and the same cult there may correspond highly different personal moods and states. This diversity is a proof of the independence of the subjective element. The relation between the objective and the subjective elements (or perhaps better: between the social and the individual elements) in religion is constantly varying, so long as there is life in religion. In the historically shaped religions it is the objective, social elements that most strongly strike the eye. And if the expression ``positive religion'' is to be understood in analogy with expressions like ``positive law'', what is meant thereby is a religion that appears as an institution, with a shaped cult and a fixed tradition. Psychologically viewed, however, what motivates man's adherence to and holding fast to what has thus been fixed, and what under certain conditions leads to making changes therein, must always be the essential thing. The significance of this subjective element can rise so that the objective element is as good as eliminated. We then come to
% --- p. 68 ---
\setcounter{footnote}{0}%
\opage{68}a limit where \textit{there} can be dispute whether the concept of religion can still be applied. It may easily become a dispute about words.
It may lead to raising profound questions like that which a man put: ``Is the star that the astronomers call Uranus really Uranus?'' --- If one lays the weight on the fundamental motives, not on the external marks, it will be justified to extend the concept of religion beyond the sphere of the so-called positive religions. Religion then becomes a generic concept, within which religion as positive religion, built on tradition and provided with fixed forms of cult, denotes only a particular species. Thereby the oppositions are not underestimated; for within one and the same generic concept the strongest oppositions can appear, just as in reality a struggle of life and death can be waged between species that belong to the same genus. ---

After these introductory remarks I now go on to discuss some works in the philosophy of religion from the most recent times in order thereby to give a critical supplement to my own attempt.
I shall choose such works as can be said both to be typical of the philosophy of religion of our time and to be instructive for those who are interested in the religious problem. --- There are, however, some authors whom it might in and of itself be natural to bring forward here, but whom I would rather discuss on another occasion, because they will best be understood in connection with the most recent philosophy in general. To these belong especially Guyau and Nietzsche, whom I shall in no long time have occasion to discuss in a work on modern philosophers\footnote[1]{See \textit{Moderne Filosofer} (Copenhagen 1904). p. 107---147.}, on which I am at work, and to which the present treatise can also be regarded as a supplement.

In the first section of the treatise I deal with some systematic attempts; in its second section I bring forward a series of more specialized investigations in the psychology of religion, especially from American literature.

% --- p. 69 ---
\setcounter{footnote}{0}%
\opage{69}

\section*{I.\\ SYSTEMATIC ATTEMPTS.}
\streg

\subsection*{1. HERMANN SIEBECK.}

SIEBECK (b. 1842), who works as Professor of Philosophy in Giessen, has made himself known by an excellent work on the history of psychology in antiquity and the Middle Ages and by his aesthetic studies. His chief work is the Philosophy of Religion that appeared in the year 1893. Originally proceeding from Herbart's school, Siebeck has later received influences from Lotze; but above all the effect of zealous Kant studies is felt in his chief work.

The task of the philosophy of religion is for Siebeck posed by the relation between religion and culture. Modern culture is inclined to regard itself as an end; religion it can, at best, acknowledge only as one of its imperfect predecessors. From the religious side, on the other hand, it is maintained that the life of culture --- like life and work in the world of experience in general --- can only be a means for a higher, ``supramundane'' end. The dispute that thus arises it is the business of the philosophy of religion to decide. It is therefore not a religiously attuned philosophy, like Neoplatonism in antiquity or like Schelling's later doctrine; nor is it, like the scholasticism of the Middle Ages, a presentation of a certain determinate religion in philosophical form. It regards religion in its various shapes as a fact in human life and investigates what significance this fact has for the development of human life. It seeks to determine in what domain of human spiritual life religion particularly belongs, and what influence it exerts from there on other sides of life. Philosophy here investigates religion in a similar manner to that in which in its other disciplines it investigates science, art, law and morality.

According to Siebeck it must essentially be an indirect method that the philosophy of religion can employ. It must test whether the very existence and development of the life of culture do not become worthless and
% --- p. 70 ---
\setcounter{footnote}{0}%
\opage{70}self-contradictory when it cuts itself off from the supplementary points of view and motives that religion affords. A positive determination of the higher world to which religion will open access, philosophy cannot give. Further than to demonstrate the necessity of a supplementation it cannot reach; but it can test whether the forms in which the religious consciousness has sought to express the Divine are suited to furnish the supplementation that a critique of cultural development leads one to demand. In advance Siebeck declares as his conviction that in the indicated sense of the word a demonstration can be given of ``the normality of the essential content of the religious consciousness''.

For the philosophy of religion, when its task is determined in this way, it will be natural to keep to the highest forms of religion, since only they become of significance for the relation between religion and culture in modern intellectual life. Siebeck distinguishes between three main forms of religion: natural religion, ethical religion (``Moralitätsreligion'') and religion of redemption. In the first form the gods are conceived as powers of nature, and the relation to them is determined essentially by the need of self-preservation. In the second form the gods have a calling to fulfill, and their worship is motivated by the dignity this calling gives them. The third form, the religion of redemption, which appears typically in Buddhism and in Christianity, is characterized by the significance that prophetic personalities acquire. Myth and cult, which rule in the earlier forms, here recede before the original personalities. There is in all independent personality, especially in the personality of genius, something irrational, something new, which does not find its full explanation in what precedes, and which perhaps only slowly finds recognition and adherence. But most of all this holds of the great personalities from whom the new value concepts proceed, whose validity is to be tested through long times in the life experiences of the race. To them applies \textit{the concept of revelation}. This concept is, according to Siebeck, itself a value concept, i.e. a concept whose validity does not rest on the manner in which its content can be demonstrated in reality, but on the practical significance of the content for life. This is seen by comparing different utterances of
% --- p. 71 ---
\setcounter{footnote}{0}%
\opage{71}Siebeck. In one place (p. 192) it says that the religious content that is proclaimed by the great personalities is called revelation on account of its ``supramundane'' character. And as regards the concept of ``the supramundane'', it says in an earlier place (p. 18): ``The concept of the supramundane can first arise on the basis of a world-view that not only takes the world with its existence and content as a given fact, but by virtue of a spontaneous mental activity independent of the given subjects this fact to a \textit{valuation}, the principle of which is drawn from the interior of the mind itself.''

It is really first within this third form of religion that religion and culture come into opposition with each other. Redemption, which in Buddhism is proclaimed in a more negative, in Christianity in a more positive manner, consists in the authority of personality over against the world; it is given with the appropriation of the thought that the world has value only on account of the existence of personal life. Personality has at once emerged from the world and is independent over against the world; it wishes not merely to understand the world and assert itself in it and over against it, but also to work in it for that which, from the standpoint of personal life, stands as the highest. Hence it comes that it falls so naturally to religious faith to conceive the ground of the world itself in analogy with the personal. For faith personality is the first object, whereas it is the very last object of science, because it is the most composite and at the same time most central subject-matter with which knowledge can occupy itself.

On this highest standpoint an opposition arises between faith and knowledge. By virtue of the value of certain phenomena the religious consciousness is inclined to place them outside the natural causal connection, whereas on the other hand science is inclined to declare itself to be the sufficient foundation for view of life and conduct of life. Siebeck maintains against the first assertion that belief in miracles is not a necessary element in religious faith. The clearer the religious consciousness becomes about itself, the more it will --- he holds --- see that belief in miracles does not glorify but precisely limits God, in that it sets nature as a barrier that is to be broken through
% --- p. 72 ---
\setcounter{footnote}{0}%
\opage{72}by God. Against science's assertion of being a sufficient foundation for life he maintains that this assertion itself can only be a matter of faith, since a scientific proof of it is not possible, and since the rejection of the necessity of religion takes place through an act of personality, which is motivated by the presupposition that with science itself the highest good is given. The very need for scientific research rests on this unprovable presupposition.

There is in all religion a subjective and an objective element. The subjective element, which one might call religiosity, expresses itself in mood, the objective element, religion as opposed to mere religiosity, is thing, work, performance, visible and tangible shape for the subjective mood. The depth of religiosity is the measure of value for the objective being of religion --- sufficiency and vitality; but the subjective mood would never have developed except under the influence of the objective elements, so that religiosity never in time precedes religion. There is a constant relation of reciprocal action between them. But the more a system of cult is organized, the more are new moods and new thoughts excluded from the religious life. It is cult that always makes the greatest resistance against religious renewal. The higher a stage of development religion reaches, the more pronounced individual variations will the religious life present, so that the great human differences can assert themselves not only in the relation between the religions, but also within one and the same religion. ---

This is the most essential content of the first division of Siebeck's work, the descriptive part, which treats of the essence and development of religion. As he now goes on to discuss the truth of religion, he declares that what is here at stake is not the truth of any single religion, but the validity of the religious principle in general. This validity can be grounded only by revelation --- in the sense of the word indicated above --- intervening in human culture and, although it does not itself spring from this, yet affording a necessary supplement to it.

At this point it is not quite easy to understand
% --- p. 73 ---
\setcounter{footnote}{0}%
\opage{73}Siebeck's line of thought. In his descriptive part (which, incidentally, is not merely descriptive) he has sharply distinguished between valuation and causal explanation, so that these need not come into opposition with each other, but denote different ways of viewing existence. But now he sets up the assertion that an ethical action, an action that aims at asserting and carrying through the valuable in the world of reality, becomes possible only if this is in the last instance subject to an ethical power, and he finds in this a ground for the concept of God. Every ethical action rests, he holds, on the presupposition that it is possible to realize the content of the ethical ideas within the connection of nature, and this presupposition is a religious presupposition. Of the ethical power that is here presupposed one can form an idea only by the way of analogy, but such a use of analogy Siebeck regards as justified when it avoids both the dualistic relation between God and world of ``deism'' and the abolition of every difference between them by ``pantheism''; thereby one arrives at a ``theism''. --- Here one misses in Siebeck a closer investigation both of the significance and justification of analogy and of the relation between valuation and causal explanation. As regards the first point, this is the more remarkable in that Siebeck himself earlier in his interesting treatise \textit{Die metaphysischen Systeme in ihrem gemeinsamen} \textit{Verhältnisse zur} \textit{Erfahrung} (in the second volume of Vierteljahrsschrift für wissenschaftliche Philosophie) has demonstrated the significance of analogy for every speculation that goes beyond experience. It remained now only to apply this consideration to the religious ideas. As regards the second point, he does indeed emphasize the difficulties of uniting the teleological and the causal way of viewing in full mutual harmony, but he holds that the religious consciousness does not regard the world as a closed totality, but as part of a higher connection of things, from which ever new creative forces can flow into it. I am not sure that the religious consciousness is here interpreted rightly; it has in any case very often a tendency precisely to regard ``the world'' as a whole for itself, and Siebeck has himself emphasized this in his critique of the
% --- p. 74 ---
\setcounter{footnote}{0}%
\opage{74}idea of miracles. It is rather for the scientific consciousness that the world does not form any closed totality. ``The world'' is not a scientific concept.

In thus acknowledging the justification of conceiving the natural causal connection as subordinate to an ethical principle, Siebeck does not overlook the objections that constantly arise against this way of viewing from the existence of pain and evil in an existence that is supposed to have its ground in such a principle. He does, to be sure, not see that this old objection must strike particularly sharply ``theism'', which conceives that principle in personal form. But he concedes that the difficulty is insurmountable for the religious consciousness, however much this consoles itself with the thought that evil must be possible in a world where free beings were to develop. He has given the chapter in his book in which this question is discussed the heading ``Theodice'' and remarks that theoretically it is God who is to be defended in an attempt at theodicy, but that what matters for the religious consciousness is to keep its hope up in spite of its trials. One might remark upon this that prior to all ``theodicy'' there goes a stage in which it is man who is to be justified before the tribunal of God. Then comes only later a stage in which God is summoned before the tribunal of human thought; then one begins to write theodicies. After that comes the third stage, where neither of these takes place, but where it clearly and distinctly concerns a struggle to keep hope up, and where the ideas one forms of what lies beyond experience expressly stand as symbols for this hope.

But when it thus turns out that the problems that science does not solve are not solved in religion either, indeed appear in it precisely in the sharpest form, then the supplementation of culture, the completion of the world-view, that the religious idea was to bring has not succeeded. The necessity of adopting this supplement cannot be demonstrated.

% --- p. 75 ---
\setcounter{footnote}{0}%
\opage{75}

\subsection*{2. AUGUSTE SABATIER.}

It is a sign of the times that also from the theological side one begins anew to take an interest in the philosophy of religion. This interest has otherwise been as good as vanished, after the hopes that had been cherished for romantic speculation were shattered. The Protestant-theological professor in Paris who died a few years ago, AUGUSTE SABATIER, has in his work \textit{Esquisse d'une philosophie} de la \textit{religion} \textit{d'après} \textit{la psychologie} \textit{et} l'histoire (5 éd. Paris 1898) set himself the task of treating the religious problem on a psychological and historical basis.

In psychological respect Sabatier lays the weight on this, that what is aimed at in religion is not a solution of ``the riddle of the world'' but of ``the riddle of life''. It is the enigmatic and contradictory in our own being that gives rise to the religious need, and world conditions interest from the religious standpoint only so far as they are assumed to have influence on the solution of the riddles of the inner life. The religious need is a kind of higher need of self-preservation, which is awakened by the disharmonies in our inner life, especially by the opposition between dependence and freedom, or between passivity and activity. The stream of our life, which springs from the innermost in us, is broken and split like a powerless wave against the skerries and rocks of external things. This is the ground of all pain, and then especially of the highest pain of all, that which is conditioned by the disharmony between our ideal and what we can carry through in reality. A harmony between these oppositions is attained, according to Sabatier, in prayer, the proper religious act. Here there appear at once the feeling of dependence toward the mysterious power we feel ourselves related to, and the feeling of freedom by virtue of the force that streams to us from this power. The essential in prayer is not the words in which it is clothed, or the ideas to which it is attached, but the inner connection of self-surrender and self-activity, which here mutually support each other and lead over into each other. The necessity of religion rests on its springing from our innermost conditions of life. Therefore it will show itself even where there is antipathy against
% --- p. 76 ---
\setcounter{footnote}{0}%
\opage{76}every external religion. There is an inner altar that can be preserved and kept in honor even if one overthrows all external altars. There is only one true atheism: frivolity and blasé indifference.

When the need is strong and inward, the satisfaction is already given; it does not come from outside. For Sabatier the whole mystery of religion lies in Pascal's words: ``You would not seek me if you had not found me''. The answer, revelation, is present in the very possibility of prayer, in the inner force and striving from which prayer springs. At the same moment man seeks God, it is God who seeks man. Revelation is not a communication once for all of unchangeable doctrines, which it then only matters to hold fast, but consists in the manner in which the consciousness of God arises and develops in the single individual and in humanity.

Sabatier designates this concept of revelation as the psychological or the critical, in opposition partly to the mythological, partly to the dogmatic concept of revelation. It has its great historical representatives in the Israelite prophets and in Jesus, but it finds application in every religious personality. Religious inspiration is as mysterious as the poetic, but no more miraculous than this. It is an inner fertilizing force, which works in harmony with all mental forces and laws. Prophetic inspiration is only a higher degree of what stirs in every strongly moved religious state of feeling. The standard for the value of revelation lies in the enrichment, liberation and elevation of the personal life that it effects.

The ideas and actions in which religion finds expression vary with the times. It is the great fault of orthodoxy, even of the Reformers, that they have not distinguished between the religious experience and its expression. The psychological phenomena and the theological explanations have been confounded. In every dogma a distinction must be made between the religious experience that lies at its ground and the more or less imperfect expression that reflection has been able to give it. This last element is the most variable element in religion. There may be three kinds of dogmatic variation. Dogmas may decline and
% --- p. 77 ---
\setcounter{footnote}{0}%
\opage{77}disappear for want of use (par désuétude), in that religious experience no longer affords the content they are to express. Thus it goes with the assumption of demonic possession and of a personal devil. Dogmas may next acquire another content than the original, in that they are interpreted according to new experiences. Such a displacement goes on uninterruptedly. By inspiration, incarnation and Trinity later generations do not understand the same as earlier generations. New wine is constantly poured into the old wineskins. Finally new dogmas may arise as expressions of new experiences. Thus there arose within Protestantism justification by faith and the universal priesthood.

To proclaim the immutability of the dogmas of the Church would be to proclaim the immutability of the philosophy that was employed to shape them. Dogmatics must keep to the religious experiences of which it is the exposition. And when certain concepts are employed in order to be able to make this exposition, it must also acknowledge the results to which the theory of knowledge comes regarding the significance and validity of our concepts.


Religious knowledge is symbolic knowledge. To express the feelings of trust, fear and love it uses images that are drawn from human relations. The power of the symbol rests on the close connection between feeling and image. Religious feeling does not at first distinguish between content and form; it is not conscious that it expresses itself symbolically. It is critical philosophy that impresses the unavoidability of symbolism, which both orthodoxy and rationalism each in its own way misjudge. --- The symbols may become common to large circles, and there are symbols that perhaps will never disappear from the religious life.

Not only the theory of knowledge does Sabatier require to be respected by dogmatics, but also the general method of the empirical sciences. He distinguishes, as we have heard, between the riddle of the world and the riddle of life. The riddle of the world it is the business of science to discuss. The world must be taken as science shows us that it is, and the riddle of life is then posed by the question of how science's
% --- p. 78 ---
\setcounter{footnote}{0}%
\opage{78}image of the world relates to the innermost need of the spiritual life
--- whether it contains means or hindrances for the satisfaction
of this need.

The independence of empirical science must benefit psychology
too, and here we stand (as I have already pointed out in my
Philosophy of Religion §. 66) at the most difficult point in Sabatier's
philosophy of religion. He has traced all phenomena of the
religious life back to religious experience, as he describes this
more closely, and he has maintained that all religious ideas are
symbols that express this experience. Yet he nevertheless assumes
that God is the cause of religious experience, and here
there would then be a case in which a phenomenon of experience
received a theological explanation, the possibility of which
Sabatier himself has expressly rejected, since God is not himself
a link in the series of phenomena and therefore cannot be used
for scientific explanation. Moreover, the very idea of God, like
all other religious ideas, must surely have a symbolic
character.

Sabatier's theological point of departure shows itself, however,
not only in his attempt to explain religious experience but
also in his very description of it. He brings forward only a
single type of such experience and does not lay stress on the
variations, although in principle he acknowledges their possibility
and their legitimacy. As with Siebeck, we feel here too the lack of
a broader psychological material. Such a material would perhaps have
led to a less narrow determination of the nature of religion.
Yet these objections do not prevent one from parting with the
French theologian's philosophical work with the memory of a
noble, refined and enthusiastic man.

\subsection*{3. RUDOLPH EUCKEN.}

After a series of preliminary works (especially \textit{Die Einheit des} \textit{Geisteslebens}
1888, and \textit{Der Kampf} um \textit{einen} \textit{geistigen} \textit{Lebensinhalt}
1896) this brilliant thinker, who works in Jena, the old
hearth of idealism, has expressed himself on the religious problem in his
book \textit{Der} Wahrheitsgehalt der Religion (1901).

% --- p. 79 ---
\setcounter{footnote}{0}%
\opage{79}It is not a system of philosophy of religion that Eucken has here
wished to give. Such a system is, in his view, for the present impossible
on account of the unclear and uncertain conditions in the religious
domain; ``In the spiritual anarchy of our age there is
no fixed and recognized point of attachment; every more deeply
penetrating discussion must go back to the first foundation and
build up anew upon it. Thus I too had, through a
general investigation of human existence, to work my way step
by step to the place where the problem of religion
breaks forth, so as then immediately to appear as the center
of all the search for the soul and significance of our existence''. The author
declares himself a seeker, and he addresses himself
to seekers. He finds in the age a strong longing for religion
combined with a distinct consciousness of the inadequacy
of religion's present form. The new type of life that the
Renaissance asserted, and the new labor of thought thereby
initiated, led away from Christianity. But the grandiose
culture in which this type of life and this labor of thought took shape
has brought with it inner entanglements, and faith in its sufficiency
has been shaken. Thereby the thought emerges whether
there might not be something in religion that cannot and must not
be lost. The lack is felt of a great goal, which not only
could cast light over the great circle of human life but also
raise the single individual above his often petty circumstances. And at the same time
the old riddles of life, which one regarded as solved or
at least as displaced, press forward anew. Above all a sharp
contradiction is felt between man's spiritual predisposition and his actual
condition. All this leads to a struggle between religion and
culture.

But neither religion alone nor culture alone commands
the means of deciding the struggle. The decision must come about by
distinguishing, both within religion and within culture,
between the eternal and the transitory, and between the indispensable
and the intolerable. A revision must be undertaken
of the very fundamental conditions of our spiritual life, and to undertake
such a revision is the business of philosophy. Not in the sense that philosophy
can construct a religion. For religion is something factual, an
% --- p. 80 ---
\setcounter{footnote}{0}%
\opage{80}expression of a real existence, which can be pointed out when it is there
but not artificially produced. But since what is at issue here is the pointing out
of a spiritual fact, the pointing out is not a
quite simple matter; spiritual labor is needed to find
religion's peculiar place in the whole connection of life.
Philosophy cannot get further here than to point out possibilities,
which only through great personalities and often under
great upheavals in human conditions can be converted
into actuality. But if it can be demonstrated that religion
is bound up with the innermost ground of our existence, then
every attack upon it and every criticism of it can serve only
to make its innermost, eternal core more and more manifest.

The problem is posed by the position of the spiritual life in existence.
Its persistence and its development presuppose that behind
and beyond the spiritual processes that experience shows us there lies
an eternal and coherent spiritual world (eine geistige
Substanz, ein seelischer Bestand), which contains the \textit{possibilities}
for new formations and the \textit{yield} of what has already developed.
The life of the mind must be more than mere phenomenon, hence
more than psychology can point out. It must have an eternal
principle in itself. Not merely a psychological but a metapsychical
view must be adopted. For such a view essence
and value are one, in that its highest object is the valuable core
of existence itself. Only through such a view does
a central content of life become possible.

The ancient culture overvalued form; the culture founded by the Renaissance
overvalued force. The true
culture, which has a religious character, sets \textit{formation of essence} (die
Wesensbildung) in contrast both to formation of form (die Formbildung)
and to formation of force (die Kraftbildung). Both form
and force are to serve a valuable content. Formation of form
alone leads to a congealing in plastic works of art, in which the
eternal is supposed to be expressed for all times. Formation of force
leads to restless movement, so that all being is dissolved. Formation of essence,
on the other hand, maintains a constant interaction of eternity
and time, in that the eternal possibilities are
% --- p. 81 ---
\setcounter{footnote}{0}%
\opage{81}converted into actuality through labor in time, and in that, on the other hand, the yield
of this labor is preserved in the form of eternity. Preservation of the
given and production of new content must go hand in hand,
in order that we may press forward to true reality.

The line of thought that maintains the possibility and necessity of ``formation of essence''
Eucken calls the \textit{noological} (from noos, reason,
mind, which among the Neoplatonists is placed above psyche, soul), and
he sets it in contrast to the psychological view,
which merely concerns itself with the conscious life given in experience.
It is, in his conviction, nothing new that he
thereby proclaims, but an old truth, which KANT in particular
has impressed by his definite distinction between the logical,
the ethical and the esthetic mode of consideration on the one hand,
the empirical-psychological on the other. Here validity
or value stands in contrast to reality --- but in such a way
that validity is in the end maintained as the highest reality.
--- The noological line of thought differs from the speculative or metaphysical line of thought of an earlier
age, which believed it could attain an understanding of the phenomena
by theoretical abstractions.
The noological method builds on experience, on
fact. To explain noologically means to place a spiritual
manifestation of life in the whole of the spiritual life, point out its position
and its task within this whole, and thereby illuminate and
consolidate it. This explanation differs from the psychological,
which consists merely in pointing out the processes through
which man acquires and appropriates a spiritual content.
Both methods must be applied, in order that the opposition between idealism
and realism may be overcome. Idealism has maintained the independence
and intrinsic value of the content of mind; the justification of realism
rests on its emphasis on the natural conditions. The
realistic mode of consideration comes into its own where it succeeds, within
the spiritual whole that the noological conception maintains,
in finding the special causes and laws; but it is
wrong in thinking that causal explanation can furnish the last
standard of truth. Idealism has rightly demanded a
valuable content for existence, but has disregarded the real
conditions.

% --- p. 82 ---
\setcounter{footnote}{0}%
\opage{82}According to Eucken, the noological view now leads precisely
to the fact that religion maintains: the presence of an
absolute life of the mind, exalted above and yet at work in the phenomena of life
appearing in experience. There is no religion without
the living presence of a higher world in our empirical world
--- but in such a way that these two worlds stand in an irrational
relation to each other. Religion can exist without belief in God; genuine
Buddhism is an example of that. But without an outlook beyond
the world of experience religion is an empty word. The motive for
religion lies in the opposition between essence and form of existence:
in mere experience the life of the mind stands as an appendage to the
material, as a secondary matter; but this conflicts with the unconditional
demands that the life of the mind makes, and the exalted values that it maintains.

From the fundamental claim of religion, reality looks
different from before. A view of the world proceeding from religion
will direct its eye toward what abides (in contrast
to what changes), toward the free (in contrast to the merely
natural) and toward the rational (in contrast to the irrational).
It will take particular note of the points where a new
sphere begins, such as the transition from the inorganic to the
organic, from the life of nature to the mental life, from the mental life to the life of the mind,
and it will conceive this graded series as an unfolding of
the innermost in existence. But such a religious
conception of the world cannot be proved. It has the character of a conviction,
not of a science. It is the business of philosophy to determine
the relation and the boundary between the scientific and the religious
view of the world.

The motive for religion lies in a spiritual need. It is not
a question of man's empirical felicity and its persistence,
but what is decisive is the wholeness and unity of the life of the mind,
the foundation of all special and individual manifestations of life. It
is its persistence that is at stake, in that mere experience
does not confirm it. It then depends on whether there can be awakened
a strong need to hold fast a conviction on this
point. Empirically viewed there are grounds enough for doubt. Experience
by no means shows us the world of the mind as an independent world,
still less as the main thing in existence. And in the course of the development of 
% --- p. 83 ---
\setcounter{footnote}{0}%
\opage{83}culture personalities are sacrificed to society and to progress.
Within the life of the mind itself disharmonious tendencies stir;
part stands against part, and part against whole; spiritual
tendencies combat one another, and evil proves to be a
real power.

In face of the doubts that thus arise, Eucken stresses, first,
that the strong feeling of disharmony is itself a consequence
of the great demands made upon life. The bitterness of suffering and
the harshness of resistance bear witness to the depth of our
existence and to the presence of higher forces in our being.
With the oppositions pain too would fall away. It was
never in times of unhappy experiences and heavy entanglements
that mankind doubted its ideal task, but on the contrary this doubt
thrived in such times as sluggish comfort and
imagined satiety prevailed.

Second, it is Eucken's conviction that that doubt
must be combated by a more powerful and more concentrated form
of \textit{religion} than the indeterminate and universal \textit{religion} that is
described in what precedes. He points to the significance of the great personalities
in the history of religions. In them
the divine stirred in a deeper and more powerful manner than in
others, and they were able to express life's manifold experiences
in clear and glowing images. A new life broke through here
in a decisive manner and intervened decisively in the historical
development of the race. By bringing forward their significance we make
the transition from \textit{the universal} religion to \textit{the characteristic}
religion.

``Characteristic religion'' denotes in Eucken the historical
or positive religions, in so far as they --- each in its own peculiar
manner --- give a new total image of life, clear in its fundamental features,
and set this image in contrast to the world of experience.
It differs from ``universal religion''
by the determinate shaping without which such an opposition
to the world of experience cannot be maintained. It springs
from great personalities and rises above all popular culture as
an immediate utterance of the innermost essence of the life of the mind. Without
such characteristic religion independent inwardness is not
% --- p. 84 ---
\setcounter{footnote}{0}%
\opage{84}possible. It is here not a question of a mere embellishment of life,
but of the breakthrough of a new life, a new reality. Yet
the breakthroughs to which the great positive religions bear witness are only
high points and turning points within a great process. The divine,
which here breaks through, did not first come into being at
these particular places.

The higher life that thus breaks forth our concepts are not able
to express. Imagination must come to our aid
with its symbols. And here lies religion's vulnerable point.
Imagination is always determined by the given stage of culture and its
traditions, and the symbols will therefore acquire a historical character.
A time will one day come when that which perhaps for millennia
was an expression of the highest truth, and which for
human consciousness had entered into an intimate connection
with the religious content, is shaken in its significance and
must yield its place to new forms --- often only after severe tremors
and hard struggles, during which doubt, from being directed
against the symbols, may pass over to attack the innermost
core of religion.

The universal religion is constantly necessary as a corrective
of or introduction to the characteristic religions, which, left
to themselves, are easily narrowed and become mere means of consolation.
The significance of the characteristic religions rests not on the fact that
any of them should have found a conclusive expression of the
highest truth, but on their vivifying and elevating
effect amid the unrest and vicissitudes of the time. It must ever anew
be tested whether what is historically handed down really performs such a
work. But the Church has a tendency to place itself slavishly
before a model given once for all and to conceive all
truth as consisting in imitation and repetition. It isolates
the highest and thereby gives it a passive position in relation to
real life. Thereby struggle always becomes necessary
in the religious domain. It is the significance of the great religious personalities
that they compel us to choose. ---

What especially determines Eucken's place in modern
philosophy of religion is his decided conviction that the values
can be maintained only when, behind and beyond the world of 
% --- p. 85 ---
\setcounter{footnote}{0}%
\opage{85}experience, a higher reality is assumed, in which the life of the mind
that in experience appears scattered and in changing forms
persists in the form of unity and eternity. In an often profound
and atmospheric manner he maintains the necessity of this opposition,
--- now so that he recalls the old mysticism, now so that he recalls
romantic idealism, and now so that it is the exalted poetry of the personal
struggle of life that finds expression in his words. Yet he stresses
very strongly that philosophy can here only point out possibilities, and
that it is through prophetic personalities that such possibilities
gain actuality and definite form. When he --- in order to
avoid subjectivism --- maintains the necessity of a metaphysics
(or a ``metapsychics''), it is yet properly by virtue of
the \textit{value} of the true, the beautiful and the good that this metaphysics is to
be necessary. What he calls the ``noological'' method
differs from the psychological method in that it will
show that that value can be maintained only by the assumption of a
connection of the whole that bears the valuable manifestations of life
appearing in our experience scattered and fleetingly. In particular the
noological mode of consideration lays stress on the great breakthroughs,
the significant beginnings and the qualitative differences
within the life of the mind as we know it in experience. He regards
the attempts of empirical science to show continuity in
the world of phenomena as hopeless, but interprets the breaks
in empirical continuity as single outcroppings of a great
transcendent continuity, a realm of possibilities, ends and
values, which only by glimpses reveals itself in the world of
finitude. In this he is akin to romantic idealism. But noology differs from the speculative
method in that the assumption of such a connection of the whole
is expressly grounded only on experience of the
valuable. It seems then clear that the point of departure is empirical
and personal. He gets his impetus to pass over from experience to
``metapsychics'' through the tension that is produced by a
great dilemma. Either --- he maintains --- one must regard
the life of the mind as a fleeting phenomenon, an appendage to the material
nature, --- or one must assume that it in itself 
% --- p. 86 ---
\setcounter{footnote}{0}%
\opage{86}constitutes a great totality and is one with the innermost core of existence.
This dilemma comes forth in him both where he develops
the idea of ``the universal religion'', thus would make the transition
from the empirical standpoint to a religious standpoint of a more
indeterminate kind, and where he goes over to ``the characteristic
religion'', thus would show the necessity of a historically formed
religion. Characteristic are especially the following utterances: ``In
this decision it is not this or that but the life of the mind
as a whole that is at stake. In man's experience
the life of the mind appears in an untenable intermediate position; from this man must
either turn back to mere nature and confess that all
peculiar striving was only a heavy dream, or boldly go
forward and secure for the new striving a new world. Only an unclear
way of thinking can seek a middle way between these ways, and only
a feeble disposition can calmly endure the dilemma. In any case
it stands fast that without religion there is no truth in the
life of the mind and no inner greatness in man.'' (Der Wahrheitsgehalt der Religion. p. 238.) ``Is human life a
mere appendage to nature, or is it the beginning of a new world?
On the answer to this question depends all shaping
of our relations, the whole direction of our action. Religion
would raise human existence up to a height raised above the world and
thereby save our life from the nothingness into which
it otherwise surely falls. If this its attempt is only
a flight of Icarus, then all hope falls away, then the noblest and
best stands as empty imagination, then the whole ends in unreason.''
(ibid. p. 319). ---

What is to drive us to choose here is nevertheless always in the end
personal experience of life, in particular the degree to which
one is gripped by the fate of values in reality. There seems
to be no ground for introducing a new term;
``noology'' is neither a necessary nor a happy enrichment
of the already sufficiently supplied philosophical
terminology. From a private utterance of Eucken's, which I
hope to be able to cite without committing an indiscretion, it is seen
clearly that it is the personal experience of life's oppositions,
in particular of the great opposition between ideal and reality
% --- p. 87 ---
\setcounter{footnote}{0}%
\opage{87}that has determined his own choice and led him to ``noology''.
In a letter he wrote to me after reading my
Philosophy of Religion, he states, after having dwelt on the predominantly
psychological character of my attempt to treat the
problem: ``Mich treibt es, bei allem Gegensatz zur alten ontologischen
Metaphysik, immer wieder zur Metaphysik und so
auch zu einer metaphysischen Begründung der Religion; ich
fürchte ohne das dem Subjectivismus zu verfallen, und dagegen
sträube ich mich namentlich deswegen, weil ich stärker unter
dem Eindruck der grossen, in der Erfahrung unlösbaren Konflikte
unseres Daseins stehe, kurz gesagt, dualistischer denke
und denken muss.'' In this utterance Eucken points directly
to individual experience as determined by the peculiarity of the single
personality. It is this that determines his
standpoint. But the more clearly this foundation appears, the more natural
it becomes to assign to the philosophy of religion the task
of investigating the psychological soil from which the different
religious standpoints grow up, and to attach to this an examination in the
theory of knowledge and an ethical examination of the thoughts and endeavors
in which these standpoints find outlet. A quite
peculiar method the philosophy of religion then does not need, and the so-called
``noology'' is indeed itself a ``middle way'' between speculation
and practical faith. As speculation it falls to
the treatment of the theory of knowledge, as practical faith to that of psychology and ethics.
The conception of the task and procedure of the philosophy of religion
that follows from this seems to be most in the
spirit of critical philosophy, in which Eucken --- like
Siebeck and Sabatier --- would work. The ``metaphysical'' grounding
of religion that Eucken demands of philosophy proves,
according to his own utterances, in the end to rest on the need for a
decision of a great practical dilemma.

I find in Eucken's standpoint --- in the ``noology'' that
leads him to his religious metaphysics --- an example of how
personal experience of life can lead to faith in the persistence of value.
For me every faith, even though it appears in the
most philosophical forms, is always only \textit{object} of philosophy, not
itself \textit{product} of philosophy. In any case the philosophy that
% --- p. 88 ---
\setcounter{footnote}{0}%
\opage{88}produces it is itself a product of personal experience of life. Philosophy
as science has, after all, no other methods
than the logically inferring, the psychologically explaining and the
ethically evaluating.

But from my standpoint such works as Eucken's,
and the same holds of Sabatier's, acquire an exceedingly great interest
through the clear manner in which the decisive motives
for a religious view of life come forth. Two men of ideal mind, gripped by life
and equipped with energy of thought, have here expressed
themselves on one of the greatest human problems.

\section*{II.\\ INVESTIGATIONS IN THE PSYCHOLOGY OF RELIGION.}
\streg

A thoroughgoing attempt to go back to immediate experience
in treating the religious problem must in particular
direct attention to the point where consciousness
passes from immediate experience to forming a judgment, thus
to the first setting up or involuntary coming into being of a religious
proposition. What goes on here cannot be essentially different
from what always goes on in the transition
from observation to judging. Religious consciousness has neither
a psychology nor a logic that differs from the ordinary
psychology and logic, and religious judgments must come into
being as all other judgments come into being. What is peculiar
to religious judgments is that the experience on which they ultimately
rest is an experience of feeling, an experience of the value of certain experiences
as furthering or hindering, and that the experiences
that are here in question concern the relation between ideal and reality,
between the world of values and that of reality. Religious 
% --- p. 89 ---
\setcounter{footnote}{0}%
\opage{89}values are therefore values in the second power or secondary values:
they presuppose that man knows something that gives life
value (in whatever this then consists), and that he then experiences that there
are powers in the world that further or hinder this valuable thing.
At the high points of religious feeling the difference
between primary and secondary value naturally falls away; there is known
no Why and no predicates are applied. In most pronounced
form this occurs in mysticism. But where a judging can take place,
thus where there can be talk of statements of religious consciousness,
there --- I maintain --- a closer analysis will show that
what something is said about is the relation between value and
reality, as the individual believes he has immediately
experienced this. Religious judgments concern the fate of values in
the world of reality. --- Here it must, however, well be remarked
that religious judgments are not always, perhaps even very seldom, formed
at first hand.

In every formation of judgment one must (as I have more closely
developed in my treatise Det psykologiske Grundlag \textit{for logiske}
\textit{Domme}) distinguish between the initial idea and the concluding idea.
In the linguistic formulation of judgments these two ideas appear
as logical subject and logical predicate.
At the beginning of the process of judgment the initial idea has
a certain indeterminacy, which is sought to be removed by the discovery of a
predicate, a closer determination, which takes place through analysis
of the initial idea and of the connection in which it
occurs. The richer the content of the initial idea is, the more
difficult it will be to find a definitive concluding idea,
a concluding and adequate predicate. In this respect
religious consciousness, as it appears
in history, presents the greatest differences. The narrower both
the world of values and that of reality is, and the more elementary
the relation between them appears, the easier it is for
man to rest content with a predicate and thereby conclude
religious thinking. Here are all possible degrees from
the limited and childish predicates of fetishism and animism
to the sublime limiting concepts of the highest mysticism. And the
process by which the concluding idea is here linked to the initial
% --- p. 90 ---
\setcounter{footnote}{0}%
\opage{90}idea, is now involuntary association of ideas, now imitation and habit, now speculative brooding, now
anxious search and
questioning. The content of the initial idea in religious experience is
the harmonious or disharmonious relation between value and
reality, and the concluding idea is to state what it is that comes to light in this relation. If such a concluding idea is won or found, the transition has been made from experience to faith. The concept of God in its various forms is the basic predicate in which religious consciousness seeks to gain expression for and explanation of its deepest and most essential experiences. The word God is to mean ``him to whom sacrifice is made'', thus him to whom man feels himself in debt, to whom he refers his experiences of the relation between value and reality. But who is that? Qvid est quod amo, quum te amo? asks Augustine. Here lies the possibility of ever new brooding, especially since the initial idea is not constant, inasmuch as the relation between value and reality is constantly changing. There is in Goethe's Faust a monologue that suits well as an example of what is here meant by religious experience. In the scene ``Wald und Höhle'' Faust makes a reckoning. Nature's and thought's great treasures are opened to him; he can immerse himself in the inner as in the outer world in understanding and mood; but at the same time there has been given him a demonic companion whom he cannot do without, and whose coldness and impudence transform those
high gifts into nothing --- at the same time as he awakens the flaming fire of sensual desire. A concluding idea does not come about here; if one wishes to seek one in the poem, it must be found partly in the Prologue to the first part, partly in the end of the second
part. --- The great opposition that Goethe's Faust here accounts for
lies --- only in still sharper form --- at the ground of Schopenhauer's
experience of life, and Schopenhauer's philosophy was in its innermost
ground an attempt to find a concluding idea as expression
for this experience. ---

Immediate religious experience is an ideal, since every human being we know begins under the influence of a tradition and at once interprets his experiences by the help of tradition, or perhaps more exactly, at once sees them in the forms and the 
% --- p. 91 ---
\setcounter{footnote}{0}%
\opage{91}illumination that tradition has arranged. Already the growing into a
language brings about that there can here be no talk of absolutely
original experience. Every religious experience is more or less
confessional. Both the values, reality and the relation
between them man apprehends from the first under the influence
of ideas that he has not himself formed on his own,
still less himself derived from his experiences. It will then be of great
interest to investigate how the traditional ideas
are gradually altered by the individual's own advancing experiences,
and next to make a comparison between religious
standpoints and processes of development that are determined by essentially
different points of departure.

There now exists on all the topics touched here --- both on
religious judgments as value judgments, on the forms and modifications of a confessional
religiosity and on comparative psychology of religion
--- a series of works, which for the most part have appeared
after the completion of my Philosophy of Religion, and which it
will be of interest to have characterized.

\subsection*{1. ON RELIGIOUS JUDGMENTS. \textit{(MAX REISCHLE).}}

The attempts at philosophy of religion with which we occupied ourselves in the first chapter were more or less consistently built
on critical philosophy and its distinction between purely
theoretical judgments and value judgments. Within theology this
distinction has in recent years played no small role, especially with
ALBRECHT RITSCHL and his disciples. Auguste Sabatier stands
close to this tendency, and he also speaks of value judgments
(jugements d'estimation ou de dignité). The weight that is laid
on the difference between the two kinds of judgments is characteristic
of the endeavors of the most recent time to lead religion back
to immediate experience and to determine the proper
character of this experience. It is through the investigation of value judgments
that in the most recent time there has come about a contact between theology
and philosophy, which otherwise in the last half-century --- especially
on account of theology's more and more pronounced confessional
% --- p. 92 ---
\setcounter{footnote}{0}%
\opage{92}character --- have stood far apart from each other. --- I shall here as
an example draw attention to a treatise by MAX
REISCHLE: \textit{Werturteile und Glaubensurteile}. (Halle 1900).

A value judgment presupposes, according to Reischle, a feeling of value
that differs from the mere feeling of pleasure. The condition for
the arising of a feeling of value is that the satisfaction concerns the whole
self (das Gesammt-Ich), not a single part or state of the
self. Value I attribute to a thing of which I have seen by way of reflection
that it will satisfy my self as a whole.
Yet Reischle concedes that the boundary between ``the whole self''
and its single parts or states is fluid. What he wishes to
stress is especially that before the \textit{value judgment} there goes a
composite process, by which the feeling of value attains a clarity
and certainty that the mere feeling of pleasure lacks.

Value judgments can, according to Reischle, be of two kinds: natural
and ideal (or normative).

\textit{Natural value judgments} are grounded in man's nature
and conditions of life, which determine the satisfaction that is possible
for him. Here many differences make themselves felt between
the different individuals, but since interests are nevertheless to a certain degree
solidary, there can be collective value judgments, and since
conditions of life are to a certain degree common, there can be general
value judgments. Where an external standard (e.g. money) is
possible, the value judgment can acquire an objective character, so that perhaps
even its subjective origin is forgotten.

\textit{Ideal or normative} value judgments presuppose an idea or
a norm that we recognize and expect to be recognized by others. Here
belong esthetic, ethical, intellectual and religious value judgments.
These are determined not merely by the manner in which we are passively
determined from outside, but essentially by agreement
with a law that holds for the activity of our own mind. In
ethical value judgments e.g. we follow a norm or an idea that
is able to lead our will to harmony and freedom. In religious
value judgments the underlying idea is the significance it
has to submit oneself to a highest power in order by unity with it
to reach the highest freedom in face of the world. The ideal value judgments
make claim to general validity, because no one can
% --- p. 93 ---
\setcounter{footnote}{0}%
\opage{93}deny the ideas that condition them without giving up his personality,
the freedom and inner connection of his being. ---

The difference between natural and ideal value judgments is ---
what Reischle has not noticed --- evidently just as fluid
as he himself concedes the difference between feeling of pleasure and
feeling of value to be. That ``self as a whole'' which the feeling of value is to
presuppose, namely contains a norm that already holds for
the special and changing elements and states; their value is
determined by the relation to it. And besides, on the other hand, ``the self as a whole''
must itself demonstrate its value in face of
the single elements and states whose subordination is demanded;
it is by no means self-evident that the whole is in the right against
the single element or the single state. Finally, according to Reischle himself
the value of the esthetic, ethical, intellectual and
religious functions of mind rests on their procuring our being
harmony and freedom; --- but that means nothing other than that
they make a ``self as a whole'' possible, and the ideal value judgments
become, seen from this side, not different from the natural. It
is therefore by no means easy to set up measures of value in the spiritual
domain, where there is the most intimate interaction and mutual
determination of the elements among themselves and between elements
and whole. Within the individual consciousness, as
within society, the different measures of value struggle.
The value of every idea or norm must itself be demonstrated --- and therefore
the problem is posed ever anew. Of all value judgments precisely,
as I have shown elsewhere, those that lead to the setting up
and recognition of a norm present the greatest difficulties
for theory. The transition from the ``natural'' to the ``normative''
(ideal) includes great problems, which a mere classification
of the actually given value judgments easily covers. Reischle
has a clear eye for the analogous problem that arises when the
different religions, each with its basic values, stand over against
one another; but he has not seen that the struggle can already
begin in the interior of the single individual, and that a guiding idea or norm
is always the result of such a more or less violent
struggle. ---

In the last part of his investigation Reischle discusses how far
% --- p. 94 ---
\setcounter{footnote}{0}%
\opage{94}\textit{articles of faith} can be called value judgments. He concedes
that authority and tradition play so great a role both in the
Protestant and in the Catholic Church that the articles of faith that the
single individuals profess cannot be called value judgments in the sense
that these single individuals themselves have derived them from their experiences of value.
On the other hand Reischle considers it undoubted that the articles of faith
must be called value judgments when it is asked on what
their claim to validity rests in principle. They are value judgments
of ideal or normative kind; their significance consists in
their stating the conditions for blessedness, i.e. for the personality's
complete harmony and freedom. It would be wrong
to call them postulates, as if it were merely our own
need and yearning that formulated in them what could produce
satisfaction. There lie at the ground of the religious judgments definite
facts, for the Christian thus the life and
death of Jesus, and the religious judgments express trust in the power
and significance of such facts. The Christian believes in Christ,
because only in adherence to him can he find the true
life. Of all historical religions it holds that certain given
phenomena --- persons and events --- are interpreted as divine
revelations, and the religious judgments express trust
in these revelations. Striking, especially beneficent or
destructive phenomena in nature, in the conscious or the
unconscious mental life, in the people or in history awaken the involuntary
idea that higher powers are here at work for the
weal or woe of men. Often the judgment is dictated by fear, most often however
by trust. And the higher a religion stands, the more it is ethical
ideas that dictate the judgment. The very fact that something is called ``revelation'',
is thus according to Reischle a value judgment; it depends
on whether the personalities or events in question stand as
forms of the entry of great values into history, and whether they
can become object of trustful adherence and devotion.
A purely objective decision of the strife between different revelations
is therefore impossible. This is a necessary consequence
of empiricism in the philosophy of religion, and Reischle also draws
this consequence, which he denies only by putting the word
% --- p. 95 ---
\setcounter{footnote}{0}%
\opage{95}``revelation'' in quotation marks where it is not a question of
his own religion.

\subsection*{2. THE PATHOLOGY OF RELIGIOUS FEELING. \textit{(MURISIER).}}

To be able to point out the elements in religious feeling
it may be expedient to study it in such forms
in which it assumes a morbid character. The morbid consists
in elements that otherwise occur in harmonious connection with
other elements separating out and appearing with overwhelming
strength, so that they dominate the whole course of life. In his
book \textit{Les maladies du sentiment} religieux (Paris 1901) E.
\textit{MURISIER} (Professor in Neufchâtel) has set himself the task of investigating
the morbid forms of religious feeling, which he thinks he can
trace back to two, ecstasy and fanaticism, the former the extreme
expression of religious feeling in its individual form,
the latter its extreme expression in social form. These two forms
stand in opposition to each other: the one leads to absolute isolation,
the other to outward-going activity and to connection
with others. But there lies, as Murisier shows, one and the same
need at the ground of both types. The mystic sees in isolation
and in the removal of all multiplicity from consciousness the
only means to preserve unity and security, to hold fast
life's total direction toward a single goal. His mind
lacks energy to master the multiplicity of life in him and
outside him, and he therefore collects himself on a single point.
The fanatic seeks, on the contrary, to attain unity and security in life's
course and direction by a practical acting that leads as
many as possible to share his faith. He seeks to bring about
a social unity within which no disharmonious
and disturbing tendencies are felt. By social uniformity
the unity and identity of one's own self is supported. The fanatic praises
ecclesiastical unity with just as great zeal as the mystic the unity
of thought. And just as eagerly as the mystic keeps away all impressions
that could disturb his absorption in the One, which is the
highest for him, just as eagerly the fanatic suppresses every
% --- p. 96 ---
\setcounter{footnote}{0}%
\opage{96}tendency that could break the social unity. His persecuting
corresponds to the mystic's asceticism, and his proselytism to the mystic's
concentration. For many people adherence
to a religious community, especially a narrow and exclusive one, is the only
way along which they can find equilibrium and peace of mind.

Through the study of these opposite forms of religiosity,
supported on a historical, biographical and medical material, Murisier
finds in the character marks that are common to them a contribution
to the characterization of religion. Religion is due --- whether
it appears in individual or in social form --- to a
need for unity that is closely connected with the need for self-preservation.
This need for unity appears in religion particularly as
a need for guidance (besoin de direction), and in so far Schleiermacher's
description of religion as feeling of dependence was a
happy stroke. In morbid form the direction toward unity, the concentrated
dependence can be attained only through exclusion of heterogeneous
impressions and directions; in healthy form it is attained
in a more harmonious manner through ``edification'' and organization.


Religion in its social form goes before individual
religiosity; historians of religion therefore concern themselves mostly
with the former and are inclined to regard the latter as a sign of dissolution.
The religion of primitive peoples consists mostly of ideas of faith
and ceremonies that are common to the whole society. The
self of primitive man has not separated itself out from family, clan
and people, just as the child's consciousness still draws its first
content from the nearest surroundings. The very ecstasy of the
savages' sorcerers reminds more of the fanatics' state of mind
than of the mystics' ecstasy; a certain development is required, before
the properly individual form can come into being. But the sociology of religion
can at no point be quite separated from the psychology of religion.
The very need to produce social unity or join
one has a psychological cause, which must be
studied in the single individuals. Even where religion most nearly
fulfills a social demand, it is yet essentially borne by an inner
yearning in the minds of the single individuals. ---

Much can certainly be objected to Murisier's 
% --- p. 97 ---
\setcounter{footnote}{0}%
\opage{97}conception of mysticism as a purely individualistic phenomenon. Partly
the mystic begins as a rule within a confessional religion,
and partly his own self belongs to that from which he wishes to be freed;
he wishes to get out of the narrow conditions, beyond all limits, and mysticism
is often an altruism driven to the extreme. Mysticism
is a phenomenon of expansion no less than a phenomenon of concentration,
or rather, there are manifold nuances within
the mystical direction. --- There are also more forms of
morbid religious feeling than the two main forms to which Murisier
keeps. There are thus states of anxiety and despair,
where neither expansion nor concentration is possible, and
there is in contrast to this the exalted self-feeling, where honor
and power of being linked with the Highest are enjoyed in an egoistic
manner. This opposition between potentiated feeling of impotence
and potentiated feeling of power is perhaps no less significant
for the pathology of religion than the forms described by Murisier.

What does not come forward sufficiently in Murisier is
the connection of religious feeling with experiences of value. One
could have a guiding thread for finding the morbid forms
of religion by proceeding from the presuppositions to which the study
of religion will again and again lead back: experiences of
values, of realities and of the relation between values
and reality. But the most scientific standpoint would
nevertheless be that which regards the pathological forms of religion as
special cases of the ordinary forms of mental disease. Taking up
this standpoint would also lead to seeing the limitation
within which one must move when one
from the phenomena of the sick conscious life would draw conclusions
concerning the nature of the healthy conscious life. --- These objections
nevertheless do not prevent one from acknowledging that in Murisier's little
book there is found much good and fine psychology, and that it --- within
the limitation the author has chosen --- illuminates the religious
life from a characteristic side.
% --- p. 98 ---
\setcounter{footnote}{0}%
\opage{98}

\subsection*{3. CONTRIBUTIONS TO THE PSYCHOLOGY OF CONFESSIONAL\\ RELIGIOSITY.}

From the religious side it is maintained more and more that the foundation
of religion is experience, and every attempt at a psychological investigation
of religion leads to the same result. But then it becomes of decisive
importance to investigate the kind and character of the experience that
can be in question here. Otherwise one cannot know how much this foundation
is able to bear, and dogmatism once again makes its entry. But here a
difficulty now arises, which comes from the fact that every state of feeling
has a tendency quickly to enter into connection with definite ideas and to
lead to definite actions. And the ideas and actions that always lie nearest
at hand are those prevailing in the surrounding society. The individual
lives himself into them, perhaps even before he has begun to make his own
life experiences. He inherits the valuations of the race before he can begin
to value life on his own. And when he begins to make his own experiences,
he will in any case at once settle them into the received forms. To a large
extent, if not for the most part, his experiences are called forth precisely
by the impression that received ideas and actions make upon him, especially
through the captivation of those around him by them. Only with great
difficulty and under special conditions will that counter-test become possible
which is needed for a real proof from experience, --- namely the test of
whether the experience and the valuation of life to which the individual
himself, according to his nature and his circumstances, can attain, can now
really be expressed and explained solely and exclusively through the received
forms, so that content and form wholly correspond to each other. Possibly
the same content could be expressed through other forms, or other ideas and
actions could even make a richer and more valuable life experience possible.

Religious experience therefore becomes, historically viewed, always more or
less confessional. It becomes predominantly a repetition of what others have
previously experienced and lived through; it moves in ready-made expressions
and formulas, and that which at the very
% --- p. 99 ---
\setcounter{footnote}{0}%
\opage{99}outset, namely when the expression was found and when the formula was
created, was effect and result, operates in confessional experience as cause,
often the chief cause. Yet within confessional experience there are all
possible degrees between the passive living into given forms and the
independent appropriation, which at its height can go on to find new paths
and create new forms.

The psychology of religion must keep this fact in view when it gathers its
material of experience. A rich material we have here --- besides what one's own
knowledge of human nature yields --- in the utterances of religious personalities,
especially in autobiographies. It is above all on such material that I have
relied in the psychological section of my Philosophy of Religion, in that I
at the same time compared the types I found in this way with the types
presented by the great popular religions, especially their founders. But it
will naturally be of importance to obtain a broader basis, and this could be
attained if information and utterances could be collected from larger circles,
also from non-literary circles. Attempts to procure such meet with many
difficulties, partly because of the helplessness that will prevent many from
expressing themselves clearly, partly because of the shyness that will prevent
many from speaking about what is most inward in their lives; the mere question
will perhaps even at times be regarded as a profanation. But since psychological
research, like all other research, is pursued for the sake of truth and clarity,
it will be fully justified to go this way and try how far it can lead. It is
especially in North America that this has been attempted. In North America the
religious life stirs from the practical side, and the prevailing Methodism with
the revival movement attached to it has naturally directed attention to the
mental processes that go on during religious experiences. It is also the psychology
of conversion that the American psychologists have especially occupied themselves
with. The leaders of the revival movement had already begun a statistics of the
frequency of conversions, and the first attempts from the purely psychological side
were in part made in cooperation with them. By means of circulated question lists
an attempt has then been made to obtain
% --- p. 100 ---
\setcounter{footnote}{0}%
\opage{100}information about the states before, during and after the religious breakthrough,
as well as about other types of religious development that did not have the character
of breakthrough, about crises and doubt, etc. Through analysis of the utterances
obtained an attempt has then been made to arrive at a description of certain main forms
of religious development --- naturally within the given historical presuppositions
(namely American Methodism). Very large numbers could not be in question here, and
therefore not a strict statistical method. One has cast the net out at a single place
in the sea of religious phenomena and described the species one found in one's haul.
There has, as the following will show, really been gained new material for comparative
psychology in this way.

The first in the series of these investigators is JAMES LEUBA with his treatise
\textit{Psychology of religious} \textit{phenomena} (American Journal of Psychology.
April 1896). From his childhood Leuba had moved in circles where religious breakthroughs
were frequent, and later he has, from psychological interest, extended his experience
by conversations with people who had experienced such, by attending meetings and by
sending out question lists concerning the circumstances before, during and after
conversion. ``The question lists,'' he says, ``were sent to persons who thought they
had `experienced religion,' to ministers and to leaders of the Inner Mission, and
through them also to a large class that it had not been possible to reach directly.
They were moreover announced in `The Presbyterian' in Philadelphia, in `The Christian'
in London and in `The Outlook'\dots{} A small percentage of the answers received
contained the phrases that are current at revivals and in mission houses and smacked
so decidedly of religious affectation (religious cant) and so little of personal
experience that they had to be set aside; but most of the answers bore inner testimony
of serious endeavor to give an exact description of an experience to which immense
importance was attached\dots{} The yield of our questions was greatly limited by the
lack of confidence of certain people in the usefulness of applying psychology to religious
phenomena, by the difficulties which ordinary people, who are not practiced in introspection,
feel in having to give an account of their own experience, by the hostile attitude of some
ecclesiastical
% --- p. 101 ---
\setcounter{footnote}{0}%
\opage{101}authorities, and finally by the aversion that the invitation to investigate the nature
of religious experiences sometimes aroused.'' Regarding this last point Leuba writes in a
later treatise as follows: ``The ground tone of deep, sometimes even tragic seriousness that
runs through most of my documents would move even an ironic skeptic to awe. The real value
that is attributed to religion reveals itself when the intellectual assumptions on which the
religious life is thought to depend are threatened\dots{} A lady of my acquaintance, to whom
I gave a question list, recoiled in great horror when she had read the first questions. Her
unclear grounds for her refusal to answer showed that she did not dare to look her faith
straight in the face; the mere thought of it struck terror in her. Another lady, who with
striking vehemence had declined a personal request for an answer, wrote a year later of her
own accord to defend her conduct. Her faith was, at the time when the question came, in a
precarious state, and she fought with the energy of despair that her heaven should not lose
its God. Her chief means of defense was to drive from her consciousness anything whatever
that could threaten her faith. The critical investigation I proposed wounded her to the quick.
These moving examples of a class of our fellow men who prefer to stick their heads in a bush
rather than let the hard facts rob their world of the soul that gives it grace, awaken to
significant reflection on the connection between faith and need.''

In a series of later treatises Leuba has continued his investigations (\textit{The}
\textit{personifying} passion in youth. The Monist. July 1900. --- \textit{Introduction}
\textit{to a} \textit{psychological study of religion}. ibid. January 1901. --- \textit{The
content of} religious consciousness. ibid. July 1901. --- Les tendances \textit{religieuses chez}
\textit{les mystiques} \textit{chrétiens}. Revue philosophique. Juillet—Novembre 1902). What
characterizes Leuba is his interest in the individual cases in their multiplicity and
peculiarity. He has brought together a collection of religious self-characterizations, which
show how many nuances and transitions there are in this field, and how little our usual
rubrics suffice. ---

% --- p. 102 ---
\setcounter{footnote}{0}%
\opage{102}Related in design and in significance to Leuba's investigations is a recently
published treatise by LUCIEN ARRÉAT: \textit{Religion} in \textit{France}. (\textit{The}
Monist. January 1903). In his investigation of French Catholicism he gives an interesting
counterpart to Leuba's investigation of American Methodism. ---

EDWIN STARBUCK (The \textit{psychology of} religion. An empirical study \textit{of} the
growth of religious consciousness. London 1892 disposes of a very extensive statistical
material, which he has subjected to a thoroughgoing analysis. He conceives conversion as a
process through which the young person is to make what is socially handed down into full
seriousness and reality in himself, and therefore as a sociological phenomenon that can be
demonstrated among all peoples and in all mental domains. Of particular interest is his
demonstration of different kinds of conversion and his comparison of the process of conversion
with other forms of religious development.

GEORGE COE (\textit{The spiritual life.} Studies in the science \textit{of}
\textit{religion}. New-York 1900) has by a similar path as Leuba and Starbuck especially
investigated the influence of temperament on the character of the religious development and its
relation to other co-operating elements.

GEORGE SANTAYANA has in his work \textit{Interpretations of}
\textit{Poetry} and Religion (London 1900) especially brought forward the poetic
element in religion and has thereby been led to place himself in opposition both to orthodoxy
and to the current criticism. For him religion is the highest poetry, while he at the same time
seeks to show how poetry at its height passes over into religion. ---

On the whole the standpoint of the American psychologists toward the religious life is a
sympathetic one. They energetically assert the independence of psychology in relation to all
theology and speculation, but they immerse themselves with enthusiasm and interest in the
religious phenomena. Their own religious standpoint does not seem to be markedly anticonfessional,
even though they think they can demonstrate the necessity that the confessional limitation give way
to the valuable multiplicity of paths and types that unbiased observation discovers. As far as possible
I shall in the following
% --- p. 103 ---
\setcounter{footnote}{0}%
\opage{103}discussion of their investigations also seek to determine to which type they themselves belong,
in order thereby to determine the personal equation and indicate a limitation for the significance of their
investigations.

\subsubsection*{\textit{a. Breakthrough and Continuity in Religious Development.}}

STARBUCK found in the material he collected two types of religious development. The one moved
through a more or less sudden decision on to a new direction of life; the other proceeded more
continuously.

\textit{The breakthrough type} again presents differences of interest

The age at which the breakthrough takes place can vary. It is most frequent in the young years and
occurs somewhat earlier for women than for men; later than at the age of 40 it occurred rarely (within
the 1265 cases where the age could be ascertained). --- The occasion for the breakthrough varied
A large part of the cases was due to revival meetings (revivals). Some took place during ordinary
church service. In others again no particular external occasion could be demonstrated. --- As regards the
motives, fear occurred far more frequently than hope, and altruistic motives (ideality and love of
humanity) were rare. It is characteristic that the feeling of sin was rarer at revivals than on other
occasions; Starbuck thinks that at revivals one does not get time to immerse oneself in oneself and brood
over one's conduct and one's past.

Most important are the differences with regard to the character of the breakthrough itself. Two forms
appear here, which, however, can be demonstrated in their purity only approximately. In the one form it is a
habit that is broken, a vicious state from which one escapes; a deep feeling of sin then precedes. In the
other form it is a new light that breaks in; a lack is filled and a positive fulfillment is reached; a feeling
of incompleteness then precedes. The former form has a negative, the latter a positive character. The negative
form requires complete overcoming and repression of the past in order to come about; the
% --- p. 104 ---
\setcounter{footnote}{0}%
\opage{104}positive form requires only a small impulse to take effect. Starbuck remarks that within the limits of his
material the positive type (spiritual illumination, preceded by a sense of incompleteness) is more frequent than
the negative (escape from sin, preceded by a feeling of wretchedness), and he stresses the questionable character of
the tendency prevailing in the Church to make the negative type the one right and normal one. Those who struggle
for fuller light must be treated differently from those who struggle to be freed from darkness. In the material
collected by Leuba too the positive type occurs most frequently. Starbuck and Leuba do not discuss the cause of this.
It is not quite certain that the positive type is in reality the most frequent; perhaps the negative type is more shy
and disinclined to express itself, so that therein lies the reason why fewer cases of it can be ascertained.

Of very great interest is the interaction of the conscious and the unconscious in the breakthrough. There are many
different degrees of consciousness here, which probably is connected with there being different degrees in the inner
hindrances that are to be overcome. A presupposition for the breakthrough is --- within both forms --- that a growth
has taken place beneath the threshold of consciousness (within the sub-liminal consciousness). Now it is not given that
the direction of the conscious striving coincides with the direction of the unconscious growth, --- and since the assent
of the will to what is unconsciously laid down is nevertheless necessary for the breakthrough to occur, it is no wonder
that so many psychological riddles present themselves here. What is to be brought about is, after all, that what has
hitherto stood only as object, as a distant and exalted ideal, is to become that very thing which lives and stirs in the
individual. What has hitherto been only object is to become subject.

Once the breakthrough has occurred, a feeling of newness arises, as a rule also of joy and peace, and frequently of
self-forgetfulness and love of humanity. What has hitherto held the self back within narrow barriers is displaced, and an
expansion can take place.

Psychologically viewed, conversion in both its forms is not a
% --- p. 105 ---
\setcounter{footnote}{0}%
\opage{105}unique phenomenon. There is a series of analogous experiences, as when uncertainty and dejection are suddenly replaced
by cheerful decision, or when a long effort to acquire a skill is as with one stroke replaced by mastery over the difficulties,
or when a state of feeling announces itself suddenly through contrast or prepared by inner organic processes. Such phenomena
occur especially in the
age when the young person is to live himself into the habits and traditions of the race and make them his own --- is to expand
his self to embrace larger interests, break down the barrier between the individual and the social will. Physiologically viewed,
this task is set at a time when growth is approaching its completion, and when new tasks and functions begin to demand their due.
It is a time of movement and of oppositions; both in the organism and in consciousness the opposition between the old and the new
emerges sharply. But through religious influence this opposition is often strengthened in an unnatural and unhealthy way. One does not
always recognize, as already remarked, the difference between the two types, but treats them alike. Starbuck cites an example of how
it can lead to despair when all are to be forced into one definite scheme of development and are compelled to make experiences that are
impossible for them. ---

In opposition to the breakthrough type in its two forms Starbuck sets \textit{the continuity type}. There is here no consciousness of
decisive change, of opposition between old and new. An involuntary awakening takes place, which indeed has a certain likeness to the
positive form of breakthrough, but lacks the experience of sudden change. Development proceeds continuously from childhood through the
young years. Yet a period of ``storm and distress'' (storm and stress) can come, which is connected with the organic changes of the years of
transition and can bring anxiety and tension. It is the period when the young person is to live himself into a larger world and make the
acquisitions of the race his own. Manifold demands and new possibilities present themselves. In intellectual form (and somewhat later) the
unrest thus awakened appears as doubt. The point is to hold fast with clearer consciousness to what one had hitherto
% --- p. 106 ---
\setcounter{footnote}{0}%
\opage{106}accepted on authority and with childlike credulity. Such a period was very frequent among those whom Starbuck came to know more closely;
it could begin, run its course and conclude without abrupt transitions. Often it led over to a time when the individual stood estranged toward what
had been handed down and regarded it as wholly worthless. This period of estrangement (period of alienation) was often due to relaxation after the unrest or doubt;
but it could also have its ground in the individual's now wishing to assert himself against the forms society wished to lead him into, or in delight in the
unrestrained use of faculties and powers. In place of the religious tradition aesthetic, intellectual or other interests often step into the foreground in such a
period. At the end of the period of youth a reconstruction as a rule sets in: the individual has won a definite faith, which does not always agree with the received
one, but as a rule is a mixture of independent production and resumption of the old. Starbuck is astonished to find this period so definitely marked. He explains it
by life's circumstances demanding a definite standpoint in the practical respect, so that perhaps even in the race a tendency to such reconstruction after the agitated
years of youth has developed. --- Through all these periods development within this type proceeds so continuously that it is often only through comparison of points
lying far from one another that one becomes aware of them.

Characteristic of the conclusion of development is --- and this holds for both main types --- a certain change in the religious ideas and in the whole conception
of religion. Even though the confessional ideas are essentially retained, a pervasive shift can nevertheless as a rule be demonstrated from the outer to the inner,
from the dogmatic to the ethical, from the isolated personal to the consciousness of participation in a more comprehensive life. Religion is conceived more and more
as an inner life that is part of a universal order of things.

\textit{If we now compare in conclusion} the breakthrough type with \textit{the continuity type}, the former appears, as is natural, with far greater unrest and with
``storm and distress'' to a more marked degree than
% --- p. 107 ---
\setcounter{footnote}{0}%
\opage{107}the latter. The breakthrough type is more receptive to outward influence; the opposition between the old and the new asserts itself more strongly in it, and the life of
feeling is more vehement. The breakthrough type moreover adheres far more than the continuity type to the definite confessional forms, to the organized institutions. Its inner
experiences it prefers to express in biblical and confessional terms and has no need to form independent forms and expressions or find new analogies. Such a need, on the other hand,
is very frequent within the continuity type, where also the conception of religion as an inner life occurs more frequently. --- Starbuck's general result is that the development of the
religious life under confessional forms can be illuminated from general psychological and physiological points of view. Dispositions and impulses of general human nature unfold in religion,
and also the different forms in which the religious development takes place are known from other mental domains. He himself regards religion ``in its highest form'' as a radiation, an interplay
and an idealization of tendencies that are already present beforehand in human nature. The forms and expressions the race has found for this are what the young person is seriously led into in the
years of transition, and the course and crises of development are conditioned by the difficulties to be overcome here. He is inclined to think that the inconveniences that appear in the course of this
development are due far more to the older generation's lack of understanding of the demands of human nature than to oddities and faults in the young. --- With this last remark may be compared a statement
by Professor PEABODY in his treatise \textit{The religion} of a college \textit{student}. (Forum. June 1901. p. 447): ``For one young man who is removed from religion through the influence of science,
philosophy or art, there are quite certainly ten who go thus because of unwise or unpractical religious instruction\dots{}. It is not the University that is to renew the agreement with the Church; it is the
Church that must open its eyes to the wonderful expansion of the intellectual horizon which every student nowadays has access to experience.''

% --- p. 108 ---
\setcounter{footnote}{0}%
\opage{108}

\subsubsection*{\textit{b. The Influence of Temperament.}}

GEORGE COE too stresses the importance of a psychological study of religion in its present state. And it is also
the period of youth to which he has especially directed his attention. This period is, he maintains, far less known
and understood than childhood and mature age. What takes place in this period is a transition from being taken up
with the particular to grasping the universal, from being overwhelmed by the outer to immersing oneself in the inner,
from letting oneself be led by authority to being determined by reasons, and from stopping at the given to striving toward
the ideal. Questions are asked more searchingly about the significance of life, and the social interests begin to awaken.
All this brings it about that a need is felt for a more personal appropriation of the content of religion than there is
possibility for in childhood. Conversion or a corresponding process of inwardization is therefore a natural link in the
development of the period of youth, and Coe also in fact finds that such a change is extraordinarily frequent (especially
at the age of 16). He recalls that among peoples around the earth it is common to place a religious initiation at this time.
Often the young man is sent out into solitude to wait for a revelation from the god of the tribe and by its help to become
clear about his calling. In agreement with Starbuck he maintains that among young Americans within Christianity it is by no
means merely fear or remorse that are motives in religious development, but also more indeterminate states: unrest, emptiness,
feeling of incompleteness, longing toward the infinite, --- even though most express their experiences in the confessional
language. It further appears throughout that purely intellectual problems and motives are subordinate; the reasons are found
afterward; one seeks first to live the life that expresses the deepest self, and afterward one perhaps looks to see how one can
by the help of reason show its justification. And it is not any single experience that leads to the goal. Most of what youth
demands decisively all at once, one can only live oneself into gradually, just
% --- p. 109 ---
\setcounter{footnote}{0}%
\opage{109}as life unfolds with its work, its disappointments and its victories.

Coe wishes to study these phenomena without theological and metaphysical prejudices. The psychology of temptation, e.g., is and
remains the same, however one may explain to oneself the origin of the impulses that carry it with them, --- whether one thinks that
they derive from the Devil, or that they are remnants of the animality from which humanity has developed. The same holds of the
psychology of conversion, with which Coe has especially occupied himself. It is his task to point out the most important factors which
co-operate in this psychical phenomenon. He goes here the way of making an exact investigation of a relatively small series of individual
cases. Besides sending out question lists he used detailed conversation with individuals, procured others' testimony about them, used hypnosis,
sought by all possible ways to arrive at an individual characterization. The number of individuals he investigated was 77 (52 men, 25 women);
they were almost all students, the men on average 24 years old, the women 22 years; most had been brought up in the Methodist church, which
indeed lays so great weight on personal religious development.

It then turned out that in such young people great differences appear with regard to conversion. In this respect three classes of individuals
can be distinguished: A) those who expected a transformation (transformation) and really experienced it; B) those who expected a transformation,
but did not experience it; C) those who were not satisfied with what they experienced, or who nevertheless strove to attain a fuller experience.

On closer investigation it now appears that class A were predominantly people of feeling (sensibility distinctly predominant), whereas class B were
predominantly people of intellect (intellect distinctly predominant). And whereas class A most closely recalled the so-called melancholic and sanguine
temperaments (the slow-intense and the prompt-weak varieties), class B recalled more the so-called choleric temperament (the prompt-intense variety). Class C presented
% --- p. 110 ---
\setcounter{footnote}{0}%
\opage{110}differences that recalled now A and now B. Temperament thus proves to be an essential factor.

With this there now combines another observation. In class B there occurred next to no automatic phenomena in connection with the religious states,
neither in the domain of the life of ideas (strong dreams, hallucinations), nor in the domain of the organs of movement (irrepressible laughter,
starting, clapping of hands, talking aloud to oneself, etc.), whereas such phenomena occurred far more frequently in class A (and in C). --- When the
classes were moreover examined with regard to susceptibility to suggestion and hypnosis, it turned out that A's suggestibility is predominantly passive,
in that the individuals here quite conform to the prompting given and follow it in detail, whereas class B's suggestibility can be called spontaneous,
in that the individuals either give themselves their promptings or independently add something to the original prompting.

Besides the aroused expectation, temperament, tendency to automatism and suggestibility thus play a role in conversion. The character of the religious
experience is determined by these factors. ---

Like Starbuck, Coe stresses very strongly the danger of regarding a single type as alone justified. If an expectation is aroused of experiencing a more or
less sudden upheaval of mind, and the fulfillment of this expectation is, according to the individual's whole psychological nature, impossible, then it can
lead to despair. Coe thinks that really it is only the melancholic and the sanguine temperament that the Church has recognized in its worship, its hymns and
its morals. In the type of Christ that prevails in the Church this shows itself especially; it is the Man of Sorrows, the suffering sacrificial lamb. The real
Jesus suffered because in spite of all opposition he held fast his thought to the utmost; he was an active, not a passive individuality; there was a choleric
element present, which has vanished from the traditional image. --- The Church must be as broad as human nature is. But it has accustomed itself to look at things
with feminine eyes: therefore it is especially the men who turn away from it.

% --- p. 111 ---
\setcounter{footnote}{0}%
\opage{111}Coe finds, in spite of his evident sympathy for religion, no scruple about a thoroughgoing psychological discussion of the religious states. If we --- he says ---
want to understand religious experiences, we must not explain them by God's unsearchable ways, but by demonstrable peculiarities in the nature of man; we are not to speculate
about divine grace, but to study the hidden life of the human mind. Like other philosophers of religion in our day he finds the final measure for religious values in the ethical.
What matters is what is gained by religious development for a deeper, more richly contentful active life.

\subsubsection*{\textit{c. Differences in the Relation to the Confession.}}

There are here, as already mentioned, two investigations, one by Leuba on American Methodism, another by Arréat on French Catholicism.

LEUBA has studied partly the religious life as it appears among contemporaries (in his native country North America), partly the religiosity of the mystics. Both studies lead
him to make a definite distinction between the standpoint of the historian of religion and that of the psychologist of religion. The historian's subject is not so much the religious
experiences themselves as the products of these experiences present in worship, myth and dogma, as they take shape within a society. To arrive at the immediate religious experience
one must turn to the individual religious consciousness. By it all religion is in the end borne and conditioned, and in the higher popular religions this relation becomes especially
clear. It is not the gods but the relation to the gods that interests the psychology of religion. Its task is to give a religious dynamics, i.e. a doctrine of the motives of the
religious life. (The expression ``religious dynamics'' had been used by Coe before Leuba). To obtain material for this from contemporaries Leuba has by sent-out question lists obtained
utterances from a series of persons who stood at very different standpoints, but all within the Anglo-Protestant world. The questions concerned the grounds of their participation or
non-participation in
% --- p. 112 ---
\setcounter{footnote}{0}%
\opage{112}worship (public or private worship), --- the circumstances and localities that exerted a special religious influence on them, --- the most characteristic religious experiences
they knew, --- the differences they made between religious and non-religious states and actions. He has succeeded in collecting a series of characteristic utterances. Statistics is not
in question; but Leuba thinks he can demonstrate certain peculiarities of the religious life in the circle he has to do with.

A main point is for him the result that what expresses itself in religion is not a striving for knowledge, but a striving for being. Religion springs, in his view, from a feeling of
fragmentariness (unwholeness) in connection with a feeling of inner division and disruption. The former feeling brings with it a need to become a partaker in a larger, more comprehensive
life, whereas the feeling of division brings with it a need to be led by a single end (unification of motives). The feeling of sin is a form of this latter feeling. To the two motives mentioned
are joined further altruism and the feeling of dependence: the first is related to the need for unity, the second is conditioned by the inner resistance that is to be overcome in efforts to
attain unity. The great founders of religions paid no heed to theoretical problems, but strove after ideal perfection and inner unity. And there manifests itself nowadays more and more a need
to go back to the immediate religious experiences from which religion has always from the beginning grown. Although one has not yet had courage or intellect enough to do this fully, still the
first steps in this direction (to return to plain religious empiricism) have been taken, especially within Protestantism. In silence there is waged --- alongside the struggle for existence
between individuals --- a purely inner struggle, which in severity and tragic seriousness often surpasses the outer. It is waged in different ways within the different religions and finds
expression in different forms: but always there is bound up with it the need for a greater and at the same time more undivided life. The religious need is really no quite special need, but has
its deep and lasting psychological and physiological foundation.

In comparing religion as it appears
% --- p. 113 ---
\setcounter{footnote}{0}%
\opage{113}in more recent times in contrast to earlier forms of it, Leuba finds a series of differences.

In primitive societies religion is a composite product of all of man's fundamental wants and strivings, both theoretical and practical. The theoretical motives have now separated out
and find their satisfaction in science; it is not on them that the significance of religion in more recent times rests, and already in antiquity Buddhism stands as testimony that these
motives are not the decisive ones. Even if all intellectual elements hitherto attached to religion fall away, the psychological essence of religion, as it has been described above, will be
able to persist.

A further significant difference between more primitive and more modern religion Leuba finds in this, that whereas formerly the feelings of obedience and awe were predominant, these are now
--- in any case in the Anglo-Saxon world --- not so frequent as is assumed. Not the obedience that the catechism inculcates, but the clear understanding of what is right and good, is now the motive
to action, and the feeling of unity and freedom is displacing submission. Connected with this is the feature that religious truth is no longer regarded as something that can and ought to be grounded in
an objective way; the decisive thing is what satisfies the personal need and yearning. Many of those questioned expressed this with the greatest candor. The motives of feeling had decidedly the
preponderance over all intellectual interests, and there is perhaps even complete indifference toward the purely objective significance of the religious ideas, even in individuals who otherwise in
more specialized fields showed themselves in possession of critical sense and capacity. One does not know God; one does not \textit{understand} him; but one \textit{uses} him. Not God, but life, more
life, higher, richer, more satisfying life is the end of religion. In the past this shows itself
clearly in the so-called mystics, and here Leuba's studies of mysticism connect with his studies of
contemporary phenomena. The mystical life must not be interpreted and judged by the concepts and dogmas the mystics employ. Its peculiarity is only a more extreme and concentrated form of the
% --- p. 114 ---
\setcounter{footnote}{0}%
\opage{114}peculiarity of all religion: a yearning for higher and more complete satisfaction, for peace and unity in the will. The mystics' yearning is more passionate than that of ordinary
religiosity, and the concentration upon a single goal is stronger. Goal and direction are the same, but they proceed with far greater energy.

Yet another characteristic feature of the religiosity of more recent times Leuba finds in this, that so many and considerable individual differences appear, even within one and the same civilization.
Under the influence of the infinitely larger horizon life has acquired since the time of the Councils, religion too has undergone changes, and differences show themselves not only in the outer things
about religion, but in the individuals' innermost mental life. The wants are different, and the hopes are different. And yet in theology one proceeds from the assumption that one and the same food suits
all. ---

Within his collection of self-confessions Leuba demonstrates a whole series of nuances, from a standpoint where the religious and the confessional coincide to such a degree that the individual is not able to
express himself about his inner experiences except in biblical and confessional terms, to such standpoints where no meaning whatever can be connected with that kind of expression. Between both extremes lie the most
interesting cases; for neither from dogmatic position nor from dogmatic negation is there any considerable psychology to be learned. Only a few features can be cited here from the interesting self-characterizations
Leuba has collected.

A scientifically educated young woman, adherent of the Presbyterian church, distinguishes definitely between what is due to her confessional upbringing and what is due to the need for something that has eternal
value (a desire for things of permanent value). These two things do not exclude each other, but neither do they coincide. When she takes part in the worship of her church, it is partly from habit, but partly also
because she needs to enter into relation with a higher being and to be freed from selfish thoughts. The difference between religious and non-religious thoughts and feelings is for her a difference of degree and
rests on the fact that the former appear
% --- p. 115 ---
\setcounter{footnote}{0}%
\opage{115}as more significant and stronger. But the ideal need is never fully satisfied. ``Whereas the current religious assertions declare religion to be the only thing in the world that can satisfy,'' she says,
``for me its value rests on its constantly driving me forward to new exertions.'' --- Other correspondents emphasize love of humanity, childlike mind or gratitude as the essential in what they call religion.
--- A physician, who takes part in the worship of the Church, although his own private devotion has for him a more decidedly religious character, states among other things the following: ``I am far from having the
anthropomorphic idea of God that I had in my young years. By God I understand the power or the principle that brings about all development in the universe toward a higher order of things, and which by its laws raises
the stages of life from the unconscious to the conscious, from the animal to the rational, from the perishable to the imperishable. In myself I call divine all drives, tendencies and exertions that lead to developing
in me greater receptivity to beauty, truth and justice, higher reason and capacity for moral progress''. --- A related statement from another correspondent is the following: ``I have given up every attempt to imagine God
and no longer feel any need for an image or a clothing of the divine principle which it is a necessity for me to believe in. My religious relation to what one calls God I can describe from two points of view: 1) I feel
in my innermost consciousness a need to be One with a
great principle, a need for unity with Something that I cannot define, but which surrounds me, in which I am so to speak immersed, and with which I wish to be in harmony, an all-embracing
Present\dots{}. 2) I have almost uninterruptedly a deep consciousness of standing far below an inner ideal, a type that is laid down in me and that I constantly have before my eyes. I feel that I \textit{must} realize this
type, which constantly dwells in me and seeks to lead me. Experience has taught me that I become happy by following it, unhappy by moving away from it. To these two fundamental instincts or intuitions I have reduced all
my religious ideas. I am a disciple of Christ
% --- p. 116 ---
\setcounter{footnote}{0}%
\opage{116}and strive to imitate him; my ideal is to become One with the Divine, as he was. But there are many of the assumptions of the great Christian communities that I do not share.'' --- A woman who lives with husband and
child in the country in constant struggle with poverty declares that for her religion is truth, and by truth she understands the law that holds for our being. ``When I say God, I mean the Good, the law of my being\dots{}
I believe that life is development, growth\dots{} and that we have as much heaven and as much hell here in life as we ourselves will\dots{} The Good is our ideal; the Evil is our teacher, who is to make us understand the
Good\dots{} Human beings cannot be divided into good and evil, but into developed and undeveloped\dots{} I know no religious act except life itself\dots{} Perhaps I was nearest to religious experience when it first became
clear to me that I was free, free from faith, doctrine, form and sentimentality, free to grow, to become All that lies in my nature\dots{} If you cannot find any religion in this, you will hardly find any religion in
me\dots{}. I am against the doctrine of the imitation of Christ, because it conflicts with independent thinking, and because mental and spiritual development rests on independent thinking. We might just as well let
another eat for us as think for us. Enthusiasm one mind can inspire in another, but not Existence\dots{}. Christianity has been a great factor in the development of the human mind, but the Christianity of dogmas one has
now almost outgrown\dots{}. We are not about to lose our religion, but we are about to find ourselves. We conquer fear and learn to have trust in the divine law that rules over life.''

All these examples belong to the group that thinks itself positively religious, although it is often, especially as regards the last example, doubtful with what right. But they are characteristic expressions of the need
to give an account personally of what one --- the single individual --- can find of value for his life in the given positive religion. Through some transitional cases of more indeterminate form we come in Leuba's series over to
decidedly non-religious forms, which --- however far one may stretch the concept of religion --- cannot come under
% --- p. 117 ---
\setcounter{footnote}{0}%
\opage{117}it. Partly there is lacking here religious experience in any sense of the word, partly there are definite intellectual and moral grounds that lead away from religion --- admittedly because religion and the Church system are
assumed to coincide. Thus one correspondent remarks: ``I deny the moral superiority of the churches. They never take the lead of a moral reform. They were on the wrong side in the slavery question. They are cowardly, petty and
servile in the great industrial questions of our time.'' --- Leuba thinks that the number of non-religious persons is far greater than is commonly believed. When there are not many utterances from such persons in his collection, it is,
in his opinion, partly because his question lists have not reached them, partly because they are not always clear about their real standpoint.
---

Leuba has not himself concluded the analysis of his material; his treatises available so far are only results of a first treatment, and he ends really with an unanswered question. In the fundamental religious motive, the religious
tendencies he finds nothing that is peculiar to religion, but particular forms of the need for self-preservation and self-development. This conception he finds decidedly confirmed by his correspondents' utterances. He rejects every
conception of religion according to which it is to aim at knowledge, whether this knowledge is to give a revelation or to establish a mystery. (The Monist. July 1901. p. 571---573). But at the same time as he in his latest treatise
(Revue philos. Novbr. 1902.
p. 487) refers back to these statements, he states: ``What then is the element that distinguishes the religious life from the secular (for there is, after all, a difference between them)? It is the one means or the several means that man
employs to attain his wishes, and not these wishes themselves, that separate religion from the rest of life. Psychologically expressed: it is not the will (in Schopenhauer's meaning of the word), but knowledge, that separates them.
When one realizes that the difference between the religious people and those who are not (insofar as there are such) does not consist in striving, wish and tendency, thus not in the very sources of life, but only in the means
employed to reach the same
% --- p. 118 ---
\setcounter{footnote}{0}%
\opage{118}goal, religion appears in another light than that in which it ordinarily stands in people's eyes.'' --- A closer explanation of these hints will hopefully be brought by Leuba's following works. But when he definitely
rejects the intellectualist conception that it should be the \textit{end} of religion to give knowledge of any kind whatever, and yet thinks that the difference between religion and other spiritual forms of life \textit{rests}
\textit{on} knowledge, then the meaning can well only be that the difference rests on the way in which knowledge is used in religion as \textit{means}. And means must here mean means of interpretation, that is, symbol. I also
believe that Leuba would be able to draw from his material a result in this direction. There is in all religion a need to form what is experienced in thoughts.
In confessional religion this need is satisfied immediately by the received ideas; in the transitional forms other expressions are sought to interpret the significance of these received ideas themselves for life. At the base there
always lies an experience of the relation between ideal and reality, an experience that --- wherever religion at all arises ---
passes over into faith and trust toward powers that are assumed to work in the direction of a harmonious relation between ideal and reality. I find my hypothesis on the essence of religion as faith in the persistence of values
confirmed by the investigations of Leuba and his colleagues.

On the other hand it seems to me that Leuba lays too little weight on the influence of prevailing or received ideas. It is fateful for the whole character of a feeling to become attached to a definite circle of ideas; it thereby acquires a direction and a content that decide its quality, and when the connection between it and the ideas it has hitherto been
attached to is to be torn apart, it can happen, or at any rate threaten, that the feeling disappears at the same stroke. One can indeed with Leuba distinguish between ends, which can be said to be common to religion and other
spiritual forms of life, and means, which are different for the former and for the latter; but it must then surely be maintained that the very fact that an end can be attained by different means indicates a difference in the end itself.
% --- p. 119 ---
\setcounter{footnote}{0}%
\opage{119}There is no wholly external connection between end and means in such domains as this concerns. The sting of the religious problem lies precisely at this point. History teaches, after all, also how hard human beings find it to put
up with differences of quality within religion, and the dispute over the limits of the very concept of religion rests likewise on the question whether through a difference of quality we get outside religion's whole modality (to use the
terminology of the psychology of the senses). It becomes thereby of particular significance that there are natures that are able to find expression for their inner life only in \textit{common} symbols, and who, if they are compelled to give
them up, feel wholly homeless.

Yet it must be repeated that Leuba has not yet completed the working up of his material. It will be of interest to hear the words with which he concludes his preliminary communications. ``Whether the differences between religion and
non-religion are only superficial, or whether they are fundamental, and in what degree, --- whether upbringing in the home and other circumstances in and of themselves can explain these differences, or whether they have deeper roots in particular
innate dispositions, --- whether the one is destined to displace the other or is merely a preparatory stage for it or not, --- it is a question that must be passed over in silence in this preliminary treatise. Only this I wish to emphasize here:
even between individuals of the same society and of about the same intelligence and culture there are differences that are sufficient to make the one a human being for whom religion is the All of life, and to leave the other standing wholly
indifferent and inaccessible to it, even if it be interpreted in the widest sense.'' ---

LUCIEN ARRÉAT has --- by similar paths as the American psychologists --- sought to find different types within French Catholicism. One notices at once that the circumstances are different. The difference between Catholicism and Protestantism
is undoubtedly connected with racial differences. Whereas from another side (Charles Renouvier) it has been regarded as the fall of the French people that it did not join the Reformation, Arréat maintains that this conflicted with the
% --- p. 120 ---
\setcounter{footnote}{0}%
\opage{120}peculiarity of the French race. Protestantism offended the sense of beauty by its destruction of images and buildings, and what it robbed the imagination of it gave no compensation for in the domain of thought.
Protestant individualism demands of the single individual an independent discussion and understanding, which conflicts with the Latin race's disinclination to theological discussions. The Catholic sees the doctrine fixed once for all,
and his inwardness breaks forth in the personal relation of feeling, the enthusiastic adherence to the fixed eternal truth which the Church guarantees. Arréat has therefore also found in his material no hint of the transition from
feeling of obedience to feeling of freedom in the relation to God, which Leuba thought to find in Anglo-Protestantism; Catholic religiosity is essentially submission. --- According to his observations Arréat describes four different
classes within French Catholicism.

The first class, which is found especially in marked form in southern France, comprises the largest number of Catholics and has representatives in all classes of society. Piety consists here in observance of the religious customs,
which are conceived in a way that approaches idolatry, which shows itself especially in the strong flocking to the places where miracles occur. Here the racial character is at work whose inclination to superstition Caesar already mentioned.

The true Catholics make up the second class. Through participation in the means of grace of the Church they strengthen the feeling of security of being on board St Peter's ship, which in the midst of the unrest of the times is steered by an
infallible helmsman. The Catholic religion is praised for its sublimity and its captivating beauty, for its ideal of perfection and for its active love of humanity. Without need of mysticism or of independent thinking there expresses itself
here a mental equilibrium, which is nevertheless sometimes connected with an express conviction that it is only the mighty authority of the Church that keeps doubt and danger away.

To a third class are reckoned those who support their ecclesiastical faith by more or less conscious thinking. Many clergy belong here, and to this class Arréat thinks he can also reckon many who declined to express themselves. One
% --- p. 121 ---
\setcounter{footnote}{0}%
\opage{121}protests here against religion's being a matter of feeling: it grounds itself on rational knowledge, although grace brings the final decision. The Catholic doctrine stands as a whole from which nothing can be separated, and which
immediately captivates the soul. Yet it is sometimes said that the Church is only the outer and perishable garment of faith, which will be able to be transformed into new forms; the properly religious is placed in the need for the infinite.
By such statements the transition to the last class is made.

In many men and women who declare themselves good Catholics a certain state of dissolution can be traced. They choose, e.g., among the Church's doctrines and customs and hold to what they have use for. Often it is only an express act of will
that holds them back in the Church; one stresses the social motives, the opportunity to practice love of humanity, e.g., as that which binds to the Church. Or one conceives the religious as ideal striving, pays homage to a practical idealism
without metaphysics and uses the Church's worship only as a means of mood. One expresses the most inward sympathy with the believers, because they do not, like many free-thinkers, give up the sublime poetry of religion or are blasé toward
the mighty problems to which the word God points. A last step is taken when one goes over to Positivism or by another path seeks a wholly new expression for one's view of life. This is, however, unusual. Either one returns to the faith of
childhood after a time of unrest, or one sinks in spiritual respect to a lower level.

It is --- Arréat says in conclusion --- impossible at the moment to say whither the present crisis in the religious domain will lead, or what energy the religious feeling will be able to possess when it is not supported by faith in dogma or
poetic myths: ``I see what one wishes to eradicate from modern mental life; I do not see what is offered for its enrichment.''

% --- p. 122 ---
\setcounter{footnote}{0}%
\opage{122}

\subsubsection*{\textit{d. Poetry and Religion.}}

GEORGE SANTAYANA, in his work mentioned above, has discussed the
relation between poetry, religion and science in a way that
gives him a place in the ranks of the psychologists of religion. For
him it is the imagination, the image-forming faculty, that underlies
all those three spiritual manifestations of life. The images
that we can use in practice, and with the help of which we can predict
future experiences, we call ideas of the intellect, whereas
images that cannot be exploited in this way continue
to be called ideas of fantasy. Intellect is a fantasy that makes
empirical science possible. Under the influence of experience we choose
among the images that arise involuntarily. But among the
fantasy images that cannot pass over into ideas of the intellect,
two kinds must in turn be distinguished. When fantasy images,
although they do not express empirical reality, nevertheless stand as expressions
of life values, ideals and ends, they receive the character of religion,
whereas they stand as purely poetic forms when they have no
such significance for life. Religion and poetry are in their innermost
essence akin, indeed identical, and differ merely
in their different relations to life. Poetry at its height passes over into religion when it expresses the
value of our experiences and gives great, guiding and inspiring anticipations.
Conversely, religion passes over into poetry when it gives up the presupposition of containing metaphysical truths. What matters,
in order rightly to understand this, is to distinguish between the world of values and
the world of facts: value rests not on the given itself
but on its significance.

Our will is infinite, but our experience is limited.
Poetry leads us beyond this limitation and thereby makes possible a
higher development and exercise of the will, in that its highest
end can now appear with greater clarity through poetry's help
than mere experience allows. A striving, a passion
we understand rightly only when we see what is sublime in its whole and full
unfolding. But because poetry becomes religion, it does not cease
to be poetry. On this point not only the orthodox but
their critics as well are guilty of error. Already the Greek
% --- p. 123 ---
\setcounter{footnote}{0}%
\opage{123}philosophy's criticism of polytheism was unjust, in that it took no
account of mythology's close connection with life. A great
religion is a product of the human spirit in its entirety, and in
relation to that, criticism, which is guided by the special, intellectual
interest, stands as limited and piecemeal. In Christianity
historical idealism took the place of the naturalism of the ancient religion; the harmony that the old religion had
sought to represent in space and in nature, the new sought to
represent in time and in history. Christianity obtained its foothold in history
by having its root in a national faith
that had taken shape under the trials of the Jewish people; but out of this
it formed a great poetic expression of experiences that human beings
knew from the innermost depths of their souls. It was the great
autobiography of humanity that was here made the center of the life of the world.
The significance that a religion has and retains is not necessarily
the one that it had for its founders; they are so captivated
by their visions that they regard the content of these as external
facts. Religion is not different from superstition in
psychological respect; the difference rests not on origin but on value.
Even if ethical value was not the ground of the coming into being of religions,
it is nevertheless the ultimate standard
of their significance. ---

This whole line of thought, which Santayana carries out in detail
through interesting characterizations of religious and poetic
works, acquires a special interest when one compares it
with some of the observations of Arréat, Leuba and Coe, and when one
at the same time recalls the idealistic or mystical tendencies that are stirring
in the literature of the most recent time. What I miss in Santayana
is a more penetrating psychological analysis of the relation
between experience and faith; but otherwise it has been of
great interest to me to find a way of looking at things akin to my own
in a work that appeared at about the same time as
my own attempt in the philosophy of religion. For comparison I permit myself
to quote a single statement from my book (Religionsfilosofi
p. 334 f.): ``Perhaps it is the case that poetry
is a more perfect expression of the Highest than any scientific
concept could be. By poetry, then, one must not here
% --- p. 124 ---
\setcounter{footnote}{0}%
\opage{124}think of loose moods and fantasies, but of the involuntary
and vivid form in which what is really experienced clothes itself
in strongly moved moments\dots{} It is a poetry that neither
stands in opposition to will nor to thought, although it comes especially
into being where thought and will meet with barriers\dots{}. All high
art presupposes such a poetry of life and is its conversion
into clear forms.''

\subsubsection*{\textit{e. Comparison between Different Series of Types.}}

In my Philosophy of Religion I have, relying especially on a biographical
and historical material, presented a series of typical differences
in religious faith (p. 111—120) and, in another connection,
a series of differences with respect to the original dispositions
that condition the peculiarity of religious life in
individuals (p. 251—256). If one compares these series with
the types that the American psychologists of religion have set up,
one will find above all one opposition, which plays a leading role in all the presentations,
namely the one that I have described as
an opposition between faith as concentration and rest after
the balancing out of disharmonies and faith as the highest form of the unfolding of life.
Corresponding to this, among the dispositions, is the opposition
between disharmonious and expansive natures. In Leuba,
Starbuck and Coe alike this is the most prominent opposition.
But it also appeared that they lay quite special weight on
the expansive type and are inclined to regard the disharmonious
or discontinuous type as a special case
of it. There is of course always some resistance to be overcome,
some crises to be passed through, and it is only when resistance and
crises become particularly strong that, in the opinion of these investigators, we
pass over from the expansive to the disharmonious type.
As to the reason why the expansive type occurs so frequently
in their material, I have expressed a conjecture above.
One might perhaps add that the authors mentioned (and the
same evidently holds of Santayana) themselves personally decidedly
seem to belong to the expansive type, in so far as one can come to know
their psychical habitus from the way in which they
% --- p. 125 ---
\setcounter{footnote}{0}%
\opage{125}perceive and evaluate those of others. For my own part I would be
inclined to regard the difference between the two types as greater
and more fundamental than those authors assume. Perhaps
there lies in this opposition the germ of a spiritual struggle that again
and again will come to stand between spiritual directions of life.
In the deepest depths they ultimately do not understand each other. In Danish
intellectual life we have a classic example of this opposition in
the relation between Grundtvigianism and the Inner Mission.

Were I to reduce my series of types, this opposition
would nevertheless remain. I have in all seven types, namely
besides the two already mentioned [faith as 1) rest and as
2) unfolding]: 3) contemplation, 4) confident venturing, 5) resignation,
6) convulsive religiosity and 7) collier's faith (fides implicita).
Of these, 3, 4 and 7 would fall under 2, whereas 5 and
6 would belong under 1. As regards the differences in dispositions,
I set up, in the domain of the life of feeling and will,
the following series of oppositions: between the Disharmonious and
the Expansive, between the Passive and the Active,
between the Affective and the Continuous and between the Idiopathic and
the Sympathetic, --- and in the domain of knowledge the opposition
between analytic and synthetic natures and between discursive
and intuitive natures. It is easy to see that there is an analogy
between all these relations of opposition, and that they can all be subsumed
under the first, so that it is the opposition between
the Disharmonious and the Expansive that, seen from this side too,
comes to stand in the first rank.

In the conviction that in the future, to a far higher
degree than hitherto, in understanding and judging religious
phenomena, weight must be laid on individual differences,
I am in complete agreement with my American
colleagues.

\subsection*{4. COMPARATIVE PSYCHOLOGY OF RELIGION.}

A comprehensive work in the field of the psychology of religion
is to be found in the American philosopher WILLIAM JAMES's
% --- p. 126 ---
\setcounter{footnote}{0}%
\opage{126}book \textit{The varieties of} religious experience. A study in human
nature. London, New-York and Bombay. (1902). James is among
the most outstanding thinkers of our time. He unites extensive
knowledge with great power of observation, sharp criticism with
idealistic enthusiasm, freedom from prejudice with earnest conviction. The most important of his writings is \textit{Principles of psychology} (1890), which on several points of psychology brings new points of view to bear and throughout discusses the questions in a spirited
and stimulating way. In a collection of essays that
has appeared under the title \textit{The} will \textit{to believe} \textit{and other essays}
in popular philosophy (1897), he has in somewhat freer form expressed
his peculiar union of philosophy of experience and a pronounced
idealistic view of life, and already here he enters
upon questions of the \textit{philosophy of religion}. The new work is wholly
devoted to the religious problem, in particular to the purely psychological illumination
of it, yet in such a way that the other points of view from
which the problem must be treated also come forward. James makes use of
the specialized works of the younger American authors, but in addition draws
on an extensive literature of a biographical, especially autobiographical
kind. He sees clearly that the subject is too manifold
to be exhausted in a single presentation, and therefore --- also
in order to escape a dispute about words over what religion is --- he declares
that for the purposes of this his investigation he will
keep to a definite meaning of the word religion or to
a definite side of religion, and he then defines religion
as: the feelings, acts and experiences of individual human beings,
in so far as these human beings in their solitude conceive themselves
as standing in relation to the divine, however
they may conceive this. This personal religion he sets in opposition
to institutional religion with its more or
less developed theology or mythology and its ecclesiastical organization.
He maintains that \textit{institutional religion} (with all that
belongs to it) is always derived in relation to \textit{personal} religion; that
follows already from the fact that it builds on tradition. If one
were to keep to institutional religion alone, the investigation
would come to concern merely worship, sacrifice, ceremonial
and ecclesiastical organization, and one could then define
% --- p. 127 ---
\setcounter{footnote}{0}%
\opage{127}religion as the art of winning the favor of the gods. In personal
religion, by contrast, it is man's own inner states,
his conscience, his worth, his helplessness or his
imperfection, on which the weight falls. If one keeps to this
side of religion, the concept of religion acquires a greater scope
than when one regards the institutional as the main thing.

Personal religion is the way in which man
reacts to life as a whole (a man's total reaction upon
life), in opposition to the momentary and special reactions.
The experiences we have had of life and of the world act upon
our individual temperament and make us either energetic or
indifferent, either pious or mocking, either dejected or
enthusiastic toward life as a whole. Our ``reaction,'' which often
takes place involuntarily and half unconsciously, expresses itself in the most complete
answer we can give to the question: What is
the character of this existence to which we belong? --- Yet James
will call such a reaction religious only when it is borne by a
serious and exalted mood, a mood that at its height
leads to the giving up of one's own wishes, to willing submission
to the pain of reality, because one feels oneself carried by a
power that is greater than the one over which one has command oneself. There is in
all pronounced religious feeling a tragic element, which rests on the fact
that a higher feeling of happiness has displaced a lower. The religious
feeling contains or presupposes so many oppositions that
it is one of the richest human feelings. The world is richer
for there being a devil when there is a St. Michael who can
set his foot on his neck. This distinguishing mark naturally appears
most clearly in the extreme forms of religion, forms that
to the ordinary, everyday conception stand as abnormal; but in
all forms of religious consciousness it will occur to a certain degree.


This description indicates the general frame within
which James moves in his investigation. For a closer
characterization of his work I shall first show how the
three points of view that --- as mentioned in the Introduction --- must
come forward in every comprehensive discussion of the religious
problem also appear in James, although he has mainly
% --- p. 128 ---
\setcounter{footnote}{0}%
\opage{128}kept to the purely psychological mode of consideration. Next
I shall show how he conceives the nature of religious experience,
what types of religion he finds, and
to what fundamental idea religion, in his conception, can be
traced back.

\subsubsection*{\textit{a. The Three Points of View.}}

James designates his standpoint, both in a general philosophical
respect and in a philosophy-of-religion respect, as the standpoint of experience.
He wants to be a thoroughgoing empiricist. All intellectual
operations --- whether they consist in construction,
comparison or in criticism --- presuppose immediate experience.
He allies himself with the old English school, the one that began
with John Locke, and which always demanded a definite accounting for
the experiences on which we build our concepts. This school
ought in his opinion far sooner than Kant to have the credit for
the critical method in philosophy, the only method that makes
philosophy a study that can be pursued by serious human beings.
But it is a matter not only of the experiences we have had,
but also of those we must expect or consistently bring about ourselves,
if our thoughts are to be valid. The value of our thoughts rests
on the conclusions we can draw from them. If all the propositions
we set up were practically indifferent, i.e. if
nothing at all followed from them, could we then distinguish between true and
false thoughts? In vain has dogmatism sought for criteria
of truth that could dispense with appealing to the future. In
opposition to dogmatism stands here pragmatism, which tests
every proposition by pursuing its consequences. This
conception (in which James allies himself with the American thinker
CHARLES PEIRCE, from whom the name pragmatism
also derives) recalls the biological-economic theory of knowledge
that has been developed in Germany by Richard Avenarius and
Ernst Mach.

James's empirical standpoint leads him to skepticism toward
speculative and theological ways of viewing. He asserts in particular multiplicity
in existence and is critically disposed toward monism
% --- p. 129 ---
\setcounter{footnote}{0}%
\opage{129}in its various forms. Pluralism, the doctrine of multiplicity,
will, he believes, come more and more into honor --- both in science
and in religion. From experience one can reach nothing universal,
no absolute and all-embracing unity. And different
readings might hold for different regions of existence.
Religiously regarded, it is here especially the experiences of evil that become of significance:
has not the Book of Job once and for all demonstrated the impossibility
of solving the problem to which these experiences lead?
Here the monistic hypothesis will always present great difficulties,
and would it not then be the simplest way out to assume
that the world from the first has included elements that were not
in harmony with the whole? Ordinary people have
always been more or less polytheistic, and it is also
clear that from no experience can one derive the infinity of God.
What we can experience can at most be this, that
we stand in inner connection with something that is greater than we
ourselves; but that this is an infinite cannot be demonstrated.

What thinking can accomplish in the religious domain
can only be this: to order, determine and interpret the facts,
not to produce them. Philosophy is here always secondary.
Yet it can for that reason well be of use. For the religious
consciousness always expresses its experiences in forms
that derive from the intellectual tradition under whose influence
it stands. Therefore only critical comparison can find
what is really immediate and essential and separate it from
what is local and accidental. At the same time it is the task of philosophy to
exclude assumptions that are inconsistent or decidedly in conflict
with scientific experience. What survives this
purgatory can then be regarded as hypotheses whose justification and
significance can become the subject of closer discussion, once
they have been reduced to their simplest forms.

So much for the consideration from the standpoint of the theory of knowledge. However energetically
James asserts it, it is nevertheless the psychological mode of consideration
that he is most interested in, and which gives his
work its greatest significance.

The psychology of religion makes important contributions to general psychology. There are many psychological phenomena that can be studied
% --- p. 130 ---
\setcounter{footnote}{0}%
\opage{130}particularly well in the domain of religious life. James
emphasizes, for instance, that the religious phenomena demonstrate how
small a part of our spiritual life can be explained clearly and
distinctly. Consciousness here fades through many degrees
into the unconscious, or, as James prefers to say, into the
subconscious. Immediate impulses and presentiments stir,
which often, without our noticing it ourselves, give us the first
premises on which clearly conscious thoughts are grounded. Conscious
arguments often move only on the surface of our being,
and an involuntary and immediate conviction is the deepest thing in us
(the deep thing in us). Further, the study of religious
life teaches us to know great differences between human beings, both
with respect to the kind of life of feeling and with respect to the degree
of receptivity, --- differences that we should otherwise hardly
become so aware of. We learn how development
toward one and the same end can proceed very differently, now
continuously, now by leaps. Already purely theoretically or from the standpoint of natural history
a psychological study of religious life
has its significance.

The psychological mode of consideration then in turn becomes significant
for the understanding of the religious problem. The religious
processes are seen as special forms of processes that are also
known in other psychical domains, and obey laws that are also
known elsewhere. James here points to Starbuck's investigations
of conversion in its connection with organic and psychical
development in the transitional years. Quite special weight is laid
by James on what moves beneath the ``threshold'' of consciousness
or outside its ``Margin.'' He finds it very difficult
here to draw a definite boundary and rejects atomistic
psychology, for which consciousness is composed of single elements, each
fixed and definite in itself. In any case the
subliminal or ultramarginal influences make themselves constantly felt,
and James for his part is inclined to regard
these influences as the paths by which a higher order of things
works in our inner life. These inner manifestations
constitute our nearest, our real world; so inwardly do
we not belong to the outer world. Here James finds a
% --- p. 131 ---
\setcounter{footnote}{0}%
\opage{131}point of unity between religion and psychology, although he declares every
closer determination of it to be an interpretation that is no longer
purely psychological. Interpretation is not the decisive thing either;
it is always derived, and one and the same experience can
be the object of different religious interpretations. What
most people especially lack on this point is criticism
and caution, not faith (faith). They are too inclined to
let dogmatic assumption (belief) follow headlong upon every
vivid idea, especially when this is supported by involuntary interest. ---

The third point of view, the ethical, comes forward in the appraisal
of the religious phenomena. Just as definitely as
James distinguishes between experience and interpretation, just as definitely
he distinguishes between description or explanation of what is experienced and valuation of it. Even what, from a purely medical
point of view, must be called abnormal can by its content have great
value. The value of a state does not rest on how it
came to be. When we call certain states higher than others, it
is not on account of what we think we know about their organic
conditions. If there is anything of value in S\textsuperscript{ta} Teresa's
thoughts, then it is immaterial whether she was hysterical or not.
We evaluate a force by its effect, not by its source.


In the valuation no speculative or theological system
can be taken as a basis. Here too we must be empiricists. We
judge the religious phenomena by their fruits, and
that is what human beings have really always done. A worship
ceases when it exerts no effect on the mind, and when
in its whole character it comes into conflict with something whose value
we have experienced so thoroughly that we cannot deny it. The
gods to which human beings hold fast are those they can use,
and whose commandments confirm the demands they make on themselves and
on others. We always use human standards (human
standards). We ask whether the religious life proves itself
as an ideal form of human life. It is of course
a standard that varies with the times, but another standard
we simply do not have. Only the testing must be repeated again and
% --- p. 132 ---
\setcounter{footnote}{0}%
\opage{132}again, since under new conditions it might lead
to new results. And the valuation will often turn out differently
for different individuals. When psychological investigation demonstrates
that there are great differences between human beings in their
innermost nature, one must also assume that they need different spiritual nourishment, so that they come to stand differently
in religious respect.

Toward religion on the whole James stands decidedly sympathetic.
The best fruits of religious experience are in his conviction
the best things history has to show. He feels
here the inner life stirring with seriousness and energy, working with
an inwardness and a concentration, in opposition to all that
can inhibit and divide. It is as if one were raised up into
a higher and purer atmosphere. --- James refers on this point
to Sainte Beuve's valuation of the religious personalities
he has studied in his famous work on Port-Royal. ---

We thus find in James the three main points of view in
philosophical discussion of religious questions. He announces that
in a later work he will counter the attempts that even in our
day have been made at a more speculative treatment of the philosophy of religion,
--- attempts that, strangely enough, occur more frequently
in England and America than on the Continent. For the time being his
work stands as a significant contribution in favor of the view that
it is precisely those three points of view that are to be applied in the philosophy of religion.
--- Unfortunately, in my characterization of James's
standpoint I cannot let the openness, freshness and linguistic
art he commands come into their full right.

\subsubsection*{\textit{b. The Nature of Religious Experience.}}

Religious feeling is according to James not any special or
elementary feeling. Every feeling can under certain conditions
acquire a religious character, namely when it co-operates in the
reckoning of the experiences of human life that James calls ``the total
reaction toward life.'' A somewhat closer determination
of what religion is is given by the statement that religion is one of
the ways by which unity can come into our life. With
% --- p. 133 ---
\setcounter{footnote}{0}%
\opage{133}unity comes a feeling of peace and happiness that surpasses all
that can otherwise be experienced of that kind. \textit{This} unity can be reached
by sudden breakthrough or through a long development;
here the individual differences come to light. But
the main point is that the individual, who before has felt divided and inhibited,
low and unhappy, now has the feeling of harmony, freedom and
elevation. Thoughts that before stood as ideal possibilities out on
the periphery of the soul have now become the central possession of the soul. It
is as if new energy had streamed into the individual from sources
that lie outside its consciousness, and which it would not be able to
release directly by conscious effort. This subordinate
significance of the conscious and voluntary (the hopeless inferiority
of voluntary to instinctive action) in relation to the unconscious
and involuntary is a characteristic trait in all religiosity.
Morbid forms arise when intellectual power and
interest are lacking to shape and interpret such experiences.

Religious experience has a mystical character, partly because
of ineffability, since each person here can understand the others only indirectly,
partly because of immediacy, since discursive
thinking plays no role, partly because of the short duration of the
states and the difficulty of recalling
them, partly finally because of the passive relation in which the individual
stands to these experiences. The non-religious life of consciousness
knows similar states. The religious character arises,
as already remarked, from the fact that man's relation to
life as a whole is settled in such states; thereby they acquire
a unity and energy that is otherwise lacking. What is specifically called
mysticism is only a particular and often extreme form of what
appears in all religious experience.

\subsubsection*{\textit{c. The Most Important Types.}}

James declares that the most important result he has reached
in his investigations of religious consciousness is the
great differences that it presents. A significant part of
his work therefore naturally consists in tracing these differences
back to certain \textit{types}.

% --- p. 134 ---
\setcounter{footnote}{0}%
\opage{134}There is in particular one opposition that he emphasizes very strongly,
namely the one between the healthy and the sick mind, or ---
to use an expression he has from Francis Newman --- between
the once born and the twice born (the once-born and the
twice-born). The sick souls need a change, a crisis, a
conversion in order to reach unity and peace, whereas there are healthy
souls who by immediate unfolding and surrender reach
the goal. The religion of the healthy James designates as naturalism,
that of the sick as salvationism. But within each of these groups
very different personalities and tendencies can in turn
appear.

The religiosity of the healthy souls is an enthusiasm that holds fast to
the bright side of existence and perhaps even declares
the dark sides to be illusions. Even if they do not lack an eye
for what inhibits and darkens, they are nevertheless assured
that it can be overcome; so great is their immediate trust in
inner energy and in the harmony that prevails in the universe.
This tendency appears more in Greek antiquity than in Christianity,
more in Catholicism than in Protestantism, more among
the Latin races than among the Germanic. In more recent times it appears
in Spinoza and Rousseau, in the liberal tendency
within Christianity, especially in Unitarianism (Emerson, Parker).
In America it stands in sharp opposition to Methodism and
the Revival movement. In the most recent time the so-called mind-cure
movement is a characteristic example. It heals both soul
and body by constantly instilling and holding fast to the thought that
pain is an illusion, and that everything in reality is bright
and glorious. Fear and worry are banished. It works through
suggestions, not least auto-suggestions. By constantly shouting
Hurrah for the universe one strengthens oneself in the belief that the universe
deserves a Hurrah.

The sick soul succumbs to pain and disharmony.
The threshold of consciousness toward that kind of experiences is once and for
all lowered, whereas in the healthy soul it is raised. The world
stands as empty and vain; fear is nourished of what
the course of life may bring; the soul feels itself sinful and corrupt.
The feeling of emptiness, the feeling of fear and the feeling of sin are the three
% --- p. 135 ---
\setcounter{footnote}{0}%
\opage{135}main forms of mental illness. In Tolstoy the feeling of emptiness prevails,
in Bunyan the feeling of fear and the feeling of sin. There is here an original
disharmony in the temperament that demands a second
birth. And even if this is reached, natures of this kind will
hardly attain the overflowing feeling of life that the healthy
can have.

There exist of course in reality a multitude of transitional forms
between the two types. But in their extreme forms
the opposition between them is so great that they hardly understand each other.
Each will regard its view of life as
the one right one. Thus Wesley maintained conversion as the only
way to peace, whereas Emerson ascribed to the nature that needed only
one birth a higher value than to the one that had to struggle
hard with the devils in its inner life. James himself thinks that the experiences of life
on which the twice born can build are deeper and
more complete than those of the first type. No prophet can bring
a decisive message to human beings if he does not say things
that can ring as realities for men like Bunyan and
Tolstoy. The most complete religions must be those that both
know and overcome the darkest elements of life. Therein rests
the superiority of Buddhism and Christianity. ---

In relation to this great opposition the other oppositions
that can be pointed out are for James of less importance.
Thus the opposition between leaping and continuous development,
--- the opposition between natures in whom the elements
that are to be overcome must first have lost their vital force,
and other natures in whom the new life can break through
before the old is wholly exhausted, --- and the opposition between
a development that demands conscious effort (even if this
does not lead directly to the goal) and a development that runs its course
throughout involuntarily. ---

Occasionally James gives, besides this general doctrine of types,
also hints of his own purely individual faith. He
has, in a state of nervous slackness that long prevailed, sought help
from the mind-cure movement and believes himself to have been cured
by its help, which does not prevent him from taking a critical stance
toward it. When he states that mysticism gives only in extreme
% --- p. 136 ---
\setcounter{footnote}{0}%
\opage{136}form what all religious experience contains, he adds that he
himself does not personally know mystical states. His faith he characterizes
by saying that although he cannot join himself to the current
Christianity or to speculative theism, he is convinced
that in those states in which man feels himself in the
most inward unity with the Highest he knows, new power comes
into the world and new points of departure are set; so far is his
faith supernaturalistic, assuming divine influences within
the domain of natural experience. He adds that he knows very
well that the prevailing opinion in academic circles runs in
the opposite direction. --- Belief in immortality plays only a subordinate
role for him (although in a special little work he has sought
to show that it cannot be refuted). He says of this: If only
``in all eternity'' provision is made for our ideals, I do not see
why we should not be willing to hand over the care
of them to other hands than ours.

\subsubsection*{\textit{d. The Fundamental Religious Idea.}}

Although religion according to James has its ground in the life of feeling
and its experiences, and although it appears in very different
forms, there is nevertheless a fundamental idea or a hypothesis on which it
always rests. It is this, that a state of
harmony, unity and peace can be reached, even if it must happen through struggles
and crises, and that this end is reached by an influx of energy
from a more comprehensive order of things, which psychologically regarded
expresses itself in the unconscious nature in us. How we in particular
interpret what happens here is not the essential thing; here
purely individual supplements and elaborations (overbelief)
will always come in. The main point is that for the religious consciousness
tragedy, however often it may announce itself on the path of life,
is always only provisional and partial. Shipwreck and dissolution are not
the last word of existence. All religion thus rests ultimately on
a cosmic hypothesis. On the other hand the persistence of religion is not
dependent on any special dogmatic assumption whatever. Of course
--- if religion is made one with or wholly dependent
on an animistic conception of the world, then its time will be over,
% --- p. 137 ---
\setcounter{footnote}{0}%
\opage{137}when animism dies out; likewise if one makes it into magic or
fetishism. For James what is waged here easily becomes a dispute about words.
The decisive thing for him is that the inner experiences on which
religion essentially builds are the most concrete and
immediate experiences we have. In relation to them everything
we can experience about the outer relations of the world has an abstract
and general character. There is no reason to assume that the time
of such experiences should be over, even if they must be separated
from connection with elements with which they have hitherto been inclined
to fuse.

\subsubsection*{\textit{e. Remarks on James's Philosophy of Religion.}}

The study of James's book, which appeared a year after my own
attempt in the philosophy of religion, has been of great interest to me.
Not only is it long since I read a book that has so taught me
to look with new and fresh eyes at life and at human beings
as this book by the American philosopher has
managed to do. But the book acquired a special significance for me in that
the whole direction and way of treatment, despite all difference
on not unessential points (e.g. in what James calls
his ``supernaturalism''), is so closely akin. I have
therefore a personal reason for calling attention to
it. By its spirited form and its many interesting details
it will moreover reward study, even if one stands
further from the book's direction than I do. Only a few \textit{remarks}
will I put forward here in order to mark my standpoint toward
it.

It is in my eyes with full right that James distinguishes between
personal and institutional religion. It is admittedly a
difference that is easier to make for the higher than for the lower
religions, and it is perhaps for that reason that historians of religion sometimes
shrink from admitting it. But psychological
interest must, by the nature of the case, concern above all the
personal element in religion, whereas the objective or
``institutional'' side has significance only as symptom. And
what is the main point: the religious problem really concerns
% --- p. 138 ---
\setcounter{footnote}{0}%
\opage{138}only the subjective side. The question of the position of religion ---
now and in the future --- in human intellectual life is after all
primarily whether man's feeling and need will be able
to go on holding fast to the forms and symbols hitherto given,
and whether, if these should prove insufficient, new ones can
be found or created. He who fixes his gaze on the continuity
in spiritual development must decidedly lay the main weight
on the subjective, personal domain, however much this may be
veiled by the historical predominance of the outer forms.

Yet it seems to me that James goes too lightly over
the difficulties the problem of continuity here presents. It is
by no means self-evident that the personal element in religion
can be preserved if the historical forms fall away. As it
will be seen from my Philosophy of Religion, I stand here in such a way that I see
a \textit{possibility} of a preservation of the inwardness and concentration
in the conception of life that has had so essential a support in
the historical forms of religion, although I cannot adduce any
compelling proof of the necessity thereof, --- not even that
equivalents for it can be given. We face the possibility
that value will here be lost. History shows us a constant
reciprocal relation between the personal and the institutional. If
this reciprocal relation is to persist in the future, it is possible only
if the human spirit in the course of development retains
the ability to find or create great symbols for the
deepest experiences of life, symbols that can acquire common significance
for larger or smaller groups of human beings. It can
be hoped, but not proved. And yet it is one of the most important
questions with regard to the future of our race. One must look
at the greatest conflicts of life with a very light mind and a very unclear
eye if one does not feel the sting of this problem. --- Whether
under the new conditions one will be able to go on using
the \textit{word} religion is an indifferent question, which James treats
with the irony it deserves. ---

If James's description of religious experience is compared with
the one I have given, there appears an agreement that
shows that we have the same phenomena in view. Religious experience
is for him characterized by being man's ``total
% --- p. 139 ---
\setcounter{footnote}{0}%
\opage{139}reaction'' toward life; it is an experience of peace, unity
and energy, often won through a crisis that led to the displacement
of darkness and disharmony. What I miss here is a
definite emphasis on the fact that the disharmony that religious
consciousness in all its forms strives to dissolve is a disharmony between value and reality. By the relation between
the values that man knows and the reality that he knows
his religion is determined. What is valuable can for him be
the physical life itself in its persistence and unfolding, or it
can be ideal goods. His total feeling toward life is always determined
by whether the heaven of values is darkened or
not. It is from this point of view that I have come to my
hypothesis that the persistence of value is the fundamental religious idea
or the religious axiom. James lays the weight on the
inner feeling of unity and on the experience of an influx of inner
energy, without taking sufficient account of the fact that this
unity and heightening of energy for man means that the disharmony
between value and reality is removed. Yet I do not think
that James will disagree with me in my hypothesis. As already
mentioned above, he indeed suggests that religion, while built first and
foremost on subjective experience, rests more or
less consciously on the presupposition that the tragic in
the world can only be partial and transitory, --- that shipwreck
and dissolution cannot be the last word of existence. The
religious consciousness thus comes by its nature to presuppose
something about cosmic relations, about the innermost nature of existence.
It is on this point that religion has always had a tendency
to become metaphysics, just as conversely metaphysics,
without always being conscious of it, has proceeded from religious
presuppositions in its constructions. According to James every
attempt at a more definite and concrete formulation of those ultimate
presuppositions is ``over-belief,'' i.e. leads from the domain of thought and
experience over into that of fantasy and mysticism. James's
own ``supernaturalism'' is evidently an example of this. About
the new beginnings he assumes, he has no right to say
more than that their cause cannot for the time being be specified; that they
% --- p. 140 ---
\setcounter{footnote}{0}%
\opage{140}should have a supernatural cause is not motivated according to his
own conception. ---

There is still one essential point on which I agree
with James, namely the great weight he lays on
the individual differences with respect to religious experience
and faith. His masterly portrayal of the opposition between the
healthy and the sick souls' way of taking things is a great
contribution in favor of what I call the principle of personality.
This opposition corresponds to the opposition already mentioned above
between expansive and disharmonious natures. But James
has carried the characterization through with such force and clarity
that the question whether different personalities
can really have the ``same'' religion arises with still
greater necessity than before. ---

% --- p. 141 ---
\setcounter{footnote}{0}%
\opage{141}\essay[\textit{Parmenides} of Elea. (1901)]{PARMENIDES OF ELEA.}
\streg

A. Patin, \textit{Parmenides im Kampfe gegen} Heraklit. (Bes.
Abdruck aus dem 25. Supplementband der Jahrbücher fiir klassische
Philologie.) Leipzig 1899, Teubner. P. 491—660.

Parmenides' position within the history of Greek philosophy
has attracted particular attention in recent years. What
we have from his hand is, as is well known, fragments of a
didactic poem consisting of about 150 lines. In a compressed exposition,
in which the poetic form plays a very subordinate
role, he gives his strict, radical doctrine of unity
as the only real thing, with polemic against every assumption of real
oppositions, of becoming, motion and change. We have
here for the first time in the history of European thought
a world-view that has since reappeared repeatedly
in various forms. Plato and Spinoza are spiritual kinsmen
of Parmenides. In India a related conception appears
several centuries before Parmenides in the Vedanta doctrine, which is set forth
in the Upanishads, and at about the same time as Parmenides
an analogous view asserts itself --- more as a view of life than as a world-view ---
in Buddha's teaching. Yet in
Parmenides it is not merely a metaphysical world-view that
is cast as a didactic poem. His thinking brings out in a clear
way the sharp opposition that we now designate as the opposition
between formal (logical and mathematical) science and
real (empirical) science, and this opposition appears here
% --- p. 142 ---
\setcounter{footnote}{0}%
\opage{142}for the first time in principled form. Behind the dogmatic formulations
in Parmenides one can trace as a foundation the basic proposition
that only the propositions that can be formed as pure identities
can be the object of firm conviction. What,
in contrast to strict logical and mathematical inference, is built
on observation and experience can only be ``the opinions of mortals,''
--- can only be something that seems to be as it is,
--- shows us things only from the standpoint of ``opinion.'' Nevertheless
Parmenides also gives an account of how things
must look from this standpoint. Here qualitative
and quantitative differences and conflicting oppositions appear, and
there turns out to take place a great world process with changing
phases, an entire cosmogony. Parmenides wants to describe ``how
earth and sun and moon and the common aether and the heavenly
Milky Way and the outermost Olympus and the hot might of the stars
strove forth into becoming.''

Two main questions naturally arise with regard to
this whole doctrine: 1) How did it come into being, in what relation
does it stand to the preceding philosophical thinking? 2) What
is the relation between the two parts of the didactic poem, which seem to
stand in irreconcilable contradiction with each other?

The work before us, rich in content and well thought through,
makes a significant contribution to the discussion of these two questions,
both of which have more than purely historical and philological interest.
It seems to me that the author has entered more fully
into Parmenides' fundamental thought than earlier scholars. Thereby
he has been able to find a solution where his predecessors
had either to declare that they faced irreconcilable oppositions
or had to take refuge in forced explanations. Sometimes
his endeavor to find characteristic
allusions and connections perhaps carries him too far; but his energy
has enabled him to rediscover the Parmenidean fundamental thought in details
that have hitherto stood as accidents or absurdities.


As regards the first point, the relation to his predecessors,
the question of importance here is whether Parmenides was influenced
by Heraclitus, whose doctrine of absolute becoming and of the
% --- p. 143 ---
\setcounter{footnote}{0}%
\opage{143}constant flowing stands in such sharp opposition to Parmenides'
assertion of eternal and unchangeable being. --- According to the
older conception, which is maintained by an authority such as \textit{Eduard} Zeller
(still in the most recent editions of his works), Parmenides stands
as a consequence of the older Ionian school, which sought
the primal substance in everything. This search for unity is completed, on this
conception, in Parmenides, in that he gave the principle of unity the
most consequential and abstract form that was possible at the given stage of development.
He operates with the abstract concept of being itself
and draws a series of inferences from the presupposition that what is
must be identical with itself. He asserts what one would now call
the principle of identity and the principle of continuity in their strictest forms,
which make all difference, becoming and change come to stand as self-contradictory
and illusory. The positive, intuitive expression for the
absolute, identical and continuous being is for him space,
and existence stands before him as a thoroughly uniform
sphere, closed and resting in itself. Of the Ionians' search
for the primal substance only the abstract form of the material, the form of space,
has thus remained. There is then supposed to be here a direct,
positive development from Thales through Anaximander and Anaximenes
to Parmenides, so that Xenophanes' polemic against polytheism
becomes the transitional link. --- According to a conception first
put forward by \textit{Bernays} and later maintained by \textit{Gomperz and Diels},
the matter is more complicated, in that Parmenides' thinking is not
merely a consequence and completion of the older Ionian
philosophy, but is quite especially determined by a relation of opposition
to Heraclitus. A polemic against constant becoming is assumed
here to underlie the energetic assertion of eternal
being, and throughout the Parmenidean didactic poem one finds
a constant polemic against the Ephesian philosopher. It is sought to show
that the chronological difficulties Zeller finds in this conception
disappear on closer examination. \textit{Patin} adheres to this conception,
and he grounds it more deeply and more many-sidedly
than has hitherto been done.

Parmenides' didactic poem is brief. With brief but pithy
grounding it sets up a series of predicates that are to hold
% --- p. 144 ---
\setcounter{footnote}{0}%
\opage{144}for what truly is. It is a résumé, compressed in poetic form, of a long life's
energetic reflection. It would now, Patin emphasizes,
be difficult to conceive such a series of \textit{judgments} and postulates
unless they represent reactions
against an opposite line of thought and are thus rejections of conflicting
assertions\footnote[1]{It may serve to illustrate this to recall that there are two kinds of positive judgments: 1) an immediate and involuntary assertion, and 2) an assuring assertion, which presupposes doubt or negation, and which is expressed by special particles or by strong stress on the copula. See my treatise \textit{Det psykologiske Grundlag for logiske Domme.} p. 30.}. And all the more so when closer consideration
shows that the predicates that stand firm are essentially
of negative kind: unborn, imperishable, immovable, unchangeable,
motionless, indivisible, not unfinished, without past and without future!
Here a series of thoughts is rejected, and would Parmenides have
found these thoughts by his own brooding, or is it not more probable
that they confronted him as the assertions of others?
Even expressions like ``whole'' and ``simultaneously being throughout''
become properly intelligible only as rejections. And the predicates
that are thus rejected were asserted by no thinker
with such energy as by Heraclitus. It further appears
from a closer investigation, which before Patin Diels especially (\textit{Parmenides}
\textit{Lehrgedicht}, Berlin 1897) has undertaken, that the individual
linguistic expressions and turns in Parmenides' doctrine and polemic
point to Heraclitus. --- Besides Heraclitus, it is above all
Anaximander's doctrine of cosmic development as a differentiation
out of the uniform primal substance that is the object of combat
in the Parmenidean poem.

While Patin and Diels agree in finding constant allusions
to Heraclitus in this work, they nevertheless differ on a very
essential point. According to Diels, Parmenides stands
merely in a polemical relation to Heraclitus, in that he takes account
only of the latter's statements about constant becoming, statements that
according to Parmenides rest on the false presupposition that both
non-being and being have reality. Now, however, Heraclitus in fact has,
besides the assertion of absolute flux and
% --- p. 145 ---
\setcounter{footnote}{0}%
\opage{145}conflicting oppositions, also the idea of a constant, an
unchangeable, a unity, namely the law, thought, the primal being
(``the Fire''), which does not change, but precisely persists amid the surge of
the changes as ``that which always was, is and will be''
(Fragm. 20 ed. Bywater). He did not stop at the mere differences
and changes. What he reproached the multitude of men with
was precisely that they did not see how, despite the constant
arising of differences, an eternal harmony nevertheless always persisted
(Fragm. 45). Diels now holds that Parmenides --- in
a way that cannot be said to be unusual in philosophical
polemic --- has kept to the one side of Heraclitus' doctrine,
directed his attacks against it and overlooked the other side, the
positive doctrine of the Logos. ``Freilich'', says Diels (p. 69), ``dies ist die
Regel in der historischen Entwickelung. So einseitig polemisiert
in der Historie wie in der Philosophie eigentlich Jeder gegen
seinen Vorgänger. Der Nachfolger stellt immer die schroffste
Reaktion dar. So war es stets, und so muss es stets sein nach
dem Gesetze der historischen Spirale''. In contrast to Diels,
Patin lays great weight on what is common to Heraclitus
And Parmenides. He holds that Parmenides polemicizes only
against Heraclitus insofar as the latter assumes changes in existence.
When Heraclitus at the same time assumes a unity that precisely
persists and is asserted in and through these changes, there is here
a point of contact between the two thinkers. Parmenides protested
here only against the idea that unity should persist in and through
the changes, and maintained that it alone had reality. But
the thought of unity is common, and it made the two thinkers feel
polemically disposed toward the multitude of men, these ``two-heads,''
who see only the differences and the changes. As regards
this common thought, Heraclitus now stands decidedly as the
predecessor: ``Heraklit gehört die Erfindung, Parmenides folgt
seiner Einheitslehre, bekämpft seine Vielheit und wandelt nur
in der Begründung neue Pfade'' (p. 575). Heraclitus' doctrine of unity
was, according to Patin, a predecessor of Parmenides': ``Konnte
die Vorstellung: Das Eine ist unveränderlich und unteilbar,
obwohl eine scheinbare Bewegung den Schein einer Vielheit erzeugt,
gefasst werden, ehe gedacht und ausgesprochen war, dass
% --- p. 146 ---
\setcounter{footnote}{0}%
\opage{146}die von je geglaubte Vielheit der Dinge nicht eigentlich sei,
sondern nur als das Scheinerzeugnis eines ewig bewegten Einen?
Ist denn nicht die Leugnung der Vielheit Vorbedingung der
strengen Einheitslehre? Musste nicht der Glaube an sie durch
die Annahme einer der gewöhnlichen Vorstellung immer noch
näher stehenden ewigen Bewegung im Einen erschüttert sein,
ehe man so weit gehen konnte, auch diese Bewegung selbst
noch fiir Täuschung und Trug zu erklären? Und musste nicht
erst schon die Einheit des die Individuen trennenden Gegensatzes
behauptet sein, ehe man wagen konnte, die Annahme
von Gegensätzen selbst schon als den Anfang des Irrtums zu
bezeichnen? (p. 511 f.)'' --- One must surely admit that a
development such as the one here indicated is natural, even if it was not
necessary that the two steps of thought in question be
taken by two different thinkers. But since Heraclitus'
philosophy is there historically, and since chronology (according to what Diels
and Patin have sought to show) permits conceiving it as a historical
presupposition of the Parmenidean doctrine, Patin's conception has
probability in its favor, and Heraclitus has then influenced Parmenides,
not merely by provoking his contradiction, but also
by expressing a thought that could be developed further and was developed
further, even in extreme and paradoxical form.

The conception that Patin has thus arrived at regarding the
first of the two points emphasized above now also becomes
important for his conception of the second point. When
there existed between Heraclitus and Parmenides not merely a
negative but also a positive relation, it will also prove
easier to find an inner connection between the two parts of the Parmenidean
didactic poem, the exposition of unchangeable being
and the exposition of the becoming world full of oppositions.

As a rule it has been assumed that there existed a wholly external
relation between the two parts of the didactic poem. It was thought that
Parmenides in the second part presented a doctrine he did not himself
endorse, but that he made ``the mortals'' the concession
of placing himself for a while at their standpoint and showing how
things look from here. From his theory of unchangeable
and indivisible being neither a philosophy of nature nor a
% --- p. 147 ---
\setcounter{footnote}{0}%
\opage{147}cosmogony was possible. Such a concession to the sensory view
was of course an inconsistency from Parmenides' standpoint,
and Gomperz is surely right that Parmenides could not
have made this concession had there not been a wavering in his own mind, ---
had he not felt that the
deductive line of thought could not exhaust existence, and that
in any case the appearance of qualitative plurality and change
remained and demanded explanation. It is of course clear that the problem
of unity and multiplicity must not only become insoluble,
but lead to a glaring paradox, when one begins by
declaring unity the only thing valid and real. In the most recent time
Diels has nevertheless maintained that the second part of the didactic poem neither rests
on an accommodation nor indicates a wavering in Parmenides.
Diels also finds in this second part a constant polemic
against Heraclitus and assumes that the sketch of an empirical
conception of the world given here is only meant to show how thinkers
who cannot raise themselves above the sensory differences
and changes must conceive the matter. By ``the mortals,'' whose
opinions the second part of the poem presents, certain philosophers, especially Heraclitus, are accordingly
to be thought of, according to Diels.
``The second part of the poem,'' says Diels (p. 63), ``is nothing other
than a critical survey of the conflicting views of the preceding thinkers, a doxography that --- as with the Peripatetics ---
is merely to serve the propaedeutic ends of the school.'' Already in an
essay in a Festschrift for Zeller (1887) some years ago
Diels had expressed this conception, according to which the relation between
the two parts of the poem becomes indeed still looser and more external than on the earlier conceptions.

Patin has here, it seems to me, the merit of
having made an attempt to bring the second part of the poem into
inner positive connection with the first part. He proceeds
from the premise that so penetrating and specialized an explanation of nature as Parmenides,
judging by the fragments, has attempted here must necessarily
have sprung from a serious and positive interest
in the subject. It cannot merely have been other mortals' opinions
that he wanted to present. Among the mortals he counts
% --- p. 148 ---
\setcounter{footnote}{0}%
\opage{148}himself too, and he has thus realized that besides the
certain knowledge that moves on the ground of pure logic there is
a conception of the sensory changes and differences of things,
a conception that can be more or less correct. It is
at the very beginning of the whole poem a goddess (the goddess of truth or of
justice) who shows the thinker the solitary path
of pure logic; and it is the same goddess who at the
transition to the second part of the poem declares that she will teach him
the opinions of mortals and show him the constitution or
development of the world in so thorough a manner ``that no human
view shall be able to dispute the first place with him.'' There is thus
even in the sensory, empirical domain a difference of rank
between views, and it cannot then be mere,
accidental appearance that rules, but there must be principles that
make it possible to distinguish between the correct and the false.
And on closer consideration it now turns out that these principles
are the same as those underlying the doctrine of identity
in the first part of the poem. In the second part of the poem too
Parmenides operates with the principle of identity and the principle of continuity (to
use a modern expression), where he suggests a critique of
the various opinions. Behind the presentation of the oppositions of the sensory
world and of the process of development there lies always as
background the thought of unity as that to which everything points back. However
great problems a complete connection of the poem's two points of view
may cause, it still cannot, according to Patin, be said
that Parmenides himself gave up the first point of view when he
set himself to make use of the second. The former constantly forms the foundation
for the latter.

Even if Patin on the whole should be right in his conception,
the philosophical difficulties in the relation between the
poem's two series of thoughts of course do not disappear. They have since
that time accompanied philosophical thinking on its way in various
forms. But undeniably greater unity enters
into Parmenides' presentation and line of thought when one follows Patin's
interpretation, which bears witness to an energetic immersion in the ideas
of the old Greek.

% --- p. 149 ---
\setcounter{footnote}{0}%
\opage{149}One may expect an interesting continuation of the discussion
of the Parmenidean poem, since Diels (Parmenides Lehrgedicht
p. 4) has expressly designated his treatment published in 1897
as preliminary and is preparing a new edition that is especially
to have the philosophical connection in view, whereas in the preliminary work he directed his attention merely to
Parmenides as poet.

% --- p. 150 ---
\setcounter{footnote}{0}%
\opage{150}\essay[Holberg as a Philosopher. (1903)]{LUDVIG HOLBERG AS A PHILOSOPHER.}
\streg

Holberg is known to all as a poet, to fewer as a historian,
and to still fewer as a philosopher. And yet, if one
could ask him whether there did not lie at the foundation of his whole grand
and manifold production, which embraced both comic drama, history
and philosophy, a common interest, he would
surely answer that he was first and foremost a philosopher.

It was, he says, his purpose to moralize in all possible ways: by historical writings, by satires, by epigrams, by the underground journey and by special philosophical treatises (in the Epistles and in ``Moralske Tanker''). Everywhere it has
been his intention ``to combat common errors and
to distinguish the reality of virtues and vices from their appearance, ---
to show how the shadow is taken for the body, and how
vices are confounded with virtues'' (the Preface to ``Moralske Tanker'').
Quite especially he regards his comedies from this
standpoint: ``As for comedies, one may say that
no more powerful way of writing has been invented to moralize by, and
in which the character of virtues and vices can be given true to life.''
``It is a work that befits a philosopher.'' ``It
is in this respect that I myself have not held it unworthy of my
station and age to moralize in this manner.'' And
Molière, who has resumed ``the art of making moral
plays,'' which had lain dormant since the days of Plautus and Terence,
can ``be reckoned among the great philosophers of our times.''
% --- p. 151 ---
\setcounter{footnote}{0}%
\opage{151}(``Moralske Tanker'' 1744 p. 15—21). ``There is required of a comedy writer,
first that he be a philosopher and have studied thoroughly
what is called the Ridiculum in the human race; secondly,
that he have the talent to hackle vices through in such a way
that he diverts at the same time.'' (Just Justesens Betænkning over Comoedie).


Holberg rejoices then in his Life that he has brought it
to the point ``that one can now read in Danish of the doctrine of virtues and vices,''
and that even servant girls now (with the help of his comedies)
can philosophize about the follies of the human race.


To us it is clear that Holberg would not have been
so great a poet, historian and thinker as he was, if
there had really lain a conscious purpose of achieving
a moral effect at the foundation of his production. Purposes
come indeed very often afterwards, as attempts to justify
what has been done from involuntary need and immediate impulse.
And that Holberg had the pure joy of artist and researcher
in his works, is evident enough from his comedies.
The interest in the portrayal of character, the rich abundance of comic details
and the superior humor have in any case in the moments of production
taken power from the moralizing purpose,
if it was present. The moral at the end of the plays
is evidently enough tacked on; it does not always correspond to the content of the plays
or their basic mood. It is the sobriety after the
intoxication of production that now takes the floor, and Holberg sacrifices on
the altar of his pedantic age, since he is himself a child of his time.
And something corresponding holds, as I shall try to show, also
for his philosophical works.

There are then also a few statements by him in which
the genial faculty and the artistic interest are emphasized. Thus
when he (in ``Moralske Tanker'') says of Plautus' comedies:
``Therein shines forth everywhere that festivity which is the
soul in a comedy or Anima Theatri, which no one can
bring about except him whom Nature has formed for it, whereas
many, on the other hand, with diligence, labor and a kind of discernment
can work out such plays as Terence's.'' Or when
% --- p. 152 ---
\setcounter{footnote}{0}%
\opage{152}he (in the same book) declares: ``The aim of my writings has
not been so much to censure faults and vices as to give Specimina in certain elegant sciences.''

In Holberg's time a clear concept was lacking of what
we now call genius, and therefore one liked to present as the fruit
of conscious plan and of a sober purpose what in reality
was due to an inner, creative need that does not know what
it does, and whose work afterwards awakens wonder in the producer
himself. The modern concept of genius in this sense of
the word worked its way forward in the English aesthetics of the 18th century
(under the influence of the renewed study of Shakespeare)
and was more definitely formed by Rousseau, Kant and Goethe ---

Such an involuntary need led Holberg to write poetry ---
to research and think. And if one seeks an intermediate link or a common
foundation for his poetizing and researching, it is far more likely
to be found in his psychological interest than in a moral intention.
The psychological interest connects the poet, the historical researcher
and the thinker. The poet can work only through the image of the single thing,
the single portrayal, without being able to decide, or needing
to decide, whether it is an example or an exception
that he presents. It is the joy and interest in the image itself,
in the plastically formed fullness, on which the interest of the work of art
rests. The historian shows human characters in their
whole historical connection, as effects and as causes,
and he is able to supply material for comparison. The thinker's
task is then to set up types and find general laws
or problems. In all these directions Holberg's psychological
interest has been at work. There appears, as Georg Brandes in
his book on Holberg (p. 135) has pointed out, a certain opposition
between his psychological presentations in the comedies and
his psychological observations and reflections in the philosophical
works. As a philosopher he is interested especially in psychological
riddles and in composite and seemingly paradoxical
characters. He mentions (in ``Moralske Tanker'') the miser,
the poor spendthrift, the simultaneous increase of piety and ungodliness,
the inclination to do what is forbidden. His researcher's interest
is set in motion to find the solution of these

% --- p. 153 ---
\setcounter{footnote}{0}%
\opage{153}riddles. But this interest in the psychologically riddling
is not noticed in his dramatic works; the characters in his
comedies are of one piece. Composite and complicated characters
Holberg has not been able to form in his art, although
he had an eye for them in life, and although they interested him
as a thinker. Perhaps he has here been guided by the taste of the time:
he says expressly (Epist. 66) that ``a play without outré
characters has no effect.'' But an ``outré'' character
is precisely a character in which a single trait dominates the
whole, without coming into opposition or reciprocal relation
to other traits.

Holberg's view of human life was not bright. He found
much emptiness and vanity around him, much shell and little
kernel. ``The heroic virtue,'' which defies harm and trouble, is
rare. And most people are masked and take on,
like comedians, strange forms for as long as the play
lasts. Sincere friendship is rarer than sincere hatred. Moreover
moral diseases are almost incurable; he who suffers from
vices does not himself believe that he is at fault in anything. The moralists themselves
are no better than the others, only more irritable: it is just as easy
to anger a saint as it is difficult to make
him good again, and Holberg would rather have a whole regiment of
soldiers on his neck than a single saint! They declaim
more than they teach; ``the rouged morality of our times consists more
of words and phrases than of reality.''

Whether it has always been so, Holberg will not decide.
He will not take sides in the dispute whether the world goes forward or
backward. If he were taken as judge, he would, after
having heard both parties, dismiss the case and rule that the costs
be set off on both sides.

What can be done, and what Holberg designates as his
proper task, is to expose self-deception and show how often
one takes the shadow for the body. This he designates as his
``chief aim.'' Again and again the expression about taking
the shadow for the body returns in him. He has --- as Hamlet
demands of the art of acting --- both in his art and in his
research wanted to hold a mirror up to nature. In work on
% --- p. 154 ---
\setcounter{footnote}{0}%
\opage{154}this task his art and his thinking, his poetry
and his prose, his genius and his reflection could meet. What
he sought in the end and in all domains was --- to use a
modern expression --- personal truth. It was a very
external expression of this that he found when he designated it as
his proper purpose ``to moralize.'' ---

From Holberg's philosophy in the widest sense, in which
he uses this word, I now go on to his thinking about
topics that in a narrower sense can be called philosophical.

Holberg was influenced especially by two foreign thinkers,
by John Locke, ``the great English philosophus Locke,'' who so
energetically took up the cause of tracing all knowledge back to experience,
and by Pierre Bayle, in whose writings he found ``a profound
philosophy, a thorough critique and an uncommon accuracy;''
and whom he especially praised for the sharpness with which he held
fast to the sting of the problems against hasty attempts at solution. ---
It is characteristic of Holberg that precisely these two men have
meant so much to him. Against Scholasticism and dogmatism,
which still ruled in the learned world, he set experience
and the great riddles. He asserts at one and the same time the sobriety
and the necessity of constantly continued search. Cautiously
he approaches the limits of science, and not even there
does he lose his composure, but maintains the right to and necessity
of unprejudiced investigation. The authors of the nearer contemporary period
he did not esteem highly. He returns several times
(both in his Life, in the ``Moralske Tanker'' and in the Epistles)
to the claim that there has been a decline in the
scientific and the literary world, a claim that Rousseau
some time later put forward anew, and which underlay
his demand for a new culture.

In going through Holberg's statements on philosophical topics
I find that they can be assigned to three groups. The first
group contains statements on more special ethical and religious
questions; the second group gives indications toward
a general world-view; the third group contains
statements that impress the necessity of constant search and
striving in the face of the highest questions of existence and thought.

% --- p. 155 ---
\setcounter{footnote}{0}%
\opage{155}1. A proposition of Bayle's to which Holberg returns with particular fondness
is that men's actions
are not determined so much by their faith, their principles and
their reason as by their temperament, or by their
``affects and passions.'' ``Most people's pursuits are governed
not by reason, but each one's temperament is the wheel that
drives it forward\dots{} Learning, instruction and good upbringing
are indeed not useless, but yet do not give man another
heart and another complexion: for it touches only, as it were,
the outer skin.'' (Ep. 424 cf. Ep. 210). If one properly
applied oneself to knowing oneself, one would --- as Holberg
thinks (``Moralske Tanker'' 1744, p. 91) --- find that what one
is inclined to take as holy zeal is due to purely
organic states. This shows itself especially when zeal has
its ebb and flow according to bodily condition and is not
the same in prosperity and in adversity, in health and in sickness.
One should therefore beware of fetching proofs of the truth of a doctrine
from the lives of its adherents; for people have lived nobly
and beautifully with unbelieving opinions, and orthodoxy is not always
connected with virtue. (Ep. 464). --- This opposition between
doctrine and life belonged to the psychological observations that especially
interested Holberg, not least because the motive for religious persecution
must fall away when one cannot infer from certain definite
religious opinions any definite practical conduct.
A more penetrating investigation of the relation between the
various elements of mental life, especially of the mutual relation
of the life of thought and the volitional life, the psychology of the time had no conditions for.
Holberg stopped at the paradox he found here, and contented himself
with drawing a beneficial practical consequence from it. And
it was of great importance that it was demonstrated that the relation
between view of life and conduct of life was more complicated than
had hitherto been assumed.

Nevertheless Holberg lays great weight on rational grounding
of morality. Moral philosophy belonged, by his own declaration,
to his favorite studies in his older days. Morality
ought, he thinks, to precede theology, and children ought to be taught
morals before they learn the catechism. He wishes --- at the risk
% --- p. 156 ---
\setcounter{footnote}{0}%
\opage{156}of being called a heathen --- to show man's duty from the light of Nature.
A proposition such as that we must not do to
others what we want them not to do to us can hardly
be said to be innate; in this Holberg grants Locke right; but
``this great commandment of Nature'' everyone will yet involuntarily acknowledge
as soon as it is explained to him. (Ep. 144). Very deep
Holberg's ethics is not, moreover. He has a sharp eye
for extremes and impresses: ``All that is too much
harms; for all virtue consists in mediocrity, and as soon as it
goes beyond those bounds, it is metamorphosed into
vice.'' (``Moralske Tanker'' p. 433). Yet on closer
investigation it appears that Holberg by his praise of the middle
way is thinking mainly that each individual and each particular
society according to its peculiar, individual nature requires a certain
definite development, and that what is too much for the individual
may be too little for the other, ``just as [in cases of illness]
one body has enough in one pill, but another will have two
or three of the same weight'' (``Moralske Tanker'' p. 429). This
also agrees entirely with his strong emphasis on the significance of temperament
and natural basis. Also in the judging
of characters and actions he applies this mode of consideration,
in that he demands that account be taken of the
conditions and hindrances that the individual has in his own nature.
He assumes that God passes a milder judgment for the same misdeed
``upon a Cholericum than upon one who is born with
a tempered complexion,'' and demands more of him who is equipped
with great natural gifts than of a thick-skinned person. (``Moralske
Tanker'' p. 363). One excuses a gouty man for not being able to dance,
but not a hypochondriac for not being able to be
merry and agreeable in company. (Mor. T. p. 408; 412). He who has no
gall cannot be praised for composure; to govern a
wild and spirited horse indeed costs greater skill than to ride a
feeble ass (Mor. T. p. 538). ``Although a misdeed is equally great in
two persons, it is yet not equally punishable,
since one has violent affects to contend with, but the other
sins of free will.'' (Ep. 467). --- What Holberg here, supported
by psychological experience, objects to the usual
% --- p. 157 ---
\setcounter{footnote}{0}%
\opage{157} demand of morals that all be shorn by one comb, that has
not yet received its full recognition and implementation in
ethics and in jurisprudence. There lies in this demand the necessity
of an individualization, which both in practice and in theory
sets us the greatest tasks.

2. Just as Holberg in his morals seeks to rest on experience,
so he strives to build his world-view
on a foundation of experience. He polemicizes constantly against scholastic
and speculative metaphysics. Cautious as he is, he does
not always speak out plainly; but it is clear enough to which side
he leans.

The astronomical discoveries of modern times have according to Holberg
given men a greater idea of God than they had before,
when they regarded the earth as a great part of the world and as
that for whose sake the other heavenly bodies existed. Now
it is clear that the earth is only a grain of sand in relation to
the whole world, and God's greatness and incomprehensibility then becomes evident,
if it was not so before. One must now admit that
there can be no poorer undertaking than to want to describe
God! (Mor. T. p. 134---137.) --- Indeed, Holberg is even evidently
attracted by the opinion that by ``the world'' in the Bible understands only
our solar system. Thereby it will, he thinks, be easier
to understand how God could have delayed so long in
creating the world. He says of this (Ep. 1): ``Many reasonable
men hold that one cannot say with perfect certainty
whether the creation of the now existing world, as Moses describes it,
is God's first building, or whether it is with the world in general
not of the same nature as with its parts
in particular, namely that from the destruction of one flows the rise of another
and generation. Some are even of the opinion that
under the formation of the world, as described in the history of creation,
is comprised only the vortex in which the earth with the other planets
revolves about the sun; for the wide world, of which we, even with
spyglasses, perhaps see only the smallest part, is greater than the
common man imagines: In this I decide nothing.''\footnote[1]{To the hypothesis that the history of creation in the first book of Moses concerns only \opage{158}our solar system, a statement obviously alludes in the first letter about his life, written about 20 years before the Epistles (dated 31 December 1727). After having mentioned how he had come over certain religious doubts, he says: ``There were indeed still some Dubia left, which Studium Physicæ coelestis, in which I had a singular taste, caused me; it seemed to me that Vortices Cartesiani \& Mundorum Pluralitas were contrary to the information that Moses gives about the history of creation; but some sound interpreters of the books of Moses took also these scruples from me, in that they rationally and happily united the philosophy of our times with the history of creation.''}.


% --- p. 158 ---
\setcounter{footnote}{0}%
\opage{158}This statement is characteristic of Holberg, in that it shows his striving to pursue the thoughts that empirical science
calls forth, to their last consequences, even though he thereby comes
into conflict with current and traditional assumptions. The world
has expanded for him, and it becomes more difficult for him to
attain to secure and conclusive ideas.

In a similar way he took his stand toward the dispute between
Bayle and Leibniz over whether the evil in the world could be reconciled
with the origin of the world from an almighty and all-good being.
[He] thinks indeed that the assumption of a Fall, first in the world of the angels
and afterwards in that of men, helps to keep the difficulties away from life.
But he is clear ``that all contention cannot
thereby be removed.'' Nor here ``does he venture''
to proceed to any decision. (Ep. 320). With a certain sympathy
he mentions (Ep. 225) an opinion put forward by a friend of his, B. S. (who is presumed
to be the Rector Bernhard Schnabel in Roskilde),
which solved Bayle's problem by assuming an original difference
between an active principle, the Godhead, and a passive principle,
matter. The Godhead's plans could be realized only imperfectly
because of the resistance of matter. This doctrine, which in
antiquity Plato and the Stoics endorsed, stands in opposition both
to ecclesiastical theology and to Manichaean dualism, which
latter according to Bayle's opinion was the necessary consequence
when one did not take reason captive in the obedience of faith. ``To
say my poor thoughts about it,'' says Holberg, ``it seems
to me that this argument could make Monsr. Bayle sweat,
if he had to do with a heathen philosopher.'' But
--- he adds --- Bayle's opponents were Christians, who from

% --- p. 159 ---
\setcounter{footnote}{0}%
\opage{159}revelation know that the proposition ``From nothing comes nothing'' does not
hold! --- The solution Holberg mentions here was the one Rousseau
and Voltaire, and in the 19th century Stuart Mill, adhered
to. Holberg has thus regarded it as tenable when one
proceeded exclusively from human reason.

3. Holberg finds in general that when one disregards
commonly accepted opinions and grounds one's judgments solely on
experience, which is the most trustworthy teacher, there are
many principles whose validity can be disputed. (Mor. T.
p. 493; 505). Philosophy becomes for him --- what it later
was for men like Carlyle and Stuart Mill --- a struggle against habit.
When he so often declares that he will not ``decide,'' it is not
merely out of regard for the prevailing opinions and authorities; it
is also because he has a greater eye for problems and difficulties
than his contemporaries, and because the very searching, ``inquiring,''
researching was for him something valuable. The foundation of the Reformation,
he says, was the Christian freedom to inquire
before one believes; but this foundation was not held fast,
and the consequence was that one got many popes instead of one.
(Mor. T. p. 335)\footnote[1]{The same complaints are found, and in similar terms, in Giordano Bruno and Thomas Hobbes.}. It is in a man's power to inquire,
but not to believe. ``I can,'' declares Holberg, ``for my own part
say that if I go astray, it has cost me
more to go astray than it has cost many others to believe right.''
(Mor. T. p. 349). In the dispute that was conducted in his time between
freethinkers and the orthodox, he is more often offended by the latter than by the
former: ``I have read most of the new Deistic writings without
any notable temptation: but I have been offended by
several answers.'' ``I have found in the Deists' writings
objections by which I have groaned, and knots that I have
not so easily been able to untie. Such things have yet not taken
my courage from me, but rather encouraged me to investigate everything
carefully, which I still do daily.'' (Ep. 2). To appeal
to the imperfection of reason and the inscrutable ways of God
serves every religion equally well and easily leads to favoring
laziness (Ep. 16—17). ``Those who say they see no difficulty
% --- p. 160 ---
\setcounter{footnote}{0}%
\opage{160}treat everything cavalierly and indicate that they have
never troubled themselves to examine. They easily become comfortable,
fat and prospering conformists,'' who ``with impudence condemn
all those who do not blindly follow their opinions, which they themselves
have never at bottom inquired into.'' (Ep. 25).

Holberg had an open eye for the fact that there are questions
that cannot be solved. By dealing with them as they
present themselves in the world of experience, e.g. by investigating
mental life as it appears in its expressions, we learn
nevertheless always something, whereas abstract speculations, e.g. about
the nature and origin of the soul, are only ``Researches with which the
mind is plagued in vain.'' (Mor. T. p. 152 f.). Holberg wants to let the
problems unfold on the ground of experience and follow them
as far as the human power of thought can; but no further. He declares that he is
no hater of the metaphysical study, unless by metaphysics one understands ``the
manifold terms and subtle distinctions that one misuses in
disputations to entangle one's opponent in.'' But
there is another meaning of the word, in which metaphysics is
``a study that befits a philosopher,'' even if one
here constantly moves in the domain of ``the uncertain things.'' (Third
Letter about his Life). One must only keep oneself free from false
certainty. The investigation has its value in itself. Holberg
expresses himself sometimes on this in such a way that one almost thinks one
hears Lessing or Stuart Mill. ``I write,'' he says (Ep. 52)
``not so much to dogmatize as to make known
certain opinions that I have fallen upon by meditation, and which
are not to be regarded except as problemata.'' A skeptic he will not
be; but with Bayle he regards the skeptical philosophy as a
beneficial counterweight to the dogmatism and
self-confidence of theologians and philosophers (Ep. 333). In the domain of the mind strife
is beneficial; peace means in the long run dullness and ignorance. ``It
is more fitting for man to quarrel and dispute than, like dumb
creatures, to live in a brutal concord. Therefore several have
not without reason compared the inconveniences that are born of
discord to storm winds, which knock down houses and
% --- p. 161 ---
\setcounter{footnote}{0}%
\opage{161}trees, but at the same time cleanse the air and prevent its rotting from
too much stillness.'' (Mor. T. p. 501). ---

In Holberg's philosophy we find as main features interest in
experience, bold outlook to larger horizons than the usual,
pronounced ethical interest and heartfelt conviction of the significance of
constant research and of the struggle of the mind. On individual
points one notices in the sober author a quiet enthusiasm
and a strong need to investigate. Above all he wants to
avoid taking the shadow for the body, and his philosophizing thus connects
with the motive of his poetic production.


In the founder of Danish literature we find traits that
have also been characteristic of Danish thinking in the
following time, everywhere that it has unfolded naturally. Perhaps
one may see in these traits an expression of the character of the people itself.
A comparison with the history of Swedish philosophy confirms
this. (See my essay ``Svensk og dansk Filosofi'' in
``Norden 1902'').

% --- p. 162 ---
\setcounter{footnote}{0}%
\opage{162}\essay[Henrik Steffens (1902)]{HENRIK STEFFENS.}
\begin{center}{\small SPEECH AT A CELEBRATION AT ELERS' COLLEGE}\\
{\small 11 NOVEMBER 1902.}\end{center}
\streg

When Henrik Steffens, as an old man, was writing the last
part of his memoirs and looked back on his fate and his work,
the question arose in him: what have you actually accomplished?

``In my quiet, solitary youth I was seized by a longing
that led me toward religion and nature in their full abundance.
A restless, indeed wild temperament drew me in the opposite
direction, and often enough all my life long I succumbed in
the violent struggle; but the quiet, inner striving did not leave
me. At a momentous moment I was carried onto a larger
stage; a great spiritual movement seized me. I was drawn
into the seething center of a philosophical struggle; but the master
whom I myself acknowledge would not call me a philosopher. I took part
in one of the most grandiose wars that millennia have seen, but came back as the brooder I was when I went out, and a soldier
I did not become. The condition of the German people compelled me to reflections
of various kinds; but a statesman I did not become. Am I
now to say: Your life was entirely without significance; all that you
wanted to accomplish was only a series of failed attempts?''

His answer is: ``My life was wholly an inner life. I plunged
without hesitation into the wild movements of the age; I
swam in various directions on the swelling waves of research and
% --- p. 163 ---
\setcounter{footnote}{0}%
\opage{163}action: but always only to return
to myself and to devote myself to inner contemplation in a larger connection.
Thence sprang what I have done. And has my work
then been so wholly without fruit? \textit{Has} \textit{it not at} very
\textit{different times set the minds of the young} in motion?\dots{} Has not
the seed that I have sown in foreign soil led to a more significant
growth than my own soil could produce?''

Steffens has here himself stated that on which his significance in
the history of intellectual life rests, and on which, quite especially, his significance
for Danish intellectual life rests. Here among us he is known
only through the momentary but far-reaching movement
that he aroused in those who were young a hundred years ago. Their
minds were so moved that they never forgot it, and so that through
their gratitude and through their work it was kept in living
memory among our people and stands as one of the most celebrated episodes
in its life. The awakening was deep and lasting, and it is perhaps
to it that we owe it that we are still a people. But he who brought
the awakening tarried here only briefly. And our history of literature
also mentions him only for the sake of this brief visit. The
eventful days when our greatest poet, with his help, discovered
his own inner wealth, we have had described clearly and vividly,
and they stand before us in a radiance that again and again attracts
us. But now that we celebrate the hundredth anniversary of this
event, we shall show our gratitude to him who
called it forth by showing that it was no accident
for himself that he came to intervene in the way
he did. The awakening he brought was bound up --- as we have already
heard him say himself --- closely with his character and
with the course of his life.

When Steffens began his lectures here on 11 November
1802, he was 29 years old. He was born in Stavanger
on 2 May 1773. His father was a Holsteiner, his mother
was Danish. Thus he belonged from the first to three countries, and
he loved them dearly all his life. Denmark, however, he always called
his fatherland. And his memories of childhood and youth
are tied to it. He went to school in Elsinore
and Roskilde and became a student in Copenhagen. It was the study of nature
% --- p. 164 ---
\setcounter{footnote}{0}%
\opage{164}that especially attracted him, but he had early immersed himself in Goethe
and Spinoza, and the new philosophy that was developing in Germany
under the influence of Romanticism soon aroused his attention.
When a scientific study trip to Norway
did not bring the yield that was expected, he suddenly went to
Germany, and after a hard struggle with poverty and
adversity he became a lecturer in natural science in Kiel. Soon afterwards
the Danish government sent him on a journey abroad
for a scientific purpose. It was on this journey that he
became wholly absorbed in the world of Romanticism. There rests a glorious
and fairy-tale radiance over the volume of the autobiography in which
Steffens describes his stay in Jena, Weimar and Freiberg.
Now he immersed himself in nature on lonely mountain walks,
now he studied geology under Werner in Freiberg, now he
philosophized with Schelling in Jena, now he was Goethe's guest
in Weimar, now he received deep, religiously attuned impressions
of Novalis' personality and poetry. Observation, thinking, fantasy and
emotional life worked together into an overflowing fullness,
whose individual elements he could hardly hold
apart. A timbre had been produced in his mind,
and when he gave what stirred within him words in inspired
speech, many different strings had to be set in
motion in the listeners. It was by this combined power that he worked,
with the power at once of the thinker, the poet, the naturalist and the prophet.
He himself recognized that it was not a barbaric country
that he came back to. The enthusiasm that the battle of 2 April of the
year before had aroused had also seized him in his loneliness
abroad and awakened his pride at being a Dane.
And he found in the circle of the Ørsteds and the Mynsters an independent
and understanding interest in what was stirring in the fatherland's
intellectual life. But to get the young generation awakened to
intellectual work there was required a concentrated energy such as
he came with, irradiated as he was by the glow from the homeland of Romanticism.

Only faintly is this concentration, radiance and enthusiasm felt
in the printed lectures. The half-German language and the
many technical terms harm the effect. We must constantly
% --- p. 165 ---
\setcounter{footnote}{0}%
\opage{165}keep in mind that Steffens belonged to those for whom the spoken word stands
far above the written, and the personality far above the work.

Steffens has himself somewhere designated longing for the Infinite
as the peculiarity of the romantic mind. Nothing
single and finite could satisfy it. Observing and
critical research stops at the particulars, compares and calculates
them. But from this no whole comes out, nothing that can wholly
fill and satisfy. A view of the whole is not a result,
but presupposes an original organ and comes as a
happy gift.

That the highest in the world of the mind stood open to everyone
whose sense was open, was above all what made an impression
on the two most gifted young men who heard Steffens' lectures,
Oehlenschläger and Grundtvig. It had to awaken their
hope and trust to hear this said by a man who himself had
looked about in the world of the mind. He bore witness that the springs of poetry and religion
were not dried up, but only checked by the prose and externality of the age. And he maintained that poetry and
religion sprang from the same inner well as thought and research,
as soon as these are not limited in a petty way. For it
is due --- so he further taught --- only to the dominion of the critical intellect
in modern times that the different powers of the mind --- poetry, religion, science, action --- stand as different
and as mutually conflicting. There have been happier
times when these powers worked intimately united. That was the time of
the great fullness, as it still appears to us in the
undivided works of antiquity. ``Life itself, fresh as it sprang from the Godhead,
had immediately the significance that we now only sense. No
trace of the dividing death which in more recent times has made sciences
and poetry, and through them religion itself, monstrous''.
Now this glorious time can be renewed --- even if for the present
the most beautiful and glorious thing that pervades our age is the sense
for that old, vanished time!

Such a view of the whole, which points beyond everything single
and isolated, stirs --- Steffens acknowledges it --- in every
real investigator. It is what gives the student of nature courage
to trace obscure connections and the historian to 
% --- p. 166 ---
\setcounter{footnote}{0}%
\opage{166}immerse himself in distant times. But it must be carried through even
further than they carry it. We must regard nature as the prehistory of the mind.
--- And here Steffens makes use of his ``Beiträge zur inneren Naturgeschichte der Erde,'' published the year before, to
show how in the whole of nature an individualizing
tendency expresses itself, in that already inorganic nature shows hints
that only in plant, animal and especially in human life come to
full development. What appears in the mental domain
lies already as an obscure disposition in unconscious nature.
Therefore we must seek in the innermost of our own mind the key to
the secrets of nature. The riddles are not solved by a mere accumulation
of facts. --- And as we regard nature as the prehistory
of the mind, so we find in the history of man the great, naturally
necessary course of fate. No single age may be made the
measure for other ages; what matters is to immerse oneself in the
inner connection through the ages, that which not only steps
forth in the great geniuses, in whom the development at certain points
concentrates, as if they were the end that was to be reached,
--- but reveals itself in the humblest, most limited
human being, who just as well as they is determined by the historical
development. ---

Steffens' philosophy is --- like the whole so-called ``philosophy of nature'' of that age
--- built on a bold analogizing, where a more
or less deep poetry for the most part speaks on behalf of thought.
What distinguishes him from the real head of this school,
him whom Steffens called his master, Schelling, is above all the weight
he lays on the individual in nature and human life, and on
the historical course of development. Both are connected in turn
with the fact that real observation of nature plays a greater role with him
than with the other speculative philosophers. On the other hand
it was this sense for the organic, individual and historical
that made him alien to exact, realistic research into nature;
this seemed to him to conceive everything as dead, all processes as
externally mechanical. When he laid so great weight on development,
it therefore also does not mean that he was an evolutionist
in the modern sense; for he did not assume that development by means of
determinate real causes worked its way forward from step to
% --- p. 167 ---
\setcounter{footnote}{0}%
\opage{167}step; it was for him the mind of nature that went from experiment to
experiment, ever striving up toward a highest form of individuality.

It is characteristic that he, whose personality stands above
his work, strove to produce a philosophy of personality by conceiving
the innermost striving of existence as a striving for individuality.
In his memoirs he has an utterance of interest in this respect:
``It was the intuition of the all-embracing life that dominated
my youth, indeed my very childhood. What an obscure
presentiment and longing sought was not a purely external relation to nature,
but an inner, spiritual unity in everything; therefore the least thing was
so dear to me; therefore the smallest blade of grass in its determinate
form was so valuable to me and stood in an inner, so to speak personal
relation to the most insignificant insect. Dimly this hovered before
me in my studies, and while I took up and shared foreign
views, I nevertheless always returned to that which
was certainly for a long time free fantasy and had more poetic than
scientific significance. What Schelling became for me is well known;
indeed, the drive that led me to Schelling, and which gave
my life its significance, is a main theme of my recollections of life.
But for me what I call philosophy of nature is nothing other
than the conviction that a consistent
organic conception will develop, which in every object of research recognizes
a peculiar something that develops in its own way,
--- a conception which only through such recognition is able
to grasp the significance of the particular for the whole.'' ---

It did not fall to Steffens' lot to develop these thoughts
further on Danish soil. He went back to Germany and became
professor of philosophy and natural science (especially mineralogy) in
Halle. Here his association with Schleiermacher was fruitful
for him and contributed to developing the individualistic element
in his philosophy. The introduction to ``Grundziige der Naturwissenschaft''
(1806) ends with an enthusiastic eulogy of Schleiermacher,
and this work too is marked by his influence.

In his new fatherland he lived himself in deeply. He took
part enthusiastically in its new life of thought and art; he suffered with it
in the days of defeats and humiliations. And it fell to his
lot --- through the wonderful faculty he had of moving the
% --- p. 168 ---
\setcounter{footnote}{0}%
\opage{168}minds of the young --- to speak the first word that called Germany's
youth to the struggle for freedom. He was professor in Breslau when
the German armies had to take part in Napoleon's campaign in Russia ---
and he felt the degradation in this. Now came the rumors of the catastrophe
in Russia; Napoleon had fled through Silesia
a Prussian army corps had on its own initiative joined
Russia, and the Prussian court now came to Breslau, where they
felt more secure than in Berlin, which had a French garrison. The king
was long wavering, but they had persuaded him
to sign a decree that called the youth to the colors---
The evening before it was to be published, Steffens was together with
some young officers of Scharnhorst's school. They feared
that the decree on account of its indefinite terms would not
have the right effect. Then Steffens in the course of the night resolved
to intervene. He announced in his
morning lecture that later in the day he would give a
speech about the royal proclamation that would be issued. The rumor
of this spread, and a great crowd filled the auditorium, the stairs
and the street outside. Steffens declares that he cannot give an
account of what he said; it was what moved everyone's
thoughts, but now for the first time got air publicly. He ended by
calling on the young to enlist and declared that he
too himself would be a soldier. The day after he had to repeat
the speech for a still larger audience. A movement was now aroused,
which fulfilled the wish and longing of the German patriots.
Scharnhorst embraced him and cried: ``Steffens, I congratulate
you; you do not know what you have done!'' --- Steffens took part in the war as
lieutenant and adjutant.

When after the unrest of the campaign he had returned again to his
quiet work in Breslau, it was not long before the ideas of his youth
came to stand in a somewhat different light for him than
before. He did not give them up, but he saw that the road to their
realization was longer than he in the enthusiasm and exuberance of youth
had thought. It now appeared to him (he says in
a letter to Tieck) as if it were a tower of Babel that he had
been helping to build. The division of labor in the scientific,
and generally in the mental domain, made itself felt more and
% --- p. 169 ---
\setcounter{footnote}{0}%
\opage{169}more strongly; the different faculties and subjects
isolated themselves and went each their own ways, so that the unity in the world of thought
and of the mind that he had believed was about to come,
stood as a distant ideal. ``I saw only too clearly,'' he says,
``that not the enthusiasm at the beginning of the century,
but only the slow course of history would gradually bring about
what vigorous, flaming life already believes it possesses in
full maturity. But this vanished hope had transformed
itself into a religious faith that did not waver, for it did not rest
on its own strength, but on a divine strength. Thus
I could experience that everything failed between my hands,
without my becoming irresolute''. ``All development is quiet. You
do not see the grass grow, but it may well succeed for the calmly contemplating
mind to come ever nearer to the silent
course of development.''

This mild resignation is all the more remarkable
as Steffens' view of life otherwise in the later years took on a stricter
stamp and sharply emphasized the opposition of the elements in the
life of the individual and of the race. In ``Caricaturen des Heiligsten''
he finds sinfulness to consist in the isolation of the single element
from the whole to which it belongs, and such isolation
appears to him in a dubious way in modern
individualism and liberalism. But caricature is for him
also any one-sided scientific view, and it seemed to him to be man's fate now to fall under the one,
now under the other thought's dominion and thereby lose
sight of the true peculiarity, which can be maintained only
within a comprehensive whole. He took an even more decided and
sharper distance from rationalism than he had done in the
romantic days of his youth. His religious faith acquired a pronounced
confessional character, as he (in the work ``Von der falschen
Theologie und dem wahren Glauben'') joined the strict
Lutheranism and its resistance to the Union with the Reformed which
the government wanted to force through.

But his philosophy he did not essentially change. In his
``Anthropologie'' and his ``Religionsphilosophie'' are found the same
basic thoughts as in his youthful writings (only in a markedly
% --- p. 170 ---
\setcounter{footnote}{0}%
\opage{170}theological form), the same attempts to find
indications of the powers and struggles of human life on the lower stages of nature,
to read nature, so to speak, from above downward. The thoughts of Romanticism
were rather poetic than philosophical, and Steffens has
himself also stated that a poem like Dante's was constantly working
its way forward in his inner self, without ever gaining form.
Like an inner melody the view he in his youth had believed he could
form into a system was preserved up into his high old age. And when he saw
science, which he had believed he could dissolve into poetry, turn back to
the strict labor of prose, he followed this development too with interest (as
is e.g. also seen from his participation in the first Nordic
meeting of naturalists in 1840), perhaps because he --- like Hans
Christian Ørsted --- was certain that this honest prose would
be able to lead to a deeper poetry than the one that had been thought to be
won by shunning it.

His last years he lived as professor at the University
of Berlin. When one evening at a friend's he had heard a
new translation of Dante read aloud, he suddenly had a hemorrhage and
died a few days later (13 February 1845), 72 years old.

In his memorial speech over him Schelling said: ``Er starb
in seiner Jugend!'' --- A true word, for he preserved his youth undiminished
through the faith that the ideals of youth would work their
way forward along other roads than he had imagined from the first.
One says that he whom the gods love dies young. If Steffens had
died right after his appearance in Copenhagen, he would perhaps
stand in a still more romantic light for us. But do
not the gods love him still more, to whom it is granted to
preserve the spirit of youth undiminished through the struggles
and disappointments of a long life? ---

We may then, after having followed him to his death, return
to his appearance in his fatherland as the young
romantic herald. What he accomplished here, he accomplished
out of his whole personality, and it was only a link of a course of
life that was borne by an ideal faith from first to last. It was, then,
no accident that he worked here as he did. And it was
also no accident that he gave \textit{us} of the best he had
% --- p. 171 ---
\setcounter{footnote}{0}%
\opage{171}to give. For the love of his fatherland and the longing for
it extended, as his recollections of life show, through his whole
life. So shall we then also anew find true what Grundtvig sang
at his death, that he has

\begin{verse}
spoken a word that is never forgotten, \\
awakened a spirit whose power is felt, \\
while we have a mother tongue!
\end{verse}

% --- p. 172 ---
\setcounter{footnote}{0}%
\opage{172}\essay[Henrik Steffens in Danish Intellectual Life (1902)]{HENRIK STEFFENS IN DANISH\\ INTELLECTUAL LIFE.}
\begin{center}{\small SPEECH AT THE REFORMATION FESTIVAL OF THE UNIVERSITY OF COPENHAGEN}\\
{\small 20 NOVEMBER 1902.}\end{center}
\streg

It is a hundred years this month since Steffens began
the series of lectures here at the University that
were to intervene so strongly in the development of Danish intellectual life.
The memory of this has already been celebrated by a circle of students
who are especially connected with the room where Steffens
held his lectures. I had the honor of speaking
on that occasion and sought to show that the sudden and awakening
influence that Steffens exercised here at home during his brief
stay was closely connected with his whole personality,
as this appears to us in the course of his life. At
Askov Folk High School this memory has also been celebrated; there is
a historical connection between Steffens' appearance and the folk high school movement.
I shall now try to point out somewhat more closely
how his appearance has had effects in the different domains of
intellectual life at home.

Henrik Steffens came home after a two-year stay in Germany,
where he had moved among the greatest minds of the age
and had been initiated into the circle of thought of the romantic school.
Research into nature and speculative philosophy, poetry and religion --- all
these different strings sounded together in his proclamation. He
% --- p. 173 ---
\setcounter{footnote}{0}%
\opage{173}wanted precisely to awaken a view of the whole, in which that which the
Enlightenment and prose of the age were accustomed to separate worked together.
Here he wanted to show an inner connection in nature, so
that the development of nature came to stand as the symbolic
prehistory of the life of the mind. The key to the understanding of nature was to be sought in
the innermost depth of the mind. At the same time he pointed to the ideal
powers that revealed themselves in history, in art and in religion.
Even if the listeners did not understand everything in his presentation,
they still felt that he was enthusiastic for the great, and
it dawned on them that there were spiritual treasures in them and around
them for which they had hitherto been blind.

Upon three men at home Steffens has especially worked, partly
immediately by these lectures, partly later.

Most known is his influence on Oehlenschläger, whom
he --- as he himself has said --- awakened to consciousness of his
own spiritual wealth. It was the golden age of Danish poetry
that was thereby inaugurated, and it was not merely the great
poet alone, it was our people itself that thereby came to consciousness
of its spiritual wealth. The relation between Steffens
and Oehlenschläger has been described so fully, partly so brilliantly, by
other literary historians that I dare presuppose it as sufficiently
known.

In another young man who heard Steffens it was probably
poetry that was first and foremost awakened --- but it was regarded
in the end as the introduction to a still more deep-going
movement. Grundtvig was not a little apprehensive about the romantic
philosophy that Steffens proclaimed. He saw in it ``a crooked
and strong counteraction of the poetic in man against the
constantly dissolving use of the intellect''. He rejoiced that Steffens
spoke ``with an astonishing reverence'' of Christ and
``scornfully of the Enlightenment that had hitherto been trumpeted abroad''. But
he was also glad that ``only few understood anything coherent
of his speech, decked out with many technical terms,''
--- for ``there was much ungodliness in his teaching''. Even later,
when Steffens in religious respect took up a decided confessional
standpoint, Grundtvig after a meeting with him
(during his visit here in 1824) expressed his sorrow that they were ``heaven-
% --- p. 174 ---
\setcounter{footnote}{0}%
\opage{174}wide apart''. But in spite of everything he held fast to the view that his own
``historical-poetical activity'' was in the ground only a continuation
of Steffens' appearance here at the beginning of the century.

When one thinks of what significance that which Grundtvig
here calls his historical-poetical activity has had and has
for wide circles in our people --- how the whole folk high school movement
is a fruit of it --- this utterance of Grundtvig
(which occurs in a letter to Ingemann of 26 November
1824) is of great interest. It is under the banner of Romanticism that the folk
high school has worked, even though with great ability it has on its
way taken up elements of another kind. The greatest folk-educating
movement in Denmark in the 19th century can
thus be traced back to Steffens' lectures as its first source.

It cannot be denied, and I wish especially to say it from
this place, that in university circles the significance of
this movement has for too long been misjudged. It is natural that men of
science are apprehensive and critical toward a popularization of
science, since thereby the single facts and doctrines are
easily torn out of the whole scientific connection,
without which they cannot be seen in their significance and in their
limitation. But life has its right, and life appropriates the thoughts of science
in its own way, from points of view and
ends that life itself has led to adopt. Only so can
the appropriation become living and fruitful. And here the ``historical-poetical'' conception
introduced with Romanticism has had its great
significance. But then it must also not be forgotten that life
in its continued development can and must bring about new points of view
and ends, about which the treasures of mind and thought
may be gathered, and in whose light they may be interpreted. New times, new
thoughts! Steffens too came after all as herald of a new time
and of new thoughts. Such cannot be given once for all.
If strife comes between the different outlooks, it is
only to the good when the strife is conducted in a manner worthy of the subject,
without insinuation and without fanaticism. It serves
to our people's honor that it is awakened to such a strife, so that
% --- p. 175 ---
\setcounter{footnote}{0}%
\opage{175}it is not merely in the world of the learned that the great questions of life are fought over.

The University has not the task of immediately bringing
the truths out to the great circles of the people. But therefore one has
no right to place the scientific institution of higher learning in opposition
to the people. Scientific research is carried on on behalf of the whole people.
However solitary the individual investigator may often feel,
and however incomprehensible what he finds may often for the time being be to
everyone else, he will yet be able to regard himself as a lookout,
who is to observe and investigate what can be learned at
his post, and to report what the others cannot see at first
hand. There is no absolute opposition between the life of thought
and the other forms of human life. ---

While it was by the romantic in his personality and
his teaching that Steffens worked on Grundtvig, his influence
on F. C. Sibbern was determined mostly by those sides of his personality
and his teaching by which he deviated from the other spokesmen
for Romanticism. And through Sibbern Steffens exercises a
significant influence on Danish philosophy in the 19th
century.

Danish thought has ever since Holberg's time had an empirical
and critical character. Psychological observation and analysis, discussion
of ethical questions and of the conditions and
limits of knowledge have lain closer to the Danish philosophers than pure
speculations. What distinguished Steffens in comparison
with the other philosophers of the romantic school was precisely
his strong emphasis on the significance of the individual in nature
and human life, his demand for a living view of nature
as foundation for philosophical thinking and his
disinclination for abstract speculations such as those Schelling
and Hegel reveled in. In this spirit he spoke already in 1802,
and in this spirit he spoke when he last had the floor publicly
at home, namely at the meeting of naturalists in 1840. And
it was also in this spirit that Sibbern worked all his life.
It is in aesthetics and in theology, with Heiberg and with
Martensen, that speculative philosophy has had its real
abode at home. The line of development that goes from Sibbern
% --- p. 176 ---
\setcounter{footnote}{0}%
\opage{176}through Poul Møller to Kierkegaard and from there on, denotes
another direction, and undoubtedly a more Danish direction.

Sibbern had become a student a month before Steffens began
his lectures. He is said to have tried to hear them,
but got nothing out of them. Later Steffens came
to stand in a radiant light for him, without it being known
what had called it forth. He longed to come into personal
relation to him. It was (as he writes to A. S.
Ørsted on 1 October 1811) only for this reason that he immediately at the beginning of his long journey abroad went to Halle, where he
then also, by his own statement, had his wish fulfilled, and
that ``in a more beautiful way than he ever dared hope.'' Sibbern
was at this time in strong ferment, partly through his striving
to win a firm standpoint in the world of thought, partly through
a hopeless passion (that which he has described in the first part of
Gabrielis' Breve). By his very personality Steffens has been
extraordinarily much to him: in a letter he speaks of the
heart-cheering and instructive in the association with Steffens. He has
through association with him not only felt himself furthered in his work of
thought, but also calmed and reconciled. His last visit to
Steffens he compares with Orestes' visit to the Goddess's holy
grove. His grief and passion are here transformed
into the bold enthusiasm that was the excellent man's
prevailing mood throughout life, even though he again and
again had to fight his way to it in spite of despondency and irascibility. The enthusiasm for Steffens he preserved constantly later;
it had, according to Poul Møller's statement, the character of an infatuation,
and he was wild with joy when Steffens later visited Denmark.

On his journey Sibbern also visited Schelling and was
much enthused by him. But his own philosophical work
reminds more of Steffens (as well as of Schleiermacher, who in a philosophical respect
stood close to Steffens). He had too great respect for the
empirical sciences to be attracted by the romantic philosophy of nature. Thus he writes in a letter
of 1824: ``It is advisable for every student of nature to keep as near to experience as possible and to be wary of every
% --- p. 177 ---
\setcounter{footnote}{0}%
\opage{177}attempt to elucidate and interpret natural things by a consideration
\textit{from above}; for to that the considerations of nature are as yet
far from being matured.'' What Sibbern means by an
interpretation ``from above'' can only be Romanticism's attempt to
read the prehistory of the life of the mind out of nature instead of
resting on the very laws of the origin of natural phenomena.
Sibbern himself did not lack interest in comprehensive
views of nature. In his philosophy such a one plays a decisive
role. But it is characteristic that he even in a certain sense
takes up a standpoint that is the opposite of Romanticism's. He
seeks precisely to illuminate mental development by analogies
that are taken from the development of nature. And as regards the development of nature itself,
he lays great weight on its proceeding
from scattered points of departure; in crystallization and in
organic growth he finds examples of this. But what
thus begins from separate points of departure is not
always so easily joined into a harmonic whole. This shows itself especially
where these points of departure are distinct individuals. There
is a struggle before the goal is reached. Sibbern has thus in this
his doctrine of sporadic development given Steffens' doctrine
of nature's individualizing tendency a realistic continuation.\footnote[1]{It might have been natural for Sibbern to refer to the economist Chr. Olufsen's treatise ``Om Menneskets Rolle i den fysiske Verden'' (1822), which Professor Joh. Steenstrup some years ago brought out of oblivion, and which contains an idea that in an interesting way stands as predecessor to Darwin's doctrine of development through the struggle for existence. Olufsen's doctrine amounts to this, that the ``vital force'' in nature, on the basis of the layers and forms whose time is past, works its way forward to new forms that fit better the changed conditions. It does not seem, according to Olufsen's presentation, to be necessary that there be organic continuity between the forms; the continuity is in the ``vital force'' itself. His doctrine is a vitalistic doctrine of development, which stands between the romantic philosopher of nature and Darwin's hypothesis.}
As late as 1843 he cites (in his work on the concept of philosophy)
Steffens' proposition that all scientific speculation
must lead to ``intuition of individuality'' (Steffens' Grundzüge
der philosophischen Naturwissenschaft, 1806, p. 36, where the
% --- p. 178 ---
\setcounter{footnote}{0}%
\opage{178}grounding is that in individuality all single elements
are comprehended by a law of the whole).

In a spirited way Sibbern applies his doctrine of the sporadic character of
development to the theory of knowledge. We human beings find ourselves at a single one of the many points of development in existence and can therefore not survey this in
its entirety. And besides, each single thinker finds himself
at a determinate point of departure within human development.
No wonder that a constant debate and struggle then takes
place between the different world-views. This struggle
is a sign that our life of the mind is still in its becoming
and constantly must seek its way and its direction. This standpoint
was bound to make him apprehensive about premature system-building. And
it was his conviction that a philosophy works best while
it is itself in ferment and in becoming; when it has begun to
settle, when its ``period of movement'' is over, its significance
will also be over. Thus Sibbern sets ``the rich
Schellingian spring period'' (whose child Steffens too was)
in opposition to Hegel's systematization.

It is especially on this point that one understands his influence on Poul Møller and Kierkegaard. For Poul Møller
any system easily became ``affectation,'' because for its carrying through
there must as a rule be used concepts that are set up for
symmetry's sake, without being won by real experience and independent work of thought. And Kierkegaard's life-task became to combat ``the System'' in all forms, in the
practical as well as the theoretical domain --- the Church's systematization of the inwardness of religious life as well
as speculative philosophy's systematization of thinking.

Steffens was, however, in spite of all that gave him a
special significance, a full-blooded Romantic in his philosophy, and
it is then natural in conclusion to point out the relation between
Romanticism's conception of nature and that which is maintained in the thinking of the present
day.

Romanticism came in philosophy not merely as a reaction
against the prose of the Enlightenment period, but at the same time as a reaction
against the dogmatic conception of modern
% --- p. 179 ---
\setcounter{footnote}{0}%
\opage{179}natural science, according to which the innermost essence of nature was supposed to be exhausted
by the mechanical points of view. In opposition to
the mechanical conception, which interprets everything in nature as motion,
the Romantics wanted to interpret the different natural processes and
stages of nature by analogies from the life of the mind, to read nature ``from above''
instead of ``from below.'' Beforehand the one attempt
may be just as justified as the other; but the Romantics'
attempt did not succeed. One cannot by means of psychical
analogies give a coherent interpretation of the phenomena of nature,
an interpretation that makes it possible to infer from the
past to the future. On the other hand the history of modern
natural science has shown that by conceiving that which
is in nature as motions in space, one can formulate
strict and clear laws that again and again stand their test in
experience. Therefore Romanticism in the latter half of the 19th century
had to give way to realism.

Yet it must be maintained that the philosophy of Romanticism proceeded
from a real problem, which holds good to this very day,
the problem of the place of the mental in nature and its connection
with nature. What it did not see was that this problem
scientifically can be treated only from two points of view, the epistemological
and the psychophysical. We know nature only
through observation and thinking, that is, through mental realities,
and the basic concepts and basic presuppositions of natural science itself
are won and formed by the activity of our consciousness
and according to its laws. This is the epistemological
point of view. The psychophysical point of view is taken when
physiology and psychology, each on its own side, investigate whether
and how mental states are connected with organic
states. From both points of view the problem is discussed
in our days, and different sciences strive each from its
side to illuminate it. But it lies far deeper than Romanticism
supposed.

The remarkable thing is now that in recent years one can trace
a new reaction against the mechanical conception of nature. In part new
phenomena that cannot without further ado be understood from it,
in part studies of the history of natural science, in part investigations in the
% --- p. 180 ---
\setcounter{footnote}{0}%
\opage{180}theory of knowledge have led to the clear insight that the basic concepts and principles of natural science are not
eternal ideas and truths, but instruments of thought, auxiliary ideas
and symbols, by which it becomes possible for us to bring connection
between the observations and to infer from past to future.
Because from certain basic assumptions we can draw conclusions
that are confirmed by observation, we have no right to assert
that those basic assumptions are absolute truths; the possibility
is not excluded that the same result could be reached
from other basic assumptions. If by realism one understands
the dogmatic assumption of the presuppositions of natural science
as absolute truths, one may say that realism has
fallen short just like Romanticism. The truth lies deeper
than either of them saw. The ultimate reality is for us that
which makes us understand our observations only by means of certain
determinate presuppositions. A conclusive concept of this reality
we are, by the nature of the case, cut off from forming. What
matters is to form the right working thoughts. They
are not the truth; but the truth is in the necessity of
forming them.

However remote science may often seem to stand from life,
there is yet an inner kinship between them --- at least
in this, that in both places it is a matter of working values. Life needs
thoughts that prove their value by leading forward. Neither
life nor science has need of thoughts that stand only
as objects of passive contemplation, of mystic contemplation.

This conception of thought is less romantic than that which
one had a hundred years ago. But it contains the germ
of a bold hope for thought, for life and for the relation
between thought and life. And it gives us a point of view from which
we --- in spite of the difference of the ages --- can understand the significance of
the thoughts with which Henrik Steffens a hundred years ago
made so deep an impression on some of the foremost of those who
were then young.

% --- p. 181 ---
\setcounter{footnote}{0}%
\opage{181}\essay[Søren Kierkegaard. A Characterization (1903)]{SØREN KIERKEGAARD.}
\begin{center}{\small A CHARACTERIZATION.}\\
(1902).\end{center}
\streg

Kierkegaard's significance for Danish intellectual life consists in
his emphasis on concentration on the personal life and
its needs, in contrast to the absorption in fantasy and speculation
that was peculiar to the age of Romanticism. From Germany
the great Romantic movement had come into Danish
intellectual life and had awakened great poetic and religious enthusiasm,
called forth powers of genius into activity, but had also
encouraged a propensity to revel in the world of fantasy and in abstract
speculations that often led away from the tasks and struggles
of the real world. It was more in the aesthetic and
the theological field than in the purely philosophical field
that Romanticism exerted its influence in Denmark. Danish
philosophy has, ever since Holberg's time, had a decidedly critical, ethical
and psychological character, which did not allow abstract speculations
to spread.

He who is to counteract a tendency must preferably have
experienced its influence himself. Søren Kierkegaard indeed
represents a reaction against Romanticism, but he himself stands on the ground of Romanticism,
and the tragic element in his position at the end
of his great life's struggle is determined to a considerable extent
by this. A great poetic faculty, a propensity to rock himself in
moods and to let memory and fantasy spread out as an
% --- p. 182 ---
\setcounter{footnote}{0}%
\opage{182}opposition to the prose of reality made him deeply receptive
to the works of Romantic poetry, and his untiring reflection
tussled with thoughts and worked in the service of problems
with a passion that knew no bounds. His
external circumstances in his youth favored his interest in literature
and thinking. He was the son of a wealthy merchant
who had come from West Jutland to Copenhagen and had worked
his way up there. From him he inherited a fortune that made
him independent, and without which he would not have been able to
publish his writings. He spent his student years in literary
pursuits without practical considerations, won, at the age of 28, the doctorate
with an interesting dissertation on ``Begrebet Ironi
med stadigt Hensyn til Sokrates'' (1841) and soon afterward began
his comprehensive career as an author.

But behind the faculty and the need for fantasizing and reflecting
there lay a power that held back and bound images and thoughts
to the earth. It was a soul in anguish, a melancholy that occupied itself
with the darkest possibilities and constantly made new struggle and
self-exertion necessary. It is a racial inheritance. The West Jutland
race has a propensity to a dark mind, as well as a capacity for wit and reflection.
It is not for nothing Hamlet's region. And Søren Kierkegaard's
father not only passed this propensity on to his sons by inheritance,
but also nourished it through their upbringing and by initiating
them into his own experiences. When as a little boy he went about
tending sheep on the Jutland heath and suffered from hunger, cold
and loneliness, he climbed one day up on a hill and cursed
the God who had given him this wretched life. This he could not
forget even in his old age. It stood before him that he
had committed the sin that cannot be forgiven, and that the stern
God's punishment would fall on his family. Under these Old Testament
ideas Søren Kierkegaard grew up; one
sees from his posthumous papers that they have constantly occupied him.

If one is to understand Søren Kierkegaard's activity as an author,
one must hold fast to both of the two forces that stirred in him:
the need to let thought and fantasy move freely --- and the need
to circle about his own self, to lose himself in the dark and

% --- p. 183 ---
\setcounter{footnote}{0}%
\opage{183}heavy, which was an inborn inheritance and was \textit{nourished} by the family's
tradition.

At one point of his life, namely
after the conclusion of the first, magnificent period of his
authorship (1843---46), the one from which the series of the
\textit{aesthetic} and philosophical \textit{writings} besides the first purely religious
writings stem, Kierkegaard himself felt a need to find a unifying point of view for
his manifold production. He was surprised that
in the motley series of writings he had composed in passionate
and breathless haste, so that the production was often about to
burst his strength, an ``overall design'' clearly emerged, in that
the whole could seem to be directed with consciousness toward a
definite goal. It was for him a great temptation to assume that
he had originally had a conscious intention with the whole
production --- but he had great misgivings about
attributing to himself such a thought. In his journals from 1849 we can
see how he wavered back and forth here. --- His
journals are the monologues he held when he walked up and down
in his rooms, now and then stopping at the desk to write down
a reflection. To others he never spoke of his innermost thoughts. In conversation he behaved essentially jokingly and ironizing,
--- when it was not a matter of comforting, for in this he was
by the testimony of several a master; as he was a master at
discovering the involuntary signs of sorrow or of a depressed mood,
he also knew, like few others, how to bring relief and encouragement.


Before I go more closely into what kept Kierkegaard
from attributing to himself a conscious intention with his authorship,
I shall first give a brief survey of it.

It was \textit{the aesthetic} writings that first attracted attention. The
most important of them are ``Enten-Eller'' and ``Stadier paa Livets Vej''.
They depict in poetic form different views of life in their
opposition to one another. By a ``stage'' Kierkegaard understands
not a period of life that by virtue of a natural law of development
is succeeded by another period of life. No, each single standpoint is
depicted as so definitely stamped and closed that the transition
to another standpoint can occur only by an incomprehensible leap,
% --- p. 184 ---
\setcounter{footnote}{0}%
\opage{184}by a choice \textit{that} cannot be grounded psychologically. Kierkegaard
describes especially three such views of life. --- For the ``aesthetic'' view of life enjoyment, the enjoyment of the moment, is everything. What
matters here is to possess such a virtuosity or
elasticity that one can seize the single moment and its pleasure
on the wing, without letting oneself be bound and ensnared and without depriving
oneself of the capacity to give oneself with the same freedom to the following moments.
It is a capricious hovering of arbitrariness over life
and the relations of life that is here the characteristic. --- For
the ``ethical'' view of life, on the contrary, what counts is continual work,
struggle and victory, faithfulness through time, so that there comes to be an
inner connection between the single moments. Marriage
here stands in opposition to momentary erotic love, vocational activity
in opposition to fleeting occupation. --- If
``the aesthete'' lives in the moment, ``the ethicist'' in faithfulness through time,
then ``the religious person'' lives in a disharmonious relation between
time and eternity. His goal is eternity, in which he
according to his true nature is at home, and he therefore lives in time
as in a medium heterogeneous to his nature, --- like a fish
on dry land. The sharper the opposition is between time and
eternity, the more pronounced is the religious standpoint. In
genuinely Christian religiosity the highest point is reached.

The most important of the \textit{philosophical writings} is what Kierkegaard
has called ``Uvidenskabelig Efterskrift''. Here is given a more
conceptual development of the relation between the different
views of life, which in the aesthetic \textit{writings} were depicted more in a
poetic-intuitive way and through typical representatives. Moreover,
Kierkegaard here criticizes sharply and mockingly speculative
theology and develops a theory of knowledge that is decidedly
in the spirit of critical philosophy. The difficulties and limits
of a completed knowledge are demonstrated. Experience presents oppositions
that can by no means, as Romantic speculation
thought, easily be brought under ``a higher unity''. And
above all: life is lived in time, in the midst of time and its tension!
All understanding is gained by looking backward, but life's tasks
point forward. We understand backwards, but we live forwards.
Life therefore demands far more faith and hope than science
% --- p. 185 ---
\setcounter{footnote}{0}%
\opage{185}can ground. Thus it is the relation to time that poses
to both the view of life and thinking the greatest problems.

\textit{The} \textit{religious writings} began --- which Kierkegaard lays very great weight on --- simultaneously with the aesthetic production.
It is a series of ``edifying discourses'', directed to
``the single individual'', rich in mood, sometimes too rich in reflection. Kierkegaard's religious standpoint emphasizes from the first
inwardness as the decisive thing; the personal feeling and
passion in the religious relation is more important than the objective
correctness of the content of faith. A heathen, he says, who
prays in inwardness to an idol, in reality prays to the
true God, whereas he who in an external and impersonal way
prays to the true God in reality prays to an idol.
One comes into relation to God only through the passionate
need that clings to him in order not to perish. God is an
emergency defense --- not an object for speculation. Toward the assailants of Christianity (especially Ludwig Feuerbach) Kierkegaard stands
in such a way that he grants them right against theology and the Church,
but not against ``the Christianity of the New Testament'', although he
thinks that they often give a more correct image of this than the
orthodox theologians. ---

The three groups of \textit{writings} cannot be held strictly apart
from one another. They are three streams that here flow simultaneously
forward, now separated, now united. In Kierkegaard's later production
(after 1846) the first two cease to flow, and there occurs
a complete concentration in the religious; it is also only now
that Kierkegaard himself experiences such an inner movement
that he dares to say of himself personally that now he stands
``at faith in the deepest sense''. But when he now cast a
glance back at the overwhelming and many-sided production
of the preceding years, it seemed to him as though it from first
to last was designed to lead to a single goal: to lead
human beings, the aestheticizing, speculating and worldly
people of the contemporary age, to self-understanding, to clarity about which ``stage''
they belonged to, and to a resolution to belong, through a forceful ``Enten-Eller'',
to the Highest, or at least to recognize clearly their
distance from it. But then it was that the misgivings and the
% --- p. 186 ---
\setcounter{footnote}{0}%
\opage{186}many considerations for and against awoke. If he was to be
fully honest, he could not attribute to himself such an original
intention.

The misgivings were of different kinds, and it is worthwhile
to dwell on them, since they cast light on Kierkegaard's innermost
nature.

Kierkegaard was enough of a psychologist to know that an aesthetic
and intellectual production on a larger scale cannot be brought about with
full consciousness and as a means to an end lying outside it.
And when he collected himself, he became clear about the involuntary
in his production, that which he did not fear to call the genial.
``Every genius is predominantly immediacy, knows
no Why''. Only afterward can the significance of the work produced involuntarily
become clear. Here fits Kierkegaard's own thought, expressed in
another context, that we understand backwards,
but live forwards; this latter holds especially of the creative
genius. And often the need to produce was so strong that it threatened
to burst him, to destroy him in soul and body. He
had to, with self-control, put a damper on his desire
to write, so that he worked only a few hours a day. ---
On the other hand he also did not dare to regard himself as a
chosen instrument of God, who without being conscious of it had
worked in the service of the religious, although he was clear that
the religious lay far deeper in his soul from the first than the
aesthetic and philosophical interests.

In addition, from the first he had conceived his authorship
as something merely temporary, an outlet for the
great need for literary production, before his practical activity
as a pastor began. And later he was several times on the point of
stopping; what led him further was partly the violent need
to produce, partly external circumstances of different kinds. Of these
external circumstances he uses the expression: ``The tension of reality
put a new string on my instrument''. This tension
was produced only by the attacks on him in the humor paper
``Korsaren''. It wounded him deeply that both his bodily
and mental personality was here caricatured. At that time one
was not yet accustomed to such things, and Kierkegaard was exceptionally
% --- p. 187 ---
\setcounter{footnote}{0}%
\opage{187}thin-skinned. It makes a tragicomic impression to see how
in his journals, thus in his deep solitude, he occupies
himself with this matter and feels a need to brand it as a
``lie'' that his one trouser leg was supposed to be shorter than the
other, as it was in the caricature pictures in ``Korsaren''.
But Kierkegaard saw in these attacks not merely a personal
insult but a national crime, because those with influence
in society tolerated such things without protest. Now he neither could
nor would stop his activity and go across as a pastor to the Jutland heath, which had been his wish;
that would be to yield to baseness! --- Another
string was put on his instrument by the official Church's
attitude toward his activity, especially toward his ever sterner
emphasis on the difference between original Christianity,
``the Christianity of the New Testament'', and the modern Christianity represented in the Church of the present. What he at that time
(around the year 1849) demanded was merely that the official Church's representatives
should acknowledge that difference, acknowledge that they were
far from the ideal, instead of giving it the appearance that there
was an unbroken connection between present-day Christianity and
primitive Christianity, not only as regards dogmatics (for here Kierkegaard
had no doubt), but as regards ethical conduct of life, especially
the relation to ``the world'' (i.e. culture). But this
concession people would not make. They stood haughtily ignoring
in the face of his activity. This made --- in connection
with the stronger personal breakthrough in him in these years
--- his religious activity take on new speed, with sterner
emphasis on the ideal demands and with ever sharper polemic
against the Church, until in his last year of life (1854---1855)
he conducted the revolutionary struggle against existing Christianity,
which at last put on his tongue the word: the Christianity of the New Testament
does not exist! ---

It is characteristic of Kierkegaard that he conceived
that his authorship might become an isolated episode in his
life. This points to a sharp opposition between his literary
gifts and his temperament. He felt that the life of fantasy
and of thought, as alluring as it was to him, would not
% --- p. 188 ---
\setcounter{footnote}{0}%
\opage{188}be able to absorb him wholly, because there was an obscure power in him which
it would not be able to overcome. In the innermost nature of his personality the need to produce and the melancholy have met each other, now
in strife, now under mutual influence. Only through the interaction
between them does one understand Kierkegaard. Without the melancholy, and
what it brought with it, he would perhaps have become a pure
poet or a pure thinker. As it was, he was hindered from devoting himself
to art or science for their own sake alone. He
was not allowed to go off along the tangent, but had to let all
images and thoughts circle about the deep center that constantly
drew them to itself. --- It is the journals from 1848—1849 that
have opened to us this understanding of Kierkegaard's innermost nature.

Melancholy has had a concentrating effect. In melancholy
the soul withdraws into itself or bends toward itself,
as the wounded man bends over the wound. And in melancholy all
thoughts lose their value --- except the thought that for the melancholic
can provide a firm hold, a sheet anchor in the storms of the inner life. --- It has had an isolating effect. Kierkegaard says that it
early became clear to him that for him there was no help or
comfort to be found with others. Living together with others was for him an impossibility
on account of the deep difference between his and others'
mood toward life. That was the reason why he broke off
his engagement. Beside the young girl's bright and light mind
he felt only truly how unhappy he was --- and he felt at the same time
that she would not be able to understand him, and that he would make
her unhappy by his heavy mind. That he did not take any
office he derives from the same reason: he was not master
of his mind --- and could therefore not enter a position
where he was dependent on others! It might perhaps have been a benefit
to him if for economic reasons he had been
forced to take on definite external duties. \textit{Having to} carry out
something at a definite time can help one over many a low mood
and can be a support in the struggle against its darkness and
paralysis. --- Now he had only the diversion of melancholy's
power that the literary work provided. The strong need to produce
was sharpened by the contrast to the unrest of melancholy, and
moreover the power of melancholy was weakened when as much strength
% --- p. 189 ---
\setcounter{footnote}{0}%
\opage{189}as possible was put into producing. ``Like that princess
in the 1001 Nights I saved my life by telling stories, i.e. by producing.
An immense melancholy, inner sufferings of a sympathetic kind,
I could master everything when I was allowed to produce.'' (Journals
1849). --- Finally, melancholy had to lead to ever
new brooding by posing a particular personal problem: is melancholy
guilt or sickness? Is it not presumptuous to deny
everything value? Kierkegaard himself compared his melancholy
with acedia, the medieval term for spiritual darkness
and listlessness, one of the seven deadly sins (for which reason also acediosi
are put by Dante in Hell --- because they have been ``little glad,
where the sun shone sweetly on meadows''). Here was material for constant
brooding and for the construction of dark possibilities. ---

Under the influence of the two forces --- the strong need to produce
and the deep melancholy --- Kierkegaard's literary
work has come into being. Sprung as it was from his innermost nature
and as an attempt to reach full clarity about himself,
it had to become of great awakening and educating significance for
others. It exercised a mission, not because it was originally
designed for it, but because it was itself the fruit of the inner life's
struggles in a human being strongly seized by life.

As a thinker Kierkegaard was led by his nature to lay
weight on experience and to distrust speculation, which
led away from the crises and tasks of personal human life.
His psychological interest kept him bound and led him along
the paths of observation, especially since he possessed in high degree the
art of description. His ethical interest led him to emphasize the significance
of work, struggle and suffering in time: we live indeed
forwards --- therefore struggle and painful tension are unavoidable
for everyone who does not thoughtlessly lose himself in fantastic speculation.
And his sharp critical faculty led him finally
to see the unscientific character of speculative theology, in which
many in that age sought the solution of the riddles of existence. --- In
this realistic and critical character of his thinking he does
not stand alone in Denmark. In the whole way in which the philosophy of Romanticism
was received in Denmark there manifests itself a tendency
to put psychological experience and critical testing in the
% --- p. 190 ---
\setcounter{footnote}{0}%
\opage{190}place of abstract and mystical speculation and an assertion of
the significance of development in time and of ethical work against Romanticism's
propensity to revel in the eternal, --- a tendency
in which we may perhaps see a peculiarity of our
national character. In Henrik Steffens, the Ørsted brothers, Frederik
Christian Sibbern and Poul Møller this whole
tendency appears more or less pronounced, however much they must
be said to belong to the age of Romanticism.

Kierkegaard was now, despite his reaction against Romanticism,
himself a pronounced Romantic. Indeed, he has --- strangely enough
in a man who asserts life and existence in their opposition
to thought, --- the remoteness from life that now recalls
the poetry of Romanticism, now the asceticism of the Middle Ages,
or rather: that only in these two forms finds life worth
living. I have already pointed out that he lacks an eye for life's
continuous development. The transition between his
``stages'' must always take place by ``leaps''. He expressly denies
that there exists an analogy between personal development and
organic development. And whereas from the first he had distinguished
between three essential stages or standpoints in life, he pushes
later more and more the second stage, the ethical view of life,
out: it becomes for him one with spiritlessness, philistinism.
``The ages from which one can learn in respect of ideality,
are: the child, the youth, the young girl, the old man --- from the
active man, from the busy housewife one does not learn in this
respect, and why not? because they are essentially occupied
with the ends of finitude''. Thus only youth
and old age are to have ideality in them, but not the striving and working
of the mature and vigorous age! One can think this only
when the concept of ideality is either Romantic or ascetic,
or both. Realism is here more idealistic than the
overstrained idealism. In the very connection of reality and
in the very work of reality the poetic and the religious
ideality must be found, if it is at all to be found
any more. And it must be found there if the race is not to dissolve;
for on hope and memory without intervening valuable activity
and reality one cannot in the long run live. Out of
% --- p. 191 ---
\setcounter{footnote}{0}%
\opage{191}reality itself the new ideals must constantly be produced, like the spark
from the flint; if one rejects the flint, the spark also
disappears. Here lies Kierkegaard's colossal limitation.

And here lies also the point where he who has been
under his strong influence must struggle to free himself.
Many a young man has, despite his admiration for him, thrown his
books on the floor and turned his back on him, when life's right and
natural connection dawned on thought, and when for knowledge
other problems presented themselves than those that were posed in
the paradoxes of the religious tradition. And Kierkegaard himself would
not and could not have followers. He would not! ``I used'',
he says, ``almost half of my strength to work against
myself; I would have no followers'' (Journals 1849).
He wanted to awaken a personal truth, which can be won only
by self-activity, not by direct adherence and imitation.
And he could not have followers: he saw to that
by his last word, which was this, that what he had fought
for his whole life, --- that it did not exist! During the violent
struggle with the Church in his last year he declared: ``The Christianity of the New Testament
does not exist!'' When in the most passionate
polemic that Danish intellectual life has known he had asserted
this proposition, he collapsed; his strength was exhausted,
and death closed his clear eyes. What then? Where would
he have stood if he had not died so early (aged 42)?
To this question no answer can be given.

He was a martyr for thoughts that must be carried further along
other paths than he saw. But by energetically and consistently
going forward along the path he saw, he has done us who came after
him the best service. We may deserve the name of his
successors and show our gratitude toward him, even if we
go forward along entirely different paths than he.
% --- p. 192 ---
\setcounter{footnote}{0}%
\opage{192}\essay[Fichte's ``Addresses to the German Nation'' and the Danish North Schleswigers (1899)]{FICHTE'S ``ADDRESSES TO THE GERMAN NATION''\\ {\small AND}\\ THE NATIONAL STRUGGLE OF THE NORTH SCHLESWIGERS.}
\begin{center}(1898).\end{center}
\streg

\section*{I.}
\streg

A people comes into being when common fate and common work
produce community in disposition and way of thinking. That one belongs
to the same people is known by one's sharing danger and security,
tension and hope, indignation and admiration with one another during
fate's course through history. With this is joined the inclination
to return to the same memories and to
form kindred ideals. Switzerland's example shows that a common
language is not absolutely necessary.

Under happy conditions the state is the external, firm organization,
within which the life of the people unfolds. But history
shows examples of an opposition between people and state,
especially when a population by conquest is torn loose from the
state to which it has belonged for long times. The people must then
carry on the struggle for its persistence without the means
of a state power and often in opposition to the power of
an existing state.

At the beginning of this century there was a time in
which the now so mighty German people came to need this
distinction between people and state, and one of its most energetic and
patriotic sons consoled his countrymen with this difference.
It has for us Danes --- with a view to the struggle that our
% --- p. 193 ---
\setcounter{footnote}{0}%
\opage{193}North Schleswig countrymen wage for the preservation of their language
and their nationality --- its interest to bring forward this episode.
It is this that has brought me to read anew J. G.
Fichte's ``Reden an die deutsche Nation'' (1808).

Let me first recall the historical situation in which
these Addresses were delivered.

During the wars of the Revolution and the Empire, which for a great
part fell upon Germany, Prussia had, after a brief participation,
held back. It had calmly looked on at Austria's
defeat and the dissolution of the German Empire. But then its turn
came. In the battle of Jena (1806) the Prussian military power was
struck to the ground. One fortress after the other had to surrender.
Napoleon made his entry into Berlin. The King
and the Government fled to Königsberg, and when the enemy came also
here, farther east. By the Peace of Tilsit (1807) Prussia ceded
half of its territory and came for the remainder of the kingdom
into complete dependence on the victor. Its policy in external,
in part also in internal affairs was determined by Napoleon. He
watched suspiciously over every movement in the conquered
land that he thought could become dangerous to him. Thus the
energetic statesman von Stein had to be dismissed; indeed, he was declared
an outlaw and had to go abroad.

Under these conditions there arose among the best in Germany the
conviction that what had been lost in the external had to be won in the internal.\footnote[1]{In a memorandum by Chief President von Vincke it was expressly stated that the point was ``to reconquer internally what the state had lost in external extent''. (Cf. R. Haym: \textit{Wilhelm von Humboldt.} Berlin 1856, p. 262).} It was under these conditions that JOHANN
GOTTLIEB FICHTE in Berlin in the winter of 1807—1808 delivered his
``Addresses to the German Nation''. The French still held Berlin
occupied; Fichte's words were often drowned out by the French drums
in the street, and often the rumor went that he had been
arrested. Napoleon has, however, not regarded this ``ideologue'' as
dangerous.

% --- p. 194 ---
\setcounter{footnote}{0}%
\opage{194}

\section*{II.}

Before \textit{we} go into the content of the Addresses, we must become somewhat more closely
acquainted with the personality of the speaker.

Fichte was a Saxon peasant's son, born in the year 1762. His
parents were poor, and the boy helped them in their trade,
weaving, or tended the geese. He distinguished himself
early by his clear comprehension and his good memory, and
that became the occasion for his obtaining a further education
of his abilities. A nobleman from the neighborhood came one Sunday too
late to church, and as he wanted to know what the pastor
had said, the goose boy Johann Gottlieb was fetched and
recited the entire sermon to him. The nobleman now resolved to
provide for the instruction of the bright boy and helped him on
to become a student. After his benefactor's death Fichte had some
hard years. He suffered want and was often on the point of despairing utterly.
Only the energetic will, which was the main feature of his character,
kept him up and preserved him from suicide. He succeeded
in obtaining a tutor's position in Zürich, and this Swiss sojourn
became in several respects significant for him. He
came to know free, republican conditions, and he made the acquaintance
of the reformer of pedagogy, Pestalozzi. He read
in the following years Kant's writings, which became decisive for his
entire intellectual direction, and came also during a visit in Königsberg
into personal contact with the old master. In the year
1794 he became Professor of Philosophy in Jena and developed here a
comprehensive and significant activity. His philosophy is
not easy to present in great brevity. He himself found it extraordinarily
difficult to give a satisfactory presentation of
his philosophy; his whole following life through he worked
to get a clear expression for his thought, without any form fully
satisfying him. His philosophical standpoint is characterized
by his at one and the same time taking his point of departure in the
highest problems of science and conceiving thinking as the outcome of
personality's innermost nature. Whereas reference to the
personal moment in thinking for so many means that one
should stop thinking, it meant for Fichte precisely the necessity
of an all the greater exertion: the point was
% --- p. 195 ---
\setcounter{footnote}{0}%
\opage{195}that thinking should give as clear and consistent an expression
as possible of personality's relation to experience. Consistency
in thought and character in will were for him inseparable.
Experience had for him greater significance in practical than
in theoretical respect: the world, as experience shows it to us,
is merely the material our will is to work up, --- the ground
of our duty. --- From orthodox side he was accused of atheism,
and by an undoubted encroachment upon academic freedom of teaching
he was removed from his professorship. He lived now for some
years in Berlin and became then a Prussian professor. After the battle of Jena
he went to Königsberg, and when this city too was occupied
by the enemy, he went over to Copenhagen.

In Copenhagen Fichte lived with the famous jurist Anders
Sandøe Ørsted, who in his youth was an enthusiastic adherent
of his. It has its interest to see that Fichte came into
close contact with distinguished representatives of the
intellectual life of the Nordic countries. In Denmark both Ørsted brothers were his
great admirers and understood in an independent way to work up the
influences they received from him. The Norwegian Henrik Steffens
had been his hearer in Jena and spoke of him with enthusiasm
in his famous lectures in Copenhagen. The Swedish
philosopher Benjamin Høijer had during his stay in Dresden associated with
Fichte and shown himself an energetic thinker in their
frequent discussions. A passage in ``Addresses to the German Nation''
perhaps finds its explanation through the experiences Fichte thus
had had. It must be added that his influence on Danish intellectual
life has been very considerable. His doctrine of personality
influenced partly through the Ørsteds, partly through Steffens and
Sibbern Danish thinking in a more thorough way than
later the Hegelian philosophy. Through Sibbern and Poul Møller
there run threads of thought from him to Søren Kierkegaard, and
directly he has intervened strongly in Grundtvig's personal development,
of which I shall presently cite a testimony.

When the Peace of Tilsit was concluded, Fichte returned to Berlin
and held here in the following winter the addresses with which we are now to
occupy ourselves. At the new University of Berlin he was
very active, but his activity did not last long. He
% --- p. 196 ---
\setcounter{footnote}{0}%
\opage{196}was infected with a fever his wife had contracted by
nursing the wounded, and died at the beginning of the year 1814.

Grundtvig wrote a memorial poem about him, in which he, the
orthodox pastor, declares what awakening he owes to the philosopher
accused of atheism. Two of the most characteristic stanzas of
this poem shall be cited here:

\begin{verse}
Gottlieb Fichte! At your giant's grave \\
I stand deeply moved in spirit: \\
I must stare upon your wanderer's staff, \\
That broke so suddenly for you in your hand. \\
Tears fall too from eyelids \\
Softly down on the black coffin lid. \\
Over the grave my tongue must stir; \\
But what can, what dare it sing forth?
\end{verse}

. . . . . . . . . . . . .

\begin{verse}
Fichte! Yes, the while I wander here, \\
I cannot, and I may not, forget you, \\
Nor deny that you were dear to me, \\
That your giant voice gladdened me; \\
\textit{It} woke me from my false repose, \\
It roused me to doubt and faith, \\
Roused me to remember \\
The light of truth and the bloom of love.
\end{verse}

\section*{III.}

In his ``Addresses'' Fichte asks first and foremost whether there
still is a German people, and whether its continued persistence in its
peculiar and independent nature can be maintained and deserves
to be maintained. The external independence is lost. The victor determines
what part the Germans are to take in the great events of
history. But the external subjection ought not to bring with it
the internal. The spirit the conqueror shall not be able to cow. The struggle
with weapons is over. But another struggle remains:
for spiritual independence, for the development of thought and character:

``Let us stand guard against the comfort that
may lie in serving; for it robs even our descendants
% --- p. 197 ---
\setcounter{footnote}{0}%
\opage{197}of the hope of future liberation. If our external activity is laid
in restraining fetters, then let us raise our spirit all the more boldly
to the thought of freedom, to life in this thought, to
wishing and desiring only this one thing. Even if freedom for a time
disappears from the visible world, let us give it a place of refuge
in the innermost of our thoughts, until the new
generation grows up around us, which gets power also in the external
to present this thought\dots{} Only let not our spirit be cowed
and subjected and taken captive, like our body!''

There is namely a domain which the victor's jealousy
has not closed: the domain of spiritual life and especially of education!
There the remaining forces must then be applied. To lose
oneself in pain and sorrow over what is lost is of no use.

Above all let not the victor's superiority lead the
vanquished to untimely admiration! Such an admiration
will easily get the appearance of fear. And it will conflict with the people's
old standard of greatness!

``No, brave German countrymen, may such folly remain
far removed from our spirit, --- may such defilement remain alien
to our language, which is formed for the expression of truth!
Let foreign lands exult in astonishment over every new event,
--- let them every decade form a new standard of greatness,
create new gods for themselves and speak blasphemy in their zeal
to praise men! Let us keep our old standard
of greatness: that great is only he who with enthusiasm grasps the
ideas that always bring only good fortune to the peoples! And let us
leave to posterity the judgment on living men!''

Fichte builds the hope for a better future on the German
people's peculiarity. It is a primordial people that dwells on its own
ground from immemorial times, and which has preserved its old language
--- in contrast to the Germanic tribes that settled
on Roman ground and there developed the Neo-Latin languages.
On language Fichte lays an extraordinarily great weight in national
respect. ``What an immense influence must not the nature
of language have on a people's whole human development!
Language accompanies indeed the individual right into the most
secret depth of his mind in all his thinking and willing. It restrains
% --- p. 198 ---
\setcounter{footnote}{0}%
\opage{198}him, or it gives him wings. It binds this whole multitude of human beings
who speak it to make up a single spirit. In it
the sensuous and the spiritual meet and melt together in such a way
that it is not possible to say whether language itself is sensuous
or spiritual.'' --- But above all the German people is a
people with a sense for inwardness, originality and personality.
What the German people has put into the world of intellectual work
bears witness to this. The \textit{Reformation} was a struggle for spiritual freedom,
a breakthrough of personality's need. Precisely because the
German people took the old church system with such deep seriousness, it felt
a need for liberation from it; the peoples that came to terms
with the old system in a more superficial way could
not feel such a need to be freed from it. The German \textit{philosophy}
founded by Leibniz and Kant has --- in contrast to
other peoples' philosophy --- a special eye for life's inner conditions;
it does not conceive truth as something that can be received
merely by external influence, but as something that must be acquired
through the release of free forces. Also in the field of \textit{education}
--- especially in Pestalozzi's pedagogy --- the
same tendency manifests itself to go from within outward, to awaken love
and need for the subject and not force knowledge on
from outside. It is Pestalozzi's pedagogy on which Fichte leans in
the outline of a new national education that makes up so great a
part of the content of the Addresses. It is to develop to independence,
make cultivation a personal possession, from first to last set
the young person's self-activity in motion. ---

In the course of the characterization Fichte thus gives of his
people in order to instill the courage and desire to be themselves, he remarks
occasionally that the traits of external and internal peculiarity
that he emphasizes hold also for the Scandinavian peoples.
Undoubtedly this remark did not come about
without influence from the experiences Fichte himself had
had through acquaintance with distinguished representatives of
the Nordic intellectual life. ---

After having characterized his people Fichte impresses the
significant difference between people or fatherland and state.\footnote[1]{The eighth Address bears the title: ``What a people is in the word's higher \opage{199}sense, and what love of the fatherland is'', and arrives at the result: ``People and fatherland lie far beyond the state in the usual meaning of this word''.}

% --- p. 199 ---
\setcounter{footnote}{0}%
\opage{199}What belongs to the same people or fatherland has first and
foremost the same language:

``Those who speak the same language are already, prior to all human
art, bound by nature itself by a multitude of
invisible bonds. They understand one another and are able to
gain ever clearer mutual understanding; they belong together
and make up a natural and indivisible whole. They cannot take into themselves
any people of other descent and other language and
mingle with it, without causing themselves confusion and mightily
disturbing the healthy progress of their development. By this inner boundary,
drawn by the spiritual nature of man himself, the outer limitation of
dwellings is then in turn determined, and according to the most natural
conception of things it is by no means the human beings who
dwell within certain mountains and rivers that are therefore one people, but conversely:
human beings dwell together (sheltered by rivers and mountains,
when fortune has so arranged) because they were already beforehand one
people by virtue of a far higher natural law.''

It is, according to Fichte's religious conception, ultimately a law of
development springing from the innermost source of life that constitutes
the essence of the single people. The Eternal splits into a series of
special forms of life, because it cannot be exhausted in any single
form. This particular law of development or special form of life is
what one calls a people's national character. The individual who
lives under this law will not with clear consciousness be able to
grasp everything that lies in it; but by hidden threads the
individual hangs together, through his nationality, with the original
source of life. Therefore it conflicts with the divine
order of spiritual life that a universal monarchy should take the place
of the various nationalities. Just as Fichte earlier spoke
out against those of his countrymen who forgot the national
sorrow in admiration of the world conqueror, this superman,
so he speaks here against those who thought that the time of
nationalities was now past and found consolation over their people's downfall in
the thought of belonging to a world state. There lay for Fichte
% --- p. 200 ---
\setcounter{footnote}{0}%
\opage{200}in this thought a crudeness that laid hands on the Holy, ``an
encroachment on the highest law of the spiritual world'', according to which
the essence of humanity can only be revealed in highly manifold
nuances, partly single personalities, partly peoples.

The state means, in contrast to people and fatherland, the
social order that secures right and peace and makes trade and
traffic possible. It is only means and condition or external
apparatus for the ideal content that unfolds in the people and
lives in love of the fatherland. ---

Those were the fundamental thoughts in the words Fichte addressed to his
countrymen in the time of need.

\section*{IV.}

For him who compares Fichte's Addresses
with the situation of the Danish North Schleswigers, there appear not only striking similarities,
but also significant differences. But by and large
the Danes in Schleswig can strengthen their courage and hope
in the struggle by the ideas of the German patriot. He has indeed expressly
declared that they find application not only to his own
people, but also to peoples such as the Nordic.

The North Schleswigers are a primordial people, both as regards the land and as regards
the language. They dwell on their fathers' soil, and they speak their
fathers' language. They have a fatherland to preserve, even if by
conquest they have come to belong to a foreign state. Because a country
is conquered, its inhabitants do not thereby become --- Fichte has
taught us that --- a part of the conquering people. There is still, after
the struggle of arms is over, a struggle to be fought about language and
nationality, a struggle in which some of the highest and noblest
qualities human beings can display find application. It
is a matter of loyalty to oneself, the preservation of spiritual
peculiarity.

But conditions are more difficult in North Schleswig than they were
in Prussia in Fichte's time. He could point out that there was after all
one domain on which the conqueror had not thought of laying
hand, namely the spiritual life, especially education.
It was Napoleon's brutal contempt for the ``ideologues'' that let
% --- p. 201 ---
\setcounter{footnote}{0}%
\opage{201}him overlook the significance of spiritual forces in the world. Those in power
in Schleswig know this significance. Their jealousy extends
also to this domain. Through the school and by making the foreign
language the sole language of instruction one wants to
annihilate the nationality of the conquered population. It
betrays indeed an understanding of what spiritual forces mean,
which Napoleon did not have. But it betrays a far greater
brutality to apply compulsion in the spiritual domain than
wholly to ignore it. Even in the states of antiquity founded on conquest
the vanquished were often allowed to keep their language and
in general their spiritual life unhampered.

In compensation the North Schleswigers enjoy the reforms
that were introduced in Prussia in the time of need after the Peace of Tilsit, as well
as the political rights owed to the short period of freedom
after 1848. They have civil rights to assert
as members of the state that has conquered them, and they can
send representatives to its legislative assemblies. Thereby
they have means to set certain limits to the compulsion
and to bring their thoughts forward in the forms of public life.
The struggle thereby becomes not merely a purely spiritual struggle, but
also a struggle for right.

The spirit that leads those in power to proceed in North Schleswig
as they do is undeniably wholly different from the spirit
that runs through Fichte's Addresses. It is the spirit of reaction, bureaucracy
and militarism. It announced itself very quickly
in Prussia, after the enthusiasm that Fichte and a whole series
of men of kindred spirit had awakened had been used as a
political means in the War of Liberation. When Fichte's ``Addresses'' ten years after
his death were to appear anew, they met with obstacles of censorship:
they had now become a dangerous book. The two sides in the German character,
ideality and brutality, here come into opposition to each other.
We must not be so preoccupied with the latter that we
forget how much we owe to the former. We must on the contrary
rejoice that we can appeal to the noblest and best
in the German folk spirit against the conduct that is now practiced against
our separated countrymen. We can bring their cause before the court
% --- p. 202 ---
\setcounter{footnote}{0}%
\opage{202}of ideas, to which the best in Germany took their refuge,
when it was their time of need.

It has sometimes been thought that in culture there lay a danger
for nationalities, at least for the small nationalities. But
real culture has its noblest flower in the fine sense
for peculiarities and differences, in the understanding of the new
and the foreign. It will therefore precisely spare the great and small
nuances within the human, so that life shall not become
poorer. When there is a danger to nationalities, it is not
because there is too much culture in the world, but on the contrary because
there is too much barbarism in the world. The barbarian tramples down the foreign
because he does not understand it, and the knowledge and
experience he has about the significance of the spiritual gives him only
one more means in hand to exercise compulsion. With barbarism
no negotiation is possible; but it must be possible to come
to an understanding with a people that still in the man who held
the inspired addresses in 1807—1808 sees one of its best sons.

% --- p. 203 ---
\setcounter{footnote}{0}%
\opage{203}\essay[Harnack on the Essence of Christianity. (1901)]{HARNACK ON THE ESSENCE OF CHRISTIANITY\footnote[1]{A review in ``Nordisk Tidsskrift for Filologi'', 1901.}.}
\streg

ADOLF HARNACK, Kristendommens Væsen. Authorized
translation by Edv. Lehmann. Copenhagen 1900, V. Pio's
Bookshop. X + 219 pp.

In the interesting proposal for the reform of instruction in
Greek that Wilamowitz-Moellendorff has submitted to the Prussian
Ministry of Culture, the idea has emerged, among other things,
that pupils should be introduced step by step into Greek literature,
straight from Homer to Paul. If this idea is carried out, it will become
the task of philology to discuss the problem of primitive Christianity. But
even now this problem is a philological problem, when one
takes the concept of philology in a wider sense. Just as the study of
Plato and the appreciation of his writings are not merely the concern of philosophy,
but a strictly philological foundation cannot be dispensed with, so
it stands analogously also with the study of the primitive Christian
writings, which must lie at the basis of every discussion of
``the essence of Christianity'': it belongs not merely under theology,
but also under philology, if it is not to hover in the
air. It is such a consideration that explains why
the significant work of the distinguished historian of dogma is reviewed
in a philological journal.

It is understandable that Harnack has felt the need to give
a positive and coherent presentation of his conception of the
essence of Christianity. In his scientific
% --- p. 204 ---
\setcounter{footnote}{0}%
\opage{204}career he has had to adopt, on so many points, a critical and polemical
attitude toward the prevailing theological opinions, so that it
could take on, in the eyes of outsiders, the appearance that his standpoint
was altogether negative. Just as the political opposition is often
accused of being subversive of society, so the theological
opposition is often accused of being irreligious. It
has certainly been a great satisfaction to the critical theologian
that, before an unusually numerous circle of students of the University of Berlin,
he has been able to set forth what for him stands
as the proper essence of Christianity. One notices in his
book the warmth and enthusiasm with which he has thought and
spoken on this subject. It will be welcome to many in this country
that the book is now available in a Danish translation.

In order to discover what Christianity really is, one must,
according to Harnack, look at its whole history. Primitive Christianity is
only the first period of this history, and no single period
can express the essence of Christianity completely and definitively.
Besides what Harnack calls ``the Gospel'', the kernel of Christianity,
there are in every period, the first included, other
contributing elements (``coefficients'', Harnack calls them),
which are due to the relations of the age and the culture at a given time and place,
and which may enter into a more or less intimate connection
with the actual kernel. As regards primitive Christianity,
these coefficients, which must be separated out in order that the Gospel
may continue to work, consist chiefly of the faith in the speedy
end of history by the Messiah's return and the faith in miracles. The
central elements of Christianity, by contrast, appear clearly
in the Sermon on the Mount, especially in the Beatitudes and the Lord's Prayer. To this
Harnack refers everyone who wants to know what Christianity
really is. Its guiding thoughts are the ideas of the Kingdom of God, which is to
come, of God as Father, of the infinite
value of the human soul, of a higher righteousness and of the commandment of love
as the fulfillment of the Law. But the most important thing is that behind all
these thoughts stands a living and mighty personality, deep and
clear, powerful and mild as no other. In his speech the
most inward earnestness is joined with an open eye for the multiplicity of life.
In his parables the greatest is mirrored in the least.
% --- p. 205 ---
\setcounter{footnote}{0}%
\opage{205}All oppositions fall away or enter equally into the service of the one
thing needful. Here is no morbidity and no ecstatic
overstrain, and yet the strongest and most inward absorption
and concentration. Of his personal development we know nothing,
and we know nothing at all about him before his thirtieth year. But
precisely this is, according to Harnack, a significant trait, for from it
one must conclude that no crises and inner storms preceded
his public appearance. All that he says comes
so immediately and so naturally, and is so thoroughly that
in which he lives and breathes, that it cannot have been won
laboriously.

As regards this last point, the conclusion Harnack
draws is nevertheless not quite secure. It holds true of almost
all figures of antiquity that they stand before us in their whole, fully developed
and harmonious form, without our learning anything about
how they became what they are. In many cases this may
come from a lack of historical reports; but it
surely is connected, on the whole, with the fact that in antiquity one did not
have our psychological-historical interest in seeing the characters
in their genesis. The psychology of antiquity was plastic and intuitive,
not genetic and analytical. Just as little as one can conclude
that Socrates won without inner struggles the superior self-command
with which he appears to us, just as little
can it be justified to draw such a conclusion with regard
to Jesus of Nazareth.

Since the essence of Christianity is, according to Harnack, understood only through
its whole history, he follows it on its way from primitive Christianity
through the apostolic age, early-church Catholicism,
Roman and Greek Catholicism and Protestantism. He gives here
a series of excellent characterizations. Of these perhaps the portrayal
of the Greek church is the most successful, and it will
interest especially the many who know and love modern
Russian literature. One will understand Leo Tolstoy's position and significance
better through it. Yet there ought surely also to have been a reference
to the peculiar Russian sects, which present some
of the most remarkable religious phenomena and illuminate the influence
that racial character has on the shaping of the religious.

% --- p. 206 ---
\setcounter{footnote}{0}%
\opage{206}After the Reformation, according to Harnack, nothing has come
forth that could put the religious problem in a new
form. The difficulties that the ``Gospel'' meets are and remain the
old ones. This statement must surprise every reader, since it
is surely clear that Harnack's own standpoint, the whole standard he
applies in distinguishing between kernel and husk in Christianity,
is due precisely to factors that have come forth after the Reformation.
There is a great difference between Luther's and Harnack's
conception of the essence of Christianity, and Harnack does not conceal
it, but gives a sharp demonstration of Luther's limitation.
Luther was too quick to make new territorial churches;
he was too afraid of ``the enthusiasts'', too eager to combat
them, too unwilling to learn from them. And he did not stand at the height
of the knowledge that was nevertheless already possible in his own
age. In particular he had no real historical understanding
of primitive Christianity. Countless problems he has not known
and was therefore without ability to distinguish between kernel and husk,
the original and the foreign. --- But if this criticism of the Reformer
stands firm, then something very essential must nevertheless have happened in the religious domain
since the Reformation! The difference
between Harnack's and Luther's standpoint must be due to
``coefficients'' that have come forth after the Reformation, and by
virtue of his own method Harnack ought to have gone more closely
into the relation of the ``Gospel'' to these new cultural elements
than he has done. And when Harnack separates out from primitive Christianity
the faith in miracles and the chiliastic messianic faith as coefficients
that have nothing to do with the kernel of the Gospel, he is
in reality making this separation by virtue of the modern
cultural elements that condition his method. He stands
on the ground of modern natural science, historical criticism and psychology
--- at the very moment when he declares this foundation
to be without significance!

Here arises now a problem of great significance, which
Harnack has not brought forward in the present book. Can the
New Testament idea of the Kingdom of God be wholly separated from the expectation
of its speedy, supernatural and external realization?
And is not the primitive Christian (New Testament) ethics determined at
% --- p. 207 ---
\setcounter{footnote}{0}%
\opage{207}the most decisive points --- the relation to family life, political life,
art and science --- precisely by this expectation? If
this is so, the ``Christian'' view of life and conduct of life
that becomes possible after the distinction between kernel and husk has been drawn
in the manner Harnack indicates must be of an essentially different
character from the primitive Christian.

It is surely due to the popular character Harnack has
given his lectures that he has not entered upon this
side of the matter. In his great History of Dogma --- one of the monumental
works of modern science --- he expresses himself, by contrast, on
it with all desirable clarity. He acknowledges there that with the
apocalyptic ideas the actual New Testament
ethics also falls away. Therefore the Christian view of life must be supplemented,
and to the question whence the supplement is to come
he refers to --- antiquity and Goethe! How this supplementation
is to take place, so that a harmonious whole of life can come
out of it, about that he says nothing. It has been for me a
disappointment that the new book likewise brought nothing on this, did not even
touch the problem. It seems to me that it must be an
altogether essential side of the matter when it is a question of orienting
young devotees of science in regard to ``the essence of Christianity''.

What, in Harnack's conviction, will constantly lead
back again to the ``Gospel'', to the kernel of Christianity, he expresses
quite briefly in his closing words: ``When we with firm will
assert the forces and values that at the high points of our inner life
shine forth as our highest good, yea as our true
self, --- when we have earnestness and courage enough to let them
count as the real and to shape our life after them, --- and
when we then look at the course of the history of humanity, follow
its ascending development and, striving and serving, seek out the
community of spirits in it, --- then we shall not sink into weariness
and despondency, but we shall win certainty of God, the God whom
Jesus Christ called his Father, and who is also our Father''.
--- Yes, what is here intimated is the course of all religious experience.
Religious experience concerns precisely the highest thing of value we know,
in its relation to reality. But should the firm conviction
of the persistence of value not be able to be formed and
% --- p. 208 ---
\setcounter{footnote}{0}%
\opage{208}expressed in another way and in other ideas than in those that
remain after the critical distinction between husk and kernel
in primitive Christianity? And are not these ideas themselves rather
to be called images and symbols than genuine thoughts? --- Even if
these questions are answered with Yes, Christianity is and remains
the mightiest spiritual movement that has intervened in the
intellectual life of European humanity, and its effects will
not disappear. But its relation to the other elements of intellectual life
is still more complicated than even the beautiful and
significant discussion that the present book contains
acknowledges. After all, they are not the least among books that one closes
with a question on one's lips, especially when they contribute to
the question's being able to be put more clearly than before.

% --- p. 209 ---
\setcounter{footnote}{0}%
\opage{209}\essay[Speech for Georg Brandes. (1891)]{SPEECH FOR GEORG BRANDES.\footnote[1]{At the celebration of his twenty-fifth jubilee as an author, October 1891.}}
\streg

This is not a jubilee like so many of those that
are held in our day, at which the only remarkable thing is that the
years have passed. What has gathered us here is precisely
that our guest of honor has been able to give the years that have now gone
by since his first appearance in literature a character of their own through
his works and his influence on the intellectual life. He
stands as the greatest force in our literature and as the most distinctive
contribution we have in these years given to world literature.
For this view, all other considerations must, it seems to me, give way on
a day like this. Criticism has had the floor in regard to him, a more or
less benevolent criticism (especially a less benevolent one); now
we have to do only with what will hold its own through the purgatory of criticism
and may be assumed to remain standing as the actual value of his work.

In seeking to find a single expression to designate what
stands for me as the significance of Dr. Brandes's activity,
I come to think of the hunting horn in ``Drot og
Marsk'', which has the power that it \textit{calls men}. With such a
horn Dr. Brandes has appeared among us.

It \textit{calls} \textit{men}. That means first that it awakens them
to manhood who are not yet men. Instead of the child's narrow
% --- p. 210 ---
\setcounter{footnote}{0}%
\opage{210}horizon a wider, infinite horizon now opens. The barriers
that hindered the view fall away; it is discovered how great the
world is, how countless the possibilities that exist, how many
forces are in full activity out there. But the man differs
from the child not only by his wider view, but also
by his greater independence and by the greater clarity about
himself. In Dr. Brandes's works I find both endeavors,
both the one to widen the view and make us partakers in the
foreign development, and the one to lead us back, teach us to
find ourselves. He wishes that through absorption in the foreign
and manifold we should enrich ourselves and at the same time become clearer
about our own peculiar power and have it awakened. And if
either of the two directions has the preponderance over the other, then it is
surely the endeavor to liberate and develop the individual
personality in its peculiarity. It is this endeavor
that lies behind his great portrayal of half a century of
intellectual history, as it is mirrored in literature.
His unique ability to characterize, to find
the sources of personal life in the authors he portrays, has
a constant address to the reader, reminds us that to work
out of one's innermost nature is the highest one can attain --- however
high or however low it may in itself be. There reveals itself in his whole activity a striving to liberate the individuals'
life from everything that hinders and stunts, and a
trust in what will be able to develop when freedom in earnest
comes to rule in the intellectual domain. In this constant
striving and in this unshakable trust I see the greatest
title to our thanks. More than ever it surely holds
now that the source of life we have in freedom be not stopped up.
The excesses and the caricatures, which will not fail to appear, we shall
surely get past.

But also in another sense the horn calls men.
It has a message for many who were themselves men contemporaneously with
him, and who work in other domains, but follow his activity
with sympathy and in their way feel akin to
it. A vigorous stir in one important domain of intellectual
life propagates itself to many other domains, just as the beat of the waves
% --- p. 211 ---
\setcounter{footnote}{0}%
\opage{211}propagates itself from the lake into quiet bays and inlets. It
has also repeatedly been shown how men of very
different direction and from very different domains have been able to
give their testimony to the life and the awakening Dr.
Brandes has been able to bring into our intellectual life. The circle
gathered here testifies to the same. Even those who fight
under other banners can nevertheless acknowledge the significance it
has that he is in the field and that his signal sounds.

Finally --- when the horn calls men, it can also be
to battle that it calls, in that its tones sound like a challenge.
What awakens to life often incites to contradiction, and where
independent men meet, it cannot pass without thought
coming to stand against thought, will against will. It is one thing that the signal
sounds for advance, another whether the road along which
it bids us go is the right one, really leads forward. About that
there can be strife. Most frequently and with the greatest bitterness Dr.
Brandes has been attacked by those who wish no advance.
I believe I have shown him a great recognition in assuming
that these attacks had not affected him more than that
he had force enough left over to meet an opponent who
came from another quarter than the usual. --- I owe, by the way,
to Dr. Brandes that I have come out of an
awkward situation. A criticism I directed against some utterances
of his was, as I learned, claimed as a gain from
a certain quarter that has always had a kindly eye for him. Thereby
I was ill at ease. But when Dr. Brandes's reply came, it
was taken from the same quarter as a gain against me, in that they were
glad to be able to swat two flies (admittedly not with one blow).
Thus the awkward situation fell away. --- In reminding
of the strife of the past, it is not my meaning that the strife merely
belongs to the past. Who knows whether there may not be occasion
for new strife? It is not out of love of strife that I intimate this possibility.
But Baggesen has said with right: ``In the world of mind
there ought to be strife!'' It is more comfortable to close ranks and roll
along as part of a massed body; but thus one does not
really get forward: that our guest of honor has taught us!

I for my part have come here to say Dr.
% --- p. 212 ---
\setcounter{footnote}{0}%
\opage{212}Brandes thanks for what he has done for the intellectual life in
our fatherland. To awaken living assent and living contradiction
and to lead forward in both ways, that is the mark of
the great, the significant activity. I wish on this day,
when the thanks are brought to him from the young in so rich and so well-deserved
measure, to bring him a greeting from contemporaries, from those
who had the same teachers as he, and who as men have
followed his course, now with assent, now with contradiction,
but never with indifference. I believe I speak in the name of many,
when I say that the day the sound of his horn is no longer heard
any more, one of the most considerable and noblest forces in Danish
intellectual life in our century will have ceased to work!

% --- p. 213 ---
\setcounter{footnote}{0}%
\opage{213}\essay[Speech for Henrik Ibsen. (1898)]{SPEECH TO HENRIK IBSEN.}
\begin{center}{\small AT A CELEBRATION IN THE STUDENTS' ASSOCIATION ON THE OCCASION OF}\\
{\small HIS SEVENTIETH BIRTHDAY.}\\
{\small 2 APRIL 1898.}\end{center}
\streg

Dear Dr. Henrik Ibsen!

It is the time of celebrations and triumphs for you. But before it
lies a long time of struggle. You have struggled with the
problems of life and of the age. You have struggled to satisfy the great
demands of art and to create new forms for artistic presentation.
And you have struggled to overcome the resistance of readers and spectators.

Let me speak first of this last struggle. That there came
to be a struggle had its ground in your having made greater demands
on your readers than any other author of our time. You
do not give fully finished solutions or decisions. There ever remain
questions and riddles, and the reader must often close
the book with a sting in his soul. Your poetry is enigmatic,
as life is. Everyone must work into his personality
what you give. The interpretation or the filling-out becomes his
own work and depends on which strings are set in motion
in him. When the curtain goes down on your plays, there begins
a new drama in the spectator's mind; he is compelled himself to appear,
though it is not on any external stage. You demand
a great work of readers and spectators. Therefore people have found
you severe and cold; they have thought that you forced the reader
% --- p. 214 ---
\setcounter{footnote}{0}%
\opage{214}or spectator into the struggle with disharmonies and problems,
but held yourself outside. And people regarded you as hostile to life
and hostile to the people, because you did not give harmonious solutions and conclusions,
plain answers one could carry home from the reading
or the play. Therefore you became a sign of contradiction.

Now when so many from so many camps gather around you
to do you homage, it is not merely because people have learned to
admire and love the great art in your works. It is also
because people have gained sight that these works of yours bear witness
that you have not withdrawn from the struggle with the disharmonies of life
and with the problems of the age. People have learned to regard you
as a fellow combatant in the great strife of the age, bearing your share of
the burden of the day.

You have never raved about solidarity. You have seen in it
a danger to personality. But your works show that you have
felt yourself solidary with those who suffer and struggle. You have experienced
the darkness and the doubt, the cowedness and the bitterness. In clear images,
that will never be forgotten, you have set before our eyes a series of human figures
and human fates, which are taken from the life of the age.
The artist is not to work through propositions and conclusions.
He works through images and individual portrayals,
and he must leave it to reader and spectator to decide whether
it is rules or exceptions, models or warnings, that are
given through his presentations. Thus life itself gives
us its phenomena and events: we must ourselves find out
whether they are rule or exception, whether they are to be imitated or
shunned.

When you have so seldom spoken the word of hope, --- you say
somewhere:


\begin{verse}
Doves I have none of; \\
Doves are, after all, hope's birds ---
\end{verse}

then it is perhaps because you have thought that hope is a medicine
that must be dispensed in small doses, since human beings are so
inclined to revel in it. Or it is perhaps because the
true hope must be won through understanding won by struggle. You follow
in your works the development of events and characters
step by step back to the decisive point where the ground
% --- p. 215 ---
\setcounter{footnote}{0}%
\opage{215}for everything that followed was laid. You put your finger on the place
where the sickness took its beginning. It is as if you said
to us: the stream that cannot be stopped or diverted once it
has grown strong might perhaps have been stopped or
diverted at its first source! This is your way of teaching
us hope: through understanding of the cause of disharmonies
we learn how life can be made harmonious; through understanding
of downfall we learn how life can persist in the struggle.
When you traveled in Egypt, you saw temple ruins and overturned
images of gods rising out of the sand of the desert like skeletons of
caravans that had perished in the drifting sand. You asked yourself:
why did all this perish, so that it does not even, as
so much else that had to perish does, live in the memory of men
and in their thoughts? --- and your answer was: this culture
lacked personality and therefore true vitality. This understanding
becomes then the foundation of your hope:

\begin{verse}
On buried caravans \\
I build the ways of our future!
\end{verse}

--- Your third struggle, the struggle over the artistic form,
is connected with this. In order to be able to present the development of events
and characters you have introduced a new artistic means, the conversation that points backward, explains and settles accounts.
You have introduced perspective into dramatic art. Under
the quiet but firm progress of the speeches we are led from the situation of the moment
back into the characters' past. Characters and events
unfold for us from their first disposition. Through secret passages we are
led to a point where great views over human life and its
conditions become possible. What a treasure of psychological observation and
practical wisdom of life is not laid down in these conversations!

The mighty world of images and thoughts you have created captivates
our senses and our imagination. But it is not given us merely for
contemplation. It brings stir into minds; it lets us understand
the course of fate in us and outside us; it awakens our self-activity
to resistance, to the clearing away of that which cows and
hinders. Thereby it prepares us for a new and vigorous grip
on life and its conditions. ---

% --- p. 216 ---
\setcounter{footnote}{0}%
\opage{216}Let us then now thank \textit{you} for your faithful work. You
have in the long, solitary years of work asked only about one thing, about
what true art demanded of you, and by being true
to art you have been true to yourself. We thank you
not merely for your works, but also because you, by your
unshakable holding fast to the great aim you had set yourself,
have become a model for everyone who works for truth
and beauty in an age that so often seems to be taken up with
quite other things. That you have been able to carry the struggle to its end and
have compelled recognition of the work of your life is not
merely something we congratulate you on. For us it is a
ground for hope, a gleam of light for understanding life. The drama of your own
life draws toward its close more brightly than so
many of your invented dramas do. May this gleam of light
accompany you when the days of celebration are replaced by the days of work, for
you the dearest days, and may many fruitful
days of work still fall to your lot!

% --- p. 217 ---
\setcounter{footnote}{0}%
\opage{217}\essay[Review of Bjørnson's ``Poul Lange og Tora Parsberg''. (1898)]{BJØRNSTJERNE BJØRNSON:\\ \textit{POUL LANGE OG TORA PARSBERG} [PAUL LANGE AND TORA PARSBERG].}
\streg

``Sometimes a god descends to earth'', says Poul
Lange in Bjørnson's new drama. The god here is a woman who
helps Lange to understand himself and to fight against the
inner and outer resistance to which he so easily succumbs. Of
this woman it is said, and she repeats it herself, that she always
tries the impossible. But an old saying says that
against the impossible even the gods struggle in vain. And that
Tora Parsberg also experiences.

Should the content of the new drama be given in few words,
it may perhaps be given in this way. And herewith is also expressed
the means the poet has used to let Poul Lange's
character step forth before us, --- the only means by which
a character like his could be brought to revelation. Inward,
considerate, conscientious, but at the same time anxious before
others' judgment, especially before their laughter and their scorn, he is
one of those human beings who open up only after a hard
struggle. He is a ``diffident concealer''. In his enthusiasm for great
ideas and ends he is inclined to overlook and despise
the means that human beings who call themselves practical consider
necessary for putting something through. If he has then himself, from
inexperience or in need, let himself be seduced into using means
that could be blamed or mocked, he feels all his life through ever
% --- p. 218 ---
\setcounter{footnote}{0}%
\opage{218}anew the sting of it. For on account of his inclination to
self-reflection his past, especially its dark parts,
is constantly renewed for him in painful memory. He is moreover accustomed to
compare himself with the greatest minds and therefore to judge himself severely. The paralyzing self-criticism that so easily
leads him to self-abandonment stems not merely from his being so
attentive to others' judgment and mockery, but perhaps especially from the
ideal standard he applies to himself. He has at last
come to the point that he does not regard himself as normal, ---
that he fears being ``a lout''. This cowed self-esteem
and the anxiety it betrays is a secret he broods over.
Only at certain times, when he is absorbed in work, or when
a hearty helping hand is extended to him from outside, can he
forget the inner darkness and look brightly on the future. He has
a need to act according to his actual, inner nature, and he
acts in decisive moments also according to it; but there is in
him an undercurrent of self-reflection and of listless doubt that undermines
the ground of his noble nature, so that the structure that might
have stretched its spires toward heaven is sapped and sinks
together.

It took a woman like Tora Parsberg for such a
nature to be able to open itself. She has followed him through
his life, as a spectator at close range in his youth and manhood;
she has been the confidante of his confidants; and she received
from his conduct in smaller and larger circles an ideal of a man,
which she preserved in the inwardness of her strong and faithful heart,
until the moment came when his declaration of what she had now
become for him, and his prayer that she would unite her
fate with his, gave her courage and right to intervene in
his life. In the great conversations with him in the first and third
act she teaches him to know himself. These conversations do not run
in such a way that Poul Lange himself is led to say how
he is; it is predominantly she who proclaims to him what
his actual nature is. She hardly lets him speak out,
so sure is she in her understanding of him. The deepest abyss
in him, the one it proved impossible to bridge, she has
nevertheless not discovered; the depth of anxiety, self-criticism and weakness
% --- p. 219 ---
\setcounter{footnote}{0}%
\opage{219}in his nature she did not suspect. But by her inward understanding
and sympathy she teaches him how precisely the weak and
sick in his nature is connected with the fine, noble and
great in him, how the persecution and the laughter that strikes
him prove his significance, not --- as he fears --- his
wretchedness. Over against the fanatical suspiciousness of the politicians of the moment
and their stupid hatred she shows him how firm a hold he has
in his work for thoughts and tasks that condition the future
of the race and ever anew will present themselves to human
consciences.

It is the ideal Poul Lange, as his image
had formed itself in her strong and deep mind, that she holds up
to the real Paul Lange. But it is now once and for all
so that the gravest problem in our inner life is whether we can
find the comfort to hold fast to our self as it is in our best
moments, or whether the dark, weak and low that we also
know in ourselves should be our actual self. This problem
Poul Lange had not been able to solve himself; he now gets in Tora
Parsberg the best helper for it that can be conceived. Perhaps
her superiority is portrayed with too strong strokes. Her
words become at times great, mighty tones that pass by
without properly being individualized into lines according to the dramatic
art's strictest demands. That I do not venture to
decide. But no one would surely, for the sake of greater technical perfection,
have been without the vast energy with which
the highest thoughts, such as can keep up the courage of life, here stream
forth. They are tones from an Aeolian harp, that can make up for the lack
of much filigree work one has imagined to be poetry.

My pen is not supple enough to portray the wealth of
individual small traits with which the poet has been able to make
the two main persons in the drama living and vivid. My
endeavor is only to bring forward the fundamental lines I have found
in the poetic characterization. A work of art is and remains after all
in the end incommensurable with analysis.

A wholly different image of Poul Lange from the one Tora Parsberg
draws forth presents itself of him with those who feel themselves
% --- p. 220 ---
\setcounter{footnote}{0}%
\opage{220}thwarted or disappointed by his political conduct. The head
of the ministry in which he sits has long shown himself
intriguing and faithless toward him. It affected him so
that he almost became insane and was near to committing suicide.
Now he has submitted his application for dismissal, and the question
is whether he will support his old chief against an attack in
the Storting. He feels that it accords with his nature to do so\dots{}
The head of the ministry is a man highly deserving of his people; he is moreover
an old man who has hardly long left. And above
all --- Poul Lange does not like public chastisements; it does not accord
with his nature to be an executioner. He prefers to
work positively for the new and turns his back on the bad
and base. He will preserve his strength for better times. Yet he answers
evasively an invitation from the Court to defend
the head of the ministry and the offer of the post of envoy
in London, which is made him in connection with it. And when one of
his closest friends urgently urges him to stay
away from the decisive meeting in the Storting, he lets himself be moved
to promise this. Immediately afterwards the first great conversation
with Tora Parsberg takes place, in which she opens happiness for
him, awakens his self-confidence and strengthens him in the need to
follow his innermost nature. When he intimates that people want him
to do something that actually conflicts with his nature, she says:
``So others will tell you what comes naturally to you?''
Proud of victory she leaves him without suspecting that she has given him
a false security. That he does not inform Arne Kraft that,
after closer consideration, he takes back his promise to stay
away from the meeting, although four days pass between,
must surely be explained by his overlooking in his intoxication of happiness the necessity
of it. He, who otherwise has too many doubts, has here
too few scruples.

After he has, in the decisive meeting, defended his old
chief as the one who is persistently best able to unite the people,
he takes part in the evening in a festivity where many of the Storting's
members are present --- and here he then himself becomes the object
of one of the ``public chastisements'' he had not wanted to be
party to when it concerned others. People shrink from offering
% --- p. 221 ---
\setcounter{footnote}{0}%
\opage{221}him their hand, call him strictly to account and declare
that he has lost the confidence of the best. When even his closest
friend takes part in this condemnation, he collapses.

For those who condemn him, his mode of action stands as
inexplicable by other motives than striving after personal advantage.
``He knows well enough what he wants!'' The truth, it is declared,
must in politics wait until one has need of it. There is at
every moment a single cause at hand, and he who hinders
it is to be ruthlessly thrown aside. The mildest explanation of his conduct in this circle is that it springs from
morbid sentimentality.

Masterfully it is shown in the drama through a series of details,
hints and coincidences how Lange's
character could appear to the politicians of the moment in this
way. The ground of the distorted image that has formed with
them lies partly in the fine and noble in Lange's character, partly in
his weakness, vacillation and thoughtlessness, partly in his opponents'
external and fanatical view.

\stjerner

After the violent collision Paul Lange is shaken in
his innermost. The old doubts emerge; the anxiety of being
abnormal, of being a lout, rises again in all its force,
as an unceasing echo of the fanatical words that have sounded
to him. Darkness spreads before him on all roads. Even the bright
road Tora Parsberg had opened to him he dares not trust:
how can he draw her along into his disgrace, draw
her down into his own darkness? The second great conversation with
Tora Parsberg, who now again seeks him out, lifts him indeed once more
up after a strong spiritual struggle between the two personalities,
a struggle between hope and despair, between life and death.
She leaves him again confident, so confident that she does not even try
to get him to hand over the pistols that she knows he
has in his keeping. Here again the limit of her understanding
and of her power shows itself; to the lurking despondency and the
untiring doubt that had undermined the foundation of his
personality she had not penetrated, and she could probably not either
have overcome these dark forces completely. Hardly
% --- p. 222 ---
\setcounter{footnote}{0}%
\opage{222}has she left him before these forces begin to stir
again. ``What must she after all think? Not today, but
the next day? When the feeling of victory is over? --- By God, I cannot
bear the slightest shadow of humiliation any more!'' But in the
same moment a new humiliation comes; a telegram brings
the news that his enemies have been able to inflict
on him a new public affront. Then the cup runs over, and
he directs the killing shot at his breast.

This drama will not be understood by those who have not learned
how the weak and the great, the heavy and the light, the
low and the noble can dwell together in our nature in such a way that a
human being's mode of action cannot be explained from a single
series of motives alone or be valued from a single point of view.
It is one of the shadow sides of public life as it
shapes itself in our day that explanation and judgment of human beings
take place so quickly and on the basis of external situations
determined by the relations of the moment. The poet has shown us
what tragedies may spring from it when in the depths it looks
wholly different from on the surface, although depth and surface
nevertheless belong closely together.

A poetic presentation that has reached its goal has an
individual character, has portrayed a single definite type; but
general truths cannot be derived from it. It would be
incorrect to derive from Paul Lange's fate that fidelity
to the party one has joined should not be a duty.
The drama points only to this, that it cannot be the highest
duty; the individual's striving for agreement with his
innermost nature can lead to a break with party morality. There is
a sanctuary to which the knowledge of men and the valuation of men
that can be acquired in political life alone gives no
access. Strong natures are able to preserve this
sanctuary and perhaps even to compel others to acknowledge
it; weak and wavering natures go under, because they let the cries
from outside resound so violently in the interior of the sanctuary that
the inner voice cannot be heard. But even if the walls of the temple were
% --- p. 223 ---
\setcounter{footnote}{0}%
\opage{223}too weak, they can nevertheless have enclosed real images of gods ---
and that the poet has here shown us.

If the poet has had a model for his main figure,
then he has rendered him the service the dying Hamlet
asked of his friend:

\begin{verse}
Explain me and my cause to those who do not understand, --- the whole cause, the least with the greatest!
\end{verse}

% --- p. 224 ---
\setcounter{footnote}{0}%
\opage{224}\essay[Julius Lange's Letters. (1902)]{JULIUS LANGE'S LETTERS.}
\begin{center}{\small (PREFACE TO THE GERMAN EDITION).}\\
(1902).\end{center}
\streg

Julius Lange has significance not merely as an excellent
researcher in the domain of art history, but also by his whole
bright and warm personality, which pervaded all his conduct
in life as in science. Hence it comes that his art-historical
writings not merely yield rich instruction about the works
of art, their origin and their mutual relations,
but also, partly through direct utterances, partly, and not least,
by the basic mood by which the presentation is borne, contain
a treasure of experience of life and sound thinking about human relations.
It will therefore be of great interest to get to know him
from a more personal side than his scientific works
make possible, --- to get to know him as
he was when he moved among his nearest and his friends. And thus we get to know him in his letters.

We see Lange in these letters as the sound, vigorous and courageous
human being, regarding life, as he regarded
art, with childlike immediacy and yet with the thinker's
and the tried human being's deep understanding. Many
shadows fell on Lange's way in the later years, and his view of life
gave him no dogmatic consolation. With manly understanding and
energy he held fast what had given life value for him,
and worked for it to the last, in the midst of a painful
and depressing illness. We meet in these letters not least
% --- p. 225 ---
\setcounter{footnote}{0}%
\opage{225}numerous examples of his delightful humor, which often helped him
himself and others over a dark hour or cast brightness over a
festive gathering.

The having known Lange personally and having belonged to
the circle in which he freely opened the treasures of his mind and his fancy,
is counted by his friends as one of the best things of their lives.
Through his letters some of this good can now fall to a larger
circle.

In addition, the letters naturally deal for a large part
with what lay closest to Lange's own heart, art
and its history. We hear of his plans, his journeys and
his sketches. In improvised utterances much comes
forth of what he later worked into and carried further in his life's work, or
what was to have been the great work of his life,
but became only a series of fragments, --- ``The visual arts' representation
of the human figure in the different periods of history''.

Lange once defined artistic value as the value
a subject has for the artist, and which it gets through his presentation
for others. This definition (to which, purely aesthetically, various objections
must be made) is characteristic
of Lange as art historian. Through the study of the visual arts
in their historical main periods he wished to discover
how --- that is to say, with what eyes and in what mood
--- the human figure has been conceived. And the conception
of the human figure rests in turn on the conception of life
and life's conditions. Thus the study of art acquired for Lange
its great human significance --- an essential contribution to
understanding of human life in its development.

But that definition fits moreover Lange's own works
and his very personality. His works: for through
them the visual arts gain increased value for us by virtue of
the value they have had for Lange himself. His personality:
for through it our feeling for life is strengthened and enriched,
as we see what Lange by his energy, his humor, his warmth
and his clear understanding got out of life.

% --- p. 226 ---
\setcounter{footnote}{0}%
\opage{226}\essay[Goethe (1899)]{GOETHE.}
\begin{center}(1899).\end{center}
\streg

On 28 August 1899 it is a hundred and fifty years since
Goethe was born. The thoughts of many will on that day seek to settle
the debt they owe to this great genius. His poetry has
woven itself into our life of feeling and thought, and his personality
stands before us as a mighty model of a harmonious, yet at the same time deep and earnest
conduct of life. The impressions
we receive from him are, of course, not entirely unmixed. There are many of
his poems that we must let lie, because they are due to his weak
hours, or because they move in forms and moods whose
time is past. And there is in his conduct of life much
that perhaps can be understood and excused, but in any case
not admired. And yet, despite the obscuring
power of time and criticism, his personality and his work rise
up to this day like a sunlit mountain peak and mark one of the high points
of human intellectual life.

\section*{I.}

How high he rises in comparison with other masters,
it is in a certain sense impossible to determine. Great minds cannot
simply be compared, because in the relation between
them it depends not merely on a more and a less, but also on
differences with respect to what and how: they belong to different
orders, do not all lie on the same line, so that the
% --- p. 227 ---
\setcounter{footnote}{0}%
\opage{227}distance
between them could not simply be measured. If we apply
the highest standard, we shall yet surely meet in the first rank
within the world of poetry only two great names: Homer
and Shakespeare. Goethe himself would agree with us in that. When
in the Romantic period an attempt was made to set Tieck up at
his side as an equal, he found it unjust, --- just as
unjust as if he himself were to place himself beside
Shakespeare. ``Tieck,'' he said, ``is a talent of great significance;
but when one wants to place him on a level with me, one goes
astray. I can say this straight out, for what does it matter to me?
I have not created myself. It would be as if I wanted to compare
myself with Shakespeare, who likewise has not created
himself, and who is nevertheless a being of a higher order, to whom I
look up and before whom I bow in reverence.''

Here we have at once a trait that is peculiar to Goethe:
healthy self-confidence, united with a clear self-limitation
grounded in deep understanding, and with the assurance that at the
ground of our highest activity there lie forces that we ourselves have not
given ourselves. For Goethe mastery lay in limitation, and
although the energy that works within the limitation is our
own, is one with ourselves, it is nevertheless not our own work
--- any more than the limits we find for our being are our
own work. Goethe's whole conception of life and art is borne by
this thought, which it is impossible to express in a single word.
Words on the whole easily become too heavy when one is to speak of
Goethe. He is able to unite in himself qualities, endeavors
and points of view that most often stand as incompatible, and for whose
union no exhaustive expression can be found. We
must spell our way forward in order to understand him.

\section*{II.}

His faith in himself and his striving to develop
himself and to draw nourishment from everything that meets him on his
way, from human beings as well as from nature and art, can often look
like egoism, and yet he did not regard his desire and his yearning
as justified apart from the great connection of life and ideas.
% --- p. 228 ---
\setcounter{footnote}{0}%
\opage{228}He did not regard himself as an end and everything else and everyone
else merely as means. If one wishes to call him an aristocrat of the mind,
this word needs in any case a closer explanation.
With pride he could say in his old age that
in his calling as an author he had never asked what the great
mass wanted, and how he could benefit the whole, but
had only striven to make himself more insightful and better, to raise his
own personality's value and then to express what he had recognized
as good and true. But it was not because he denied
that the value of a work is determined in the end by its lasting significance.
On the contrary: when he would not let the taste of ``the great mass''
be the decisive standard, it was precisely because it is easily
dazzled by the momentary glitter, whereas it is the
calling of creative minds, in solitude and often unrecognized, to form
and stamp that which is to have value for long ages. It is what
the poet in the Prelude to ``Faust'' says to the Theater Director
and the Clown:

\begin{verse}
Was glänzt, ist für den Augenblick geboren; \\
Das Aechte bleibt der Nachwelt unverloren.
\end{verse}

With this accords Goethe's ethical theory, in which he has
Fichte as predecessor and Carlyle as successor. ``The good,''
he says in a conversation with Eckermann, ``is no product of
human reflection, but is innate and co-created beautiful nature.
More or less it is co-created in all human beings, but
to a high degree in certain, quite exceptionally gifted minds. These
have revealed their divine inner nature through great deeds
or teachings, and it has then, through the beauty of its appearance,
won the love of human beings and powerfully led to reverence
and imitation. But through experience and wisdom the
value of the morally beautiful and good could come to consciousness, in that
the bad proved in its consequences to be something that disturbed
the happiness of the individual and of the whole, whereas the noble and
right showed itself to be something that brought about and established the particular
and the general happiness.''

For Goethe there was a constant reciprocal relation between the chosen
and the multitude of human beings. Just as the value of works and
% --- p. 229 ---
\setcounter{footnote}{0}%
\opage{229}actions
is determined by their lasting effects, which are
experienced only by turning the view from the single creative
mind outward to larger circles, so on the other hand
this mind itself is also an effect of the past. Goethe was filled
with gratitude toward significant predecessors and contemporaries.
``One is always talking of originality,'' he said in a conversation,
``but what does that mean! As soon as we are born, the
world begins to work upon us, and it goes on so right up to
the end. And what could we on the whole call our own, except energy,
force, will! If I could say what I owe to
great predecessors and fellow-livers, there would be nothing left.''
In a well-known humorous poem he has set out how the nature
of the individual is a product of the nature of parents and ancestors,
so that in the end it may become a question what is original in
the whole fellow.

\section*{III.}

There is an inner correspondence between these thoughts and
Goethe's whole conception of nature and view of life.

Each single organ and each single organic being, plant
or animal, is for his ingenious view, in all its independence
and peculiarity, a revelation of types and laws
that rule in manifold nuanced forms throughout the whole of nature.
Independence and lawful connection are not oppositions for him.
On the contrary: the independence of the single organ and of the single
being becomes possible precisely through their standing
as expressions of and effects of a comprehensive connection of nature.
Through constant and delicate transitions one form is connected
with another, the greatest with the least in nature's
infinite realm. In his scientific treatises and in
his poems on the metamorphosis of plants and animals Goethe has
developed these thoughts.

They are found again in analogous form in his religious conception. In
his poem ``Prometheus'' he gives vent to the swelling self-confidence
that sprang from the consciousness of the creative poetic power
that stirred within him and that made him independent of
% --- p. 230 ---
\setcounter{footnote}{0}%
\opage{230}external powers. With Titanic defiance he turns here against every
deity that is thought to intervene in life from outside. But with
``Prometheus'' another poem must be set together, ``The Limits of
Humanity.'' The Titanic self-confidence in the face of all external
powers rests on the assurance that the force by which he is
animated stems from the great connection of existence, from the
deepest source of all life and all activity. Thus independence
and limitation are closely connected.

The master of thought who had above all given Goethe impulses
toward his view of life and world-view was Spinoza. In beautiful
and warm words Goethe has expressed in his autobiography what he
owes to this philosopher. Life demands at so many points
self-renunciation, resignation. But most people get through
by putting one small enjoyment in the place of another: occupations, inclinations,
hobbyhorses, everything is tried, in order at last to proclaim
that all is vanity. ``There are only few people who foresee
such an unbearable feeling, and who resign once for all in order
to avoid all partial resignations. These convince themselves
of the eternal, the necessary, the lawful, and seek to form
for themselves concepts that are irrefutable, indeed that, instead of being
canceled by the consideration of the transitory, are precisely
confirmed by it.'' The great disinterestedness and the deep
peace that prevail in Spinoza's way of thinking, because he
arranges his view in such a way that the particular is always seen in
its connection with the whole, as limited and yet peculiar,
won Goethe's heart for him. He did not pose as
any scholarly interpreter of Spinoza, but took from his writings what he
could use. In reality, indeed, the task that Spinoza set
for knowledge: to see the single being in all its peculiarity
as a link in the great lawful connection of existence, was of such a nature that it
cannot be solved purely scientifically, but that its solution depends
on the capacity for artistic intuition. Between what Spinoza demanded,
and what Goethe has achieved in his greatest works, there is
therefore an even greater kinship than Goethe in his youth
was aware of. It was more for the art of life than for the art of poetry
that he learned from Spinoza. He learned an art of resignation,
% --- p. 231 ---
\setcounter{footnote}{0}%
\opage{231}that does not make one cold and poor, but makes possible a constant youth
through surrender to the great connection of ever-young nature.
It was above all in Spinoza that Goethe found the idea of a
God who does not ``push the world from outside,'' but works and stirs
in the innermost of existence.

Goethe considered it impossible to find a conclusive expression in thought
for the way in which elements and forces work
together in the whole of existence as well as in the single being. Here
he found one of the ``Limits of Humanity.'' In Faust's profession of faith,
which is undoubtedly Goethe's own, it is said that every
name is here only ``sound and smoke'' in the face of the overflowing
fullness of life, in which man feels himself borne by a
power that is greater than he himself. For
Goethe the religious feeling was a need, out of gratitude, voluntarily to surrender himself to
something higher, purer, unknown, and to find an expression for the eternally
Unnamed. (The Marienbad Elegy). But man is, in his
opinion, not born to solve the problems of the world, but to seek
out where the insoluble problem lies, and then to keep within the limits
of the comprehensible.

\section*{IV.}

Yet let us not forget that Goethe was not a philosopher, although
so many of his poems belong among the weightiest in thought and the most profound
that have ever been written.

The relation between his life and his poetry accords with
his character as a whole. What gripped him in life with joy
or sorrow, melancholy or remorse, found vent and found form in his
poetry. Here too self-assertion and surrender went together as one.
His producing was surrender: for the images often swarmed
forth in him in such overflowing fullness that he scarcely had
time to write them down. During his roaming in free
nature, the changing scenes as well as his own experiences
and thoughts called forth poems that he hummed to himself but could not
always remember afterwards. Not only the material, but also the form
made demands on him, to which he bowed. But in and
with this self-surrender self-assertion was at work: for
% --- p. 232 ---
\setcounter{footnote}{0}%
\opage{232}artistic creation was to him a self-liberation, a liberation
from what overwhelmed, weighed upon and oppressed him:

\begin{verse}
Und wenn der Mensch in seiner Qual verstummt, \\
Gab mir ein Gott zu sagen, was ich leide.
\end{verse}

In his nature poems the most delicate moods of nature vibrate.
Like no one else before or after him he was able to live and
feel himself into nature. He marks (after Ruysdael and Rousseau) a new epoch in the history of the feeling for nature. As regards the form of expression,
attention has been especially called to how many new compound,
happily formed words he has produced in order to be able to
express images of nature and moods of nature briefly and picturesquely,
words such as ``morgenschön,'' ``wellenathmend,'' ``Mondendämmerung,''
``Nebelthal,'' ``Lebensfluten,'' etc. All feeling for nature rests
on an involuntary transference (or, as the Germans aptly
call it: Einfühlung) of human mood into nature.
Goethe was able to let this involuntary transference have
its free course and yet to become clearly aware of the difference between nature
and human feeling. In the magnificent poem ``Ilmenau,''
which furnishes so many contributions to his characterization, glorious
scenes of nature are depicted, while yet the opposition between the poet's
mind and the nature that attunes him makes itself felt.
Nature's calm stands in contrast to his own worry, and
he says to the trees in the forest:

\begin{verse}
Ich sorge still, indesz ihr ruhig grünet.
\end{verse}

He cannot forget how hard people in the midst of this
great nature must toil to sustain life. The hard life of the peasant, the charcoal burner
and the miner appears before his eyes,
though compassion loses something of its bitterness through the consciousness
of what he, as the Duke's counselor, has been able to do
for the region in which he moves. Besides, he expresses his worry
over the way in which his works take effect; what he
has wanted has not worked as he imagined it:
\begin{verse}
Ich brachte reines Feuer vom Altar; \\
Was ich entzündet, ist nicht reine Flamme.
\end{verse}

% --- p. 233 ---
\setcounter{footnote}{0}%
\opage{233}And to this is added the worry for his princely friend,
whose restless mind has such difficulty working its way toward
the self-limitation that is especially necessary for one who is to
be a leader of others. But the fog that has settled over the
landscape disperses, the region lies in the full clarity of day, and
the dark mood gives place to a shining hope.

Goethe has wrongly been accused of transforming, in an egoistic
manner, all relations of life into poetic subjects or
even into means of poetic self-enjoyment. The self-liberation he
found in poetry testifies precisely against this accusation. There
runs a quiet self-accusation through a whole series of his poems,
--- Gøtz, Clavigo, Faust and more. He could not forget that
the great task that his self-development set him had in his youth
led him to a break that put an end to a
young woman's happiness in life. Friederike's fate Goethe has felt
as his misfortune and as his guilt. One hears this through
Faust's words:

\begin{verse}
Bin ich der Flüchtling nicht? Der Unbehaus'te? \\
Der Unmensch ohne Zweck und Ruh, \\
Der wie ein Wassersturz von Fels zu Felsen brauste, \\
Begierig wüthend nach dem Abgrund zu! \\
Und seitwärts sie, mit kindlich dumpfen Sinnen, \\
Im Hüttchen auf dem kleinen Alpenfeld, \\
Und all ihr häusliches Beginnen \\
Umfangen in der kleinen Welt.
\end{verse}

\begin{verse}
Sie, ihren Frieden muszt' ich untergraben! \\
Du, Hölle, musztest dieses Opfer haben!
\end{verse}

Goethe himself has said that his poems were occasional poems,
called forth by life itself, first and foremost his own life, and ``composed
under immediate intuition.'' And when he has also
spoken of his earnest way of taking life (des Lebens
ernstes Führen), he has been able to do so with the consciousness
that he has sought to let life come into its own no less
than art. He found it apt when it was said
of him that his striving aimed at giving the real a
poetic form. That was his life task, and it was his earnestness
% --- p. 234 ---
\setcounter{footnote}{0}%
\opage{234}about this life task. It is, after all, also a task that benefits
life itself, in that the image of life that poets
form in turn reacts back upon the continuing life. One therefore
has no right to speak of an opposition between the human being
and the poet. To be a poet and as a poet to give an image
of life was Goethe's life work \textit{as a human being}\footnote[1]{Compare what Goethe says about Schiller in Epilog zum Lied von der Glocke: Er wendete die Blüthe höchsten Strebens, / Das Leben selbst, an dieses Bild des Lebens.}, and if it brought
him into conflicts with other human relations and tasks, conflicts that left their sting in his inner being, then it would
in any case not be exhaustive to explain these
conflicts by a contradiction between the poet and ``the human being.''
At the ground of them lay a strife between two \textit{human} tasks,
a strife in which it is not so easy to pass judgment
as has often been done, but from which Goethe, incidentally, by his own
account did not emerge without scratches.

\section*{V.}

The study of Goethe is of such great interest from a
psychological point of view not least because we find in him all the different degrees of the consciousness by which the imagination works.
He could form with clear deliberation as well as improvisationally,
and he knew --- especially in his youth --- the entirely
involuntary, instinctive working of the imagination. Some of his
youthful works he composed in a kind of somnambulistic state,
in which he did not know what he was doing, although he noticed
that something was working within him. The images presented themselves suddenly
quite finished to him as ``free children of God'' who said:
``Here we are!'' It is what he later called the demonic
in genius. No more than nature does genius scrape together the elements
to its works as wholes. There is from first to last a
reciprocal relation between whole and elements that defies every
analysis. Therefore it is also impossible to express the content of a
poetic work in a single thought, a single idea, --- just as
% --- p. 235 ---
\setcounter{footnote}{0}%
\opage{235}impossible as it is to express the essence of existence in a single
thought. A great poetic work is, just as nature is,
incommensurable for thought. ``I am,'' says Goethe, ``of the
opinion that the more incommensurable and incomprehensible to thought
a poetic production is, the better. There they come and ask,
which \textit{idea} I have sought to embody in my Faust. As
if I myself knew it and could say it!\dots{} That a human being
who from grave errors ever strives upward toward the best
can be redeemed, that is indeed a good thought, which is effective and explains
much; but it is no \textit{idea} that was to be laid at the
ground of the whole and of every single scene. It would then also
have been a fine thing if I had wanted to string a life so
rich, variegated and manifold as I have presented in Faust
on the meager string of a single pervading \textit{idea}!''

It was not Goethe's intention to deny the right of
thought, neither in regard to the works of art nor in regard to those of nature.
But he maintained the impossibility of the work of thought being
brought to completion in an exhaustive idea, and he feared
--- as the speculative philosophy that unfolded around
him also gave cause to fear --- that one would conclude hastily,
so that abstract thoughts would be put in the place of
immediate living intuition.

Yet for Goethe's whole life and activity one could scarcely
find a better expression than the thought that, by his own
admission, expresses an essential side of the content of his
greatest poem: the thought of constant striving. Through his
long life runs an uninterrupted labor of the mind, a constant yearning for
fullness and clarity, and the warmth of the heart is preserved right up into old age.
By this the difference between the young and the
old Goethe is not to be denied. Calm contemplation quite
naturally came to play a greater role in the old man than
in the young one, who was tossed about in the period of tribulation. But already
in his youth, in the midst of the heavy seas, he looked up to the fixed
stars in order to orient himself and get a hold on himself. And the stars
did not pale when the heavy seas subsided, so that calm beholding
became possible.

% --- p. 236 ---
\setcounter{footnote}{0}%
\opage{236}

\section*{VI.}

Even the greatest has his limitation, and if he is to be
something for all times, there is much one must disregard or
alter. Neither in the mental nor in the material world
are the conditions of life the same at all times. The tasks that are to be
solved and the obstacles that are to be overcome change with the times.
It is then no wonder that much seems strange to us ---
even in the greatest of the generation that inaugurated the century
that is now about to run out.

People have at times wished to regard Goethe as one of Darwin's
predecessors. But precisely the points most characteristic of Darwin's doctrine
are lacking in Goethe: the idea of the struggle for
existence and the idea of the historical evolution of species through
descent. Goethe's conception of nature is marked by an enthusiastic
optimism. He sees the life of nature as a harmonious unfolding of
eternal types. To be sure, there is great difference between the conditions for
life on the whole; but for Goethe there is no bitter contradiction present
between what the persistence of the whole requires and what the
single being strives after. Nature is a mother who hides
herself behind her children, but in spite of her invisibility does not leave
her maternal love without witness. When it checks the need for development
at a single point, it is only to lead
development gently over to higher forms.

\begin{verse}
Doch hier hält die Natur, mit mächtigen Händen, die Bildung \\
An und lenket sie sanft in die Vollkommenheit.
\end{verse}

For Darwin nature shows us no harmonious play of forms
and forces. The way to new forms goes through inhibition,
suffering and destruction. Stronger shadows now fall over the life
of nature: it is only on the surface that nature shows us light
and peace. When Goethe in the Ilmenau poem felt nature's calm in contrast
to his own sorrow, the new conception of nature teaches that
nature's life too, no less than man's, can be checked
and subdued, so that it becomes more difficult than before to ``feel oneself
into'' nature. Even while the trees stand apparently calmly
and turn green, they fight among themselves for light and air and room to
spread.

% --- p. 237 ---
\setcounter{footnote}{0}%
\opage{237}Also in the human domain itself we have been compelled
to lay greater weight on strife and oppositions than
Goethe and his time did. For Goethe there was a deep analogy
between nature, art and social life. His ideas about these
``three great regions of the world,'' as he calls them, developed
simultaneously and in agreement with one another. But in
the social domain too the shadows now fall more strongly, and the
doubts have become more piercing. If the possibility of a harmonious
solution is to be maintained, a vigorous will is needed and a
greater contribution of hope and confidence than Goethe's calm beholding
involved or considered necessary. Goethe's doctrine of harmony
does not lack energy. But it was not set the task
of spanning such great oppositions as are set an attempt
in a similar direction in our day. His view of life becomes
naive in our eyes.

And yet --- perhaps naivety is right after all. True
greatness appears with the stamp of naivety, and in the face of life and
its problems it perhaps holds in the end that we must continue
to be like children. Even though a Faust, if it were written in our
day, would have to contain other and in part more far-reaching problems
than the Goethean work, the fundamental thoughts would nevertheless surely
continue to be the same. Goethe still stands as our
teacher and will continue to stand so for long ages. The ``children
of God'' who appeared before his thought and his imagination have
not yet become old men. Again and again we return to
his rich world of images and thoughts, to his moods of nature and
his experiences of life. When we celebrate his commemoration festival, it is not
because his mind has become a mere memory for us; it accompanies us
in our striving to feel and understand life, and everything that
has appeared since his time in the domain of art and thought
has not let us outgrow this leader.

\end{document}
